% =====================================================================
% CIF Variable, Unit, and Equation Reference
%
% This appendix is the dissertation's NORMATIVE NOTATION LAYER. Every
% CONSTANT is drawn from the verified derivation record rather than by
% cross-reference, so no VALUE here depends on another appendix. It does
% cite Appendix A for three DERIVATIONS (the skim inversion, the cubic
% identity, the scaling exponents); those citations go through \xref and
% degrade to a literal when the target is absent. It names Appendix B
% twice, in plain text: for the proven status of that identity (A.1) and
% for the corrected Xia scaling Dinh's schedule follows (A.7). So this
% file still compiles and reads correctly on its own. (Corrected at the
% v108 freeze: this sentence said "Appendix B once" and "those four
% citations go through \xref"; the A.1 mention was never an \xref, and
% the A.7 mention came with the step-2 unfreeze batch, ledger DG-24. The
% batch's corrections to this glossary are listed item by item in
% FREEZE_v108.md and recorded in DEFERRED_GLOSSARY_CHANGES.md.)
%
% (The earlier wording claimed the appendix "does not cite any other
% appendix." That has been untrue since the skim-inversion row acquired
% its provenance, and the claim is corrected rather than the citations
% deleted: the derivations genuinely live in Appendix A, and a pointer to
% them is worth more than a self-containment boast.)
%
% SUBSTANTIVE CHANGES IN THIS REVISION (detail in the Change Log):
%   foot region corrected; supplied vs refit prefactor named apart; the
%   junction campaign disambiguated from the retired common-origin claim;
%   two factors vs three model gaps recorded, with phi/G/Lambda added;
%   the f_eff retraction narrowed; the accuracy band defined; the
%   aspect-ratio symbol crosswalk to the code; the kernel-prefactor
%   pairing rule and kappa_F added; the two calibration sensitivities
%   named apart; analytic and CFD depth factors separated; values
%   generated rather than typed; cross-lineage framing corrected and the
%   Davis leg added; a selectivity-axis validation section added; the
%   +1.64 sigma closed-form distance surfaced; F split into its two
%   identities; platelet-volume terminology corrected; a Status and Role
%   legend added; a validity-envelope inventory added; the entrance-length
%   convention added; the gap symbol g deprecated in favour of G with the
%   link length ell_g added; the reduced size ratio moved off the letter u.
%
% HOUSE STYLE: no prose em dashes. The triple hyphen appears as an empty
% table cell, and INSIDE A VERBATIM QUOTATION where the em dash is the
% SOURCE's punctuation and altering it would misquote. The Lezzar anchor
% in the source-quote index is the only such case today ("the magnitude of
% this flow fraction---not the width of the gaps---that determines the
% size cutoff"). texlint.py already encodes this exemption by stripping
% quoted spans before counting; this sentence brings the stated policy
% into line with the enforced one.
%   The double hyphen appears in numeric and page ranges and in name
% compounds (Hagen--Poiseuille, Zweifach--Fung, log--log), where an en
% dash is the correct typography.
%
% TABLE STYLE: the variable tables are longtable, NOT table floats. A
% float must fit on one page; these do not, and an oversized float is
% reported as "Float too large for page" and dropped. A longtable carries
% its own caption and label and must never be wrapped in \begin{table}.
% Column fractions sum to about 0.92, budgeting the inner \tabcolsep
% gutters; summing to 1.0 produces few-point overfull \hbox warnings. No
% bare l or c column sits beside an X column WHEN ITS CONTENT IS WIDE OR
% VARIABLE, because l takes its natural width and can squeeze X below its
% minimum ("X Columns too narrow"). A SHORT, FIXED l is sanctioned: A.11's
% three-family table uses {@{}lYY@{}} for labels of a few words and compiles
% with no X-column warning. The rule was stated absolutely and the appendix
% did not follow it; the enforced policy is the one written here.
% =====================================================================
\documentclass{article}
\usepackage[margin=1in]{geometry} % wider text block: removes table/line overflow from default-article narrow column
\usepackage{float}
\usepackage{longtable}
\usepackage{graphicx}
\usepackage{booktabs}
\usepackage{amsmath}
\usepackage{amssymb}
\usepackage{siunitx}
\usepackage{xcolor}
\usepackage{url} % for \url{} in the COMSOL bibliography entry
\usepackage{array}
\usepackage{tabularx} % wide tables forced to \textwidth (wrapping X columns) instead of overflowing
% Scalable fonts where available; microtype's font EXPANSION requires them, so
% expansion is enabled only when lmodern loads. Keeps the file portable to a
% minimal TeX installation without silently dropping the typographic benefit.
\IfFileExists{lmodern.sty}{%
  \usepackage{lmodern}\usepackage[T1]{fontenc}\usepackage{microtype}%
  % BOLD SMALL CAPS. Latin Modern has no T1/lmr/bx/sc, so any \textbf{...\textsc{...}}
  % (A.8's "On the EXTERNAL label" heading, and any future one) raises
  %   Font shape `T1/lmr/bx/sc' undefined ... defaults substituted
  % and silently drops either the bold or the small caps. lmr DOES carry the b series,
  % so substituting bx -> b gives real bold small caps instead of a fallback.
  \DeclareFontShape{T1}{lmr}{bx}{sc}{<->ssub*lmr/b/sc}{}%
  \DeclareFontShape{T1}{lmr}{bx}{si}{<->ssub*lmr/b/si}{}%
}{%
  \usepackage[expansion=false]{microtype}%
}
% THE SAME GAP EXISTS WITHOUT lmodern. On a minimal installation the encoding stays
% OT1/cmr, which has no bold small caps either, so the identical warning appears as
% `OT1/cmr/bx/sc'. Computer Modern has no bold small caps to substitute, so this falls
% back to medium small caps: the small caps survive and only the weight is lost, which
% is the better of the two degradations. Declared for both encodings because the file
% must build the same on Overleaf (lmodern present) and on a bare TeX Live.
\DeclareFontShape{OT1}{cmr}{bx}{sc}{<->ssub*cmr/m/sc}{}
\DeclareFontShape{T1}{cmr}{bx}{sc}{<->ssub*cmr/m/sc}{}% protrusion/expansion: clears most under/overfull \hbox warnings in prose
\newcolumntype{Y}{>{\raggedright\arraybackslash}X} % ragged, line-wrapping X column
\setlength{\emergencystretch}{3em} % last-resort glue so stubborn lines (DOIs etc.) do not run overfull

% --- lineage shorthand ---
\newcommand{\fgap}{f_{\mathrm{gap}}}
% --- cutoff-law (selectivity-axis) decomposition ---
\newcommand{\kapp}{\kappa_{\mathrm{app}}}
\newcommand{\kster}{\kappa_{\mathrm{steric}}}
\newcommand{\kg}{\kappa_{g}}                    % E8: prefactor paired with the CUBIC branch
\newcommand{\kF}{\kappa_{F_{\mathrm{lam}}}}     % E8: prefactor paired with the EXACT kernel
\newcommand{\Kin}{K_{\text{in-plane}}} % text-mode hyphen (no \text nested in \mathrm)
\newcommand{\Kdepth}{K_{\mathrm{depth}}}
\newcommand{\Kdth}{K_{\mathrm{depth}}^{\mathrm{th}}}    % E10: ANALYTIC depth factor
\newcommand{\Kdcfd}{K_{\mathrm{depth}}^{\mathrm{CFD}}}  % E10: MEASURED depth factor
\newcommand{\dcrit}{d_{\mathrm{crit}}}
\newcommand{\Ral}{R(\alpha)}          % interpretive alias of Kdepth, see collision warning
\newcommand{\skimw}{\delta}           % dividing-streamline (skim) width
\newcommand{\acal}{\alpha_{\mathrm{cal}}}       % E9: calibration aspect-ratio array
% E19: reduced size ratio d_crit/w_c.
% WAS u_{d}. That differed from the ASME V&V-20 data uncertainty u_D by SUBSCRIPT
% CASE ALONE. A collision note documents that trap without removing it, because a
% reader scanning a dense table cannot see case: the sigma note works only because
% sigma_F and sigma(i) differ by subscript CONTENT. The letter u already carries the
% velocities u(y), u_max (standard fluid mechanics) and the V&V-20 uncertainties
% u_val, u_num, u_input, u_D (ASME's own symbols); both are inherited notations, so
% the reduced size ratio was the only one free to move. d-tilde is apt for a
% normalised diameter and clashes with no other d in this appendix: the deprecated d,
% d_crit, D=2a, D_c and D_h are all distinct.
% The macro KEEPS ITS NAME, so no subsection body changes. symlint.py carries the
% mapping \ud -> dtilde so it stops reporting a u-collision that no longer exists.
% BEFORE RELYING ON THIS: grep the remaining subsections for a literal "u_d" written
% as prose rather than through \ud; those mentions will not follow the redefinition.
\newcommand{\ud}{\widetilde{d}}
% --- gap geometry (see the G/g collision note) ---
% G is the along-flow LINK WIDTH (the opening a particle passes through, Gifford's
% "wide (~20 um) gaps"); ell_g is the LINK LENGTH (the wall thickness the fluid
% crosses); g_p is the along-flow POST WIDTH. G and g_p are both along-flow, which is
% why they sum to a row pitch. Lower-case g is the CUBIC KERNEL g(x)=3x^2-2x^3 and is
% dimensionless; it is not a length, and the earlier gap row that used it is
% deprecated in favour of G.
% --- validation-layer (V&V-20) ---
\newcommand{\uval}{u_{\mathrm{val}}}
\newcommand{\unum}{u_{\mathrm{num}}}
\newcommand{\uinput}{u_{\mathrm{input}}}
\newcommand{\uD}{u_{D}}


% ---------------------------------------------------------------------------
% \xref{label}{fallback} -- a cross-reference that degrades gracefully.
%
%   standalone build (target not in this document)
%       prints `fallback'. No warning, no `??', nothing to clean up later.
%   assembled dissertation (target present)
%       prints the live \ref, so renumbering Appendix A is still absorbed.
%
% The body text is written ONCE and is correct in both builds. Use it for every
% pointer that leaves the current file -- Appendix A results, chapter sections,
% glossary subsections not yet written.
%
% Two-pass note: on the FIRST pdflatex pass of an assembled build the .aux has
% not been read, so the fallback prints; the second pass resolves it. That is
% ordinary LaTeX behaviour and emits no warning.
% ---------------------------------------------------------------------------
\makeatletter
\providecommand{\xref}[2]{\@ifundefined{r@#1}{#2}{\ref{#1}}}
\makeatother

% =====================================================================
% VALUES ARE GENERATED, NOT TYPED.
% The macro file loaded below is emitted from the verified derivation
% record and pinned to its content hash. A number typed inline instead of
% cited from a macro is a number that will be stale after the next
% correction: the change log records three such corrections that each had
% to be chased through four places by hand. The fallback block keeps this
% file compilable when the generated file is absent; it is a SNAPSHOT and
% the generated file is authoritative.
%
% The snapshot below was re-verified against the derivation record on this
% revision: all 40 constants reproduce, including the three calibration
% depth factors, which use the EXACT ratios 80/100, 80/125 and 80/150
% rather than the rounded aspect ratios printed in the prose (0.811 at the
% widest channel is K_depth(80/150)=0.810929, not K_depth(0.53)=0.809540).
% =====================================================================
\IfFileExists{tex/notation_values.tex}{\input{tex/notation_values}}{%
  \newcommand{\kappVal}{4.9138}\newcommand{\kappFit}{4.9137}%
  \newcommand{\kappFitFull}{4.913717814706706}%
  \newcommand{\kappCIlo}{4.75}\newcommand{\kappCIhi}{5.08}\newcommand{\kappCV}{4.0}%
  \newcommand{\kstericVal}{2}\newcommand{\KinVal}{1.000}%
  \newcommand{\kappOverRootThree}{2.8370}%
  \newcommand{\plateauKappa}{4.91}\newcommand{\footKappa}{1.34}%
  \newcommand{\plateauN}{8}\newcommand{\footN}{3}%
  \newcommand{\kappaSpanLo}{1.20}\newcommand{\kappaSpanHi}{5.20}%
  \newcommand{\KdepthAnalyticShallow}{0.8879}\newcommand{\KdepthAnalyticMid}{0.9500}%
  \newcommand{\KdepthAnalyticDeep}{0.9794}%
  \newcommand{\KdepthCfdShallow}{0.891}\newcommand{\KdepthCfdMid}{0.945}%
  \newcommand{\KdepthCfdDeep}{0.981}%
  \newcommand{\KdepthCalNarrow}{0.888}\newcommand{\KdepthCalMid}{0.849}%
  \newcommand{\KdepthCalWide}{0.811}%
  \newcommand{\residualFolded}{2.4569}\newcommand{\residualSepShallow}{2.7574}%
  \newcommand{\residualSepMid}{2.5999}\newcommand{\residualSepDeep}{2.5044}%
  \newcommand{\stericFloorShallow}{1.782}%
  \newcommand{\FclosedVone}{0.5592}\newcommand{\FclosedVthree}{0.7189}%
  \newcommand{\FcomsolVone}{0.5070}\newcommand{\FcomsolVthree}{0.6281}%
  \newcommand{\FexpVone}{0.505}\newcommand{\FexpVthree}{0.649}%
  \newcommand{\FsdVone}{0.033}\newcommand{\FsdVthree}{0.036}%
  \newcommand{\alphaThreeDVone}{0.811}\newcommand{\alphaThreeDVthree}{0.660}%
  \newcommand{\ReWorstChannel}{5.2}\newcommand{\ReGapScale}{5.2\times10^{-3}}%
  \newcommand{\fgapBandLo}{3.63\times10^{-5}}\newcommand{\fgapBandHi}{1.93\times10^{-3}}%
  \newcommand{\fgapDesignCap}{5.76\times10^{-4}}%
}
% NOTE on \alphaThreeDVone: it is 0.811, the same three digits as
% \KdepthCalWide. They are DIFFERENT quantities that coincide numerically --
% the 3-D outlet resistance-ratio correction at V1, and the depth factor at
% the widest calibration channel. Do not let one be edited into the other.

% =====================================================================
% A.4 KERNEL-PAIRING VALUES.
% Generated by gen_kernel_pairing.py from the verified derivation record.
% Loaded the same way as notation_values: the generated file is
% AUTHORITATIVE, and the fallback below is a SNAPSHOT that keeps this file
% compilable when the generated file is absent.
%
% Without this block, A.4 raises "Undefined control sequence" on
% \kappSuppliedGap, \KdepthDev*, \ctwobar, \kFval and 29 others.
% Regenerate after any change to c2(), K_depth(), F_lam() or the plateau fit:
%   python3 gen_kernel_pairing.py tex/kernel_pairing_values.tex
% =====================================================================
\IfFileExists{tex/kernel_pairing_values.tex}{\input{tex/kernel_pairing_values}}{%
  \newcommand{\ctwobar}{4.177}%
  \newcommand{\kFval}{5.798}%
  \newcommand{\depthEff}{0.8475}%
  \newcommand{\kFmeanK}{5.7862}%
  \newcommand{\kFmeanKerr}{0.21}%
  \newcommand{\kFoneAlpha}{5.5342}%
  \newcommand{\kFoneAlphaErr}{4.6}%
  \newcommand{\ctwobarDecl}{4.182}%
  \newcommand{\kFdecl}{5.802}%
  \newcommand{\transferExact}{19.9}%
  \newcommand{\transferDecl}{20.3}%
  \newcommand{\dcritGGleadIn}{2.170}%
  \newcommand{\dcritGGleadOut}{1.504}%
  \newcommand{\choiceLeadGap}{0.2}%
  \newcommand{\crossUp}{18.0}%
  \newcommand{\crossDown}{-15.3}%
  \newcommand{\alphaDIn}{1.867}%
  \newcommand{\dcritGGIn}{2.174}%
  \newcommand{\dcritFFIn}{2.482}%
  \newcommand{\dcritFgIn}{2.565}%
  \newcommand{\dcritgFIn}{2.103}%
  \newcommand{\choiceIn}{14.1}%
  \newcommand{\alphaDOut}{2.692}%
  \newcommand{\dcritGGOut}{1.507}%
  \newcommand{\dcritFFOut}{1.743}%
  \newcommand{\dcritFgOut}{1.779}%
  \newcommand{\dcritgFOut}{1.477}%
  \newcommand{\choiceOut}{15.6}%
  \newcommand{\KdepthDevShallow}{0.349}%
  \newcommand{\KdepthDevMid}{0.530}%
  \newcommand{\KdepthDevDeep}{0.160}%
  \newcommand{\KdepthDevMidTh}{0.527}%
  \newcommand{\KdepthDevMax}{0.531}%
  \newcommand{\kappSuppliedGap}{1.7\times10^{-5}}%
  \newcommand{\dinhBandLo}{1.25}%
  \newcommand{\dinhBandHi}{2.17}%
  \newcommand{\dinhKappaLocal}{4.08}%
  \newcommand{\dinhCollapseLo}{1.80}%
  \newcommand{\dinhKneeRatio}{1.21}%
  \newcommand{\dinhBandWidth}{1.74}%
  \newcommand{\kappCIhalfPct}{3.4}%
  \newcommand{\xiaKappaA}{5.03}%
  \newcommand{\xiaKappaB}{5.16}%
  \newcommand{\xiaKappaC}{5.32}%
  \newcommand{\xiaKappaMean}{5.17}%
  \newcommand{\xiaAgreePct}{5.2}%
  \newcommand{\depthCorrPct}{5.3}%
  \newcommand{\chiGlobal}{30.6}%
  \newcommand{\chiPlateau}{3.8}%
  \newcommand{\chiRatio}{8.0}%
  \newcommand{\etaShortfallPct}{9.7}%
  \newcommand{\etaVal}{0.903}%
  \newcommand{\runoutFmatched}{0.519}%
  \newcommand{\runoutFmismatched}{0.265}%
  \newcommand{\residPctVOne}{9.3}%
  \newcommand{\residPctVThree}{12.6}%
  \newcommand{\squareDuctR}{28.4542}%
  \newcommand{\skimErrDesign}{0.47}%
  \newcommand{\skimErrBand}{0.86}%
  \newcommand{\DhOutlet}{75.8}%
  \newcommand{\DhInlet}{97.7}%
  \newcommand{\Uoutlet}{0.069}%
  \newcommand{\ReHighFlow}{16}%
  \newcommand{\RePRP}{2.6}%
  \newcommand{\LeInlet}{60.5}%
  \newcommand{\LeOverLamInlet}{1.23}%
  \newcommand{\LeOverLamGInlet}{1.31}%
  \newcommand{\LeOutlet}{46.9}%
  \newcommand{\LeOverLamOutlet}{0.96}%
  \newcommand{\LeOverLamGOutlet}{1.02}%
  \newcommand{\aTwelveFL}{1.42}%
  \newcommand{\DTwelveFL}{2.84}%
  \newcommand{\QgapDinh}{3.12}%
  \newcommand{\dldExcessLo}{1.23}%
  \newcommand{\dldExcessHi}{1.33}%
  \newcommand{\ihRatioA}{1.44}%
  \newcommand{\ihRatioB}{1.48}%
  \newcommand{\ihRatioC}{1.21}%
  \newcommand{\ihRatioLo}{1.21}%
  \newcommand{\ihRatioHi}{1.48}%
  \newcommand{\ihRatioMean}{1.4}%
  \newcommand{\davisKappaHi}{2.659}%
  \newcommand{\davisKappaLo}{2.459}%
  \newcommand{\inglisKappaLo}{2.040}%
  \newcommand{\inglisKappaHi}{2.449}%
  \newcommand{\dldKappaLo}{2.040}%
  \newcommand{\dldKappaHi}{2.659}%
  \newcommand{\dldConvergePct}{0.4}%
  \newcommand{\expBeta}{0.51}%
  \newcommand{\expGamma}{1.02}%
  \newcommand{\plateauUwc}{0.035}%
  \newcommand{\footUwc}{0.012}%
  \newcommand{\betaCIlo}{0.41}%
  \newcommand{\betaCIhi}{0.61}%
  \newcommand{\gammaCIlo}{0.84}%
  \newcommand{\gammaCIhi}{1.20}%
  \newcommand{\hillKmax}{5.568}%
  \newcommand{\hillUhalf}{0.01651}%
  \newcommand{\hillN}{1.791}%
  \newcommand{\dinhSaddleFgap}{1.041\times10^{-4}}%
  \newcommand{\dinhSaddleU}{0.02894}%
  \newcommand{\sibALo}{1.02}%
  \newcommand{\sibAHi}{1.47}%
  \newcommand{\sibBLo}{1.26}%
  \newcommand{\sibBHi}{1.81}%
  \newcommand{\dinhBracketLo}{1.0}%
  \newcommand{\dinhBracketHi}{4.7}%
}

% =====================================================================
% A.6 VALIDATION DISTANCES. Generated by gen_validation_values.py from the
% same F values this preamble carries, so the document and the generator
% cannot disagree about what F is. Fallback is a SNAPSHOT.
% =====================================================================
\IfFileExists{tex/validation_values.tex}{\input{tex/validation_values}}{%
  \newcommand{\EsdClosedVOne}{+1.64}%
  \newcommand{\EsdComsolVOne}{+0.06}%
  \newcommand{\EsemComsolVOne}{+0.12}%
  \newcommand{\uvalExcessVOne}{0.048}%
  \newcommand{\EsdClosedVThree}{+1.94}%
  \newcommand{\EsdComsolVThree}{-0.58}%
  \newcommand{\EsemComsolVThree}{-1.16}%
  \newcommand{\uvalExcessVThree}{0.011}%
  \newcommand{\uDoverUnumVOne}{660}%
  \newcommand{\uDoverUnumVThree}{720}%
  \newcommand{\uDoverUnum}{660}%
  \newcommand{\uvalExcessMax}{0.05}%
  \newcommand{\FexcessPctVOne}{10.7}%
  \newcommand{\FexcessPctVThree}{10.8}%
  \newcommand{\dFrecurConst}{2.35\times10^{-4}}%
  \newcommand{\dFrecurVar}{1.97\times10^{-5}}%
  \newcommand{\peMarginMin}{48}%
  \newcommand{\peHalfAgreeVOne}{2.3}%
  \newcommand{\peHalfAgreeVThree}{4.7}%
  \newcommand{\netEmergentPct}{0.12}%
  \newcommand{\siblingSqrtcPct}{5}%
}
\title{CIF Variable, Unit, and Equation Reference}
\date{}

\begin{document}

\maketitle

\appendix
\section{Variable and Unit Reference}
\label{app:units}

This appendix catalogs the symbols, dimensionless quantities, and
key equations of the dissertation's analytical, design, and
validation framework, with SI units, the practical units used in the CIF
lineage, and step-by-step dimensional derivations. It spans both
performance axes, the filtrate fraction (yield) and the particle-size
cutoff (selectivity), together with the verification-and-validation layer
that connects the analytical framework to numerical (COMSOL) and
published-device evidence. Direct quotations from the source papers
anchoring each unit convention and each adjudicated physical criterion are
collected in Section~\ref{app:units_quotes}; the selectivity axis is
validated in Section~\ref{app:units_selval} and the yield axis in
Section~\ref{app:units_validation}; an independent cross-lineage dataset
is given in Section~\ref{app:units_crosslineage}; the design envelope is
inventoried in Section~\ref{app:units_envelope}; and a consolidated master
cheat-sheet is in Section~\ref{app:units_cheatsheet}.
\newline



%==============================================================
% A.1 -- GEOMETRIC VARIABLES  (complete, corrected)
% Every locator in this subsection is verified against per-page OCR by
% locatorcheck.py (validate.sh gate 8b).
%==============================================================
\subsection{Geometric Variables}
\label{app:units_geom}

Table~\ref{tab:units_geom} fixes the geometric symbols and the in-plane
versus out-of-plane convention. $G$, $\Lambda$ and $\phi$ are included here because the
wall porosity is one of the three model gaps to which the cutoff residual
is attributed (Section~\xref{app:units_cutoff}{``Cutoff-Law Prefactor Decomposition''}).

\paragraph{Evidentiary markers.} Where a row's status is not obvious it is
marked \textsc{external} (an input the record does not derive: a source
states it, or this work supplies it and no validator discharges it),
\textsc{derived} (it follows from geometry or from an Appendix~A result), or
\emph{This work} (a convention this document sets). Unmarked rows are
\textsc{external} and carry their locator in the Source column. The status
columns of later tables use five more markers: \textsc{fitted} (calibrated
to data by a fit), \textsc{numerical} (computed by a numerical solve),
\textsc{measured} (set by a measurement), \textsc{proposed} (a
parameterization this work proposes and has not yet computed or measured)
and \textsc{bounds} (a condition that limits where the design relations
hold); \textsc{establishes} and \textsc{competes}, in the name column of
Table~\xref{tab:units_cutoff}{A.4}, say what a quantity does in the argument.

% longtable, NOT a float: this table is taller than a page, and a `table` float
% that does not fit is reported as "Float too large for page" and dropped.
% Column fractions sum to 0.92, budgeting 8 x \tabcolsep of inner gutter at 4pt;
% summing to 1.0 is what produces the few-point overfull \hbox warnings.
\begingroup\small\setlength{\tabcolsep}{4pt}
\begin{longtable}{@{}>{\raggedright\arraybackslash}p{0.115\textwidth}>{\raggedright\arraybackslash}p{0.44\textwidth}>{\raggedright\arraybackslash}p{0.065\textwidth}>{\raggedright\arraybackslash}p{0.105\textwidth}>{\raggedright\arraybackslash}p{0.165\textwidth}@{}}
\caption{Geometric variables used to specify CIF device architecture. Practical units reflect the CIF lineage convention; SI units are used in all dimensional derivations.}\label{tab:units_geom}\\
\toprule Symbol & Quantity & SI unit & Practical unit & Source \\\midrule\endfirsthead
\multicolumn{5}{@{}l}{\footnotesize\itshape Table \thetable\ (continued)}\\
\toprule Symbol & Quantity & SI unit & Practical unit & Source \\\midrule\endhead
\bottomrule\endlastfoot
$w_c$           & Center-channel width (constant-$w_c$ regime)            & \si{\meter} & \si{\micro\meter} & \cite[p.~4498, Fig.~1]{b1} \\
$w_c(i)$        & Center-channel width of segment $i$, after gap row $i$ (tapered-$w_c$ regime; $w_c(0)$ is the inlet width). Xia's width at gap $i$ is read as this $w_c(i)$ in Sections~\xref{app:units_selval}{A.7} and~\xref{app:units_equations}{A.9}, as in App.~A; the width arriving at row $i$ is $w_c(i-1)$, one taper step away & \si{\meter} & \si{\micro\meter} & \cite[p.~916]{b5}; \cite[p.~10 (Sci Rep PDF pagination)]{b2} \\
$w_s$, $w_s(i)$ & Side-channel width (generic / of segment $i$, after gap row $i$, the width that sets $R_s(i)$) & \si{\meter} & \si{\micro\meter} & \cite[p.~4498, Fig.~1]{b1} \\
$w_c(\text{ref})$ & Reference center-channel width (Xia 2022 scaling)    & \si{\meter} & \si{\micro\meter} & \cite[p.~1]{b3} \\
$w_{\mathrm{ch}}$ & \emph{Generic} channel width placeholder, $w_{\mathrm{ch}}\!\in\!\{w_c,w_s\}$, used only inside the resistance law (Eqn~\xref{eqn:gifford5}{5}) & \si{\meter} & \si{\micro\meter} & This work \\
$w$             & \emph{Diverted fluid lamina width} (empirical extracted-lamina width; cutoff reference length) & \si{\meter} & \si{\micro\meter} & \cite[p.~917]{b5} \\
$\skimw$        & Dividing-streamline (skim) width, $\skimw=w_c\sqrt{\fgap/3}$ at leading order; plays the role of $w$ in the cutoff law. \textsc{derived} (App.~A, Result~\xref{res:skim-inversion}{2}) & \si{\meter} & \si{\micro\meter} & \cite{b12a}\,(\texttt{delta\_skim}); \cite[p.~656]{b17} \\
$b$             & First-streamline (boundary-lamina) width; Inglis writes $D_c=2\beta$, rendered $b$ here (see Lamina-width collision below) & \si{\meter} & \si{\micro\meter} & \cite[p.~656, Eqn~1]{b17} \\
$h$             & Channel depth (out-of-plane height); Gifford's $d$ in the Eqn~5 sentence (``$d$ is the depth of the channel'') & \si{\meter} & \si{\micro\meter} & \cite[p.~4498]{b1} \\
$d$             & \textcolor{gray}{\emph{[deprecated]} Gifford depth notation $\equiv h$; retained for cross-reference only. Use $h$ to avoid collision with $\dcrit$.} & \si{\meter} & \si{\micro\meter} & \cite[p.~4498]{b1} \\
$L_{\mathrm{dev}}$ & Device length (along flow direction). \emph{Renamed from $L$}: Gifford's $L$ is the \emph{segment} length, not the device length & \si{\meter} & \si{\milli\meter} & This work \\
$L_c(i), L_s(i)$ & Center/side channel \emph{segment} length after gap row $i$ (between rows $i$ and $i+1$; segment $0$ is the inlet stretch, segment $N$ the outlet), the segment that carries $Q_c(i)$ and $Q_s(i)$. Gifford describes a segment as lying between rows $i-1$ and $i$, one row earlier, which is immaterial for his constant $\fgap$; the quantity Gifford's Eqn~5 calls $L$, clarified there as ``the size (along the direction of flow) of an obstacle in the device'' & \si{\meter} & \si{\micro\meter} & This work, after Gifford's Eqn~5 $L$ (\cite[p.~4498]{b1}; clarified \cite[p.~4499]{b1}) \\
$g$             & \textcolor{gray}{\emph{[deprecated]} An earlier draft used $g$ for the gap opening. That is the same quantity as $G$ below, and lower-case $g$ is needed for the cubic kernel; use $G$.} & \si{\meter} & \si{\micro\meter} & --- \\
$G$             & \textbf{Link width}: the along-flow width of the opening in the wall between center and side channel. This is the dimension a particle passes through, and the one Gifford reports as ``wide ($\sim$20~\si{\micro\meter}) gaps'' & \si{\meter} & \si{\micro\meter} & This work; value \cite[p.~4498]{b1} \\
$\ell_g$        & \textbf{Link length}: the wall thickness the fluid crosses, i.e.\ the flow path \emph{through} the opening. With $G$ it sets the gap cross-section, and hence the gap resistance $R_{\mathrm{gap}}$ of Table~\xref{tab:units_fluid}{A.2}; the hydraulic diameter of the opening is $D_h=2Gh/(G+h)$ & \si{\meter} & \si{\micro\meter} & This work \\
$g_p$           & \textbf{Post width}: the along-flow width of the wall segment between successive openings. $G$ and $g_p$ are both along-flow, which is why they sum to a row pitch & \si{\meter} & \si{\micro\meter} & This work \\
$\Lambda$       & Row pitch, $\Lambda(i)\ge\max(G+g_p,\,2p\ \text{[if folded]})$; equality only for an unfolded tight-packed row. A \emph{design constraint}, not a definition: the bound is set by the minimum feature and fold geometry & \si{\meter} & \si{\micro\meter} & This work \\
$\phi$          & \textbf{Wall porosity}, $\phi=G/\Lambda$: the open fraction of the ``wall'' the resistance kernels assume solid. \textsc{derived} from geometry. One of the \emph{three} model gaps the cutoff residual is measured against (Section~\xref{app:units_cutoff}{``Cutoff-Law Prefactor Decomposition''}). Distinct from $\phi_{\mathrm{cal}}$, which is \textsc{external} (Section~\xref{app:units_envelope}{``Validity Envelope''}), and from the fold factor $\Phi(i)$ & --- & --- & This work; \cite{b12a} \\
$N$             & Total number of gap rows. Gifford's symbol is $i_T$ (Fig.~2 caption, ``the necessary total number of rows of gaps, $i_T$''); $N$ is this work's, adopted to free $i$ for the row index & --- & --- & This work, after \cite[p.~4499]{b1} \\
$i$             & Gap row index, $1 \le i \le N$; flows and segments are labelled $0 \le i \le N$, $0$ being the inlet stretch & ---          & ---                 & \cite[p.~4498, Fig.~1]{b1} \\
\end{longtable}\endgroup

\textbf{Convention.} $w_c$, $w_s$ (and the generic placeholder
$w_{\mathrm{ch}}$) are the in-plane channel widths (perpendicular to
flow); $h$ is the out-of-plane channel depth.

\textbf{Lamina-width collision.} The bare symbol $w$ denotes
the \emph{diverted fluid lamina width} only. The rectangular-duct
resistance law (Eqn~\xref{eqn:gifford5}{5}) takes the \emph{channel} width as
argument; that argument is written $w_{\mathrm{ch}}$ here, never bare $w$.
The theoretical lamina width is $\skimw=w_c\sqrt{\fgap/3}$; the empirical
lamina width is $w$; the Inglis first-streamline width is $b$. Inglis
writes this as $D_c=2\beta$ with Greek $\beta$ (\cite[p.~656, Eqn~1]{b17});
it is rendered $b$ here to free $\beta$ for the cutoff scaling exponent.
These three denote the same boundary lamina, and their identification
$b\equiv\skimw$ is a \emph{proven} result, not a definition: the CIF skim
cubic is algebraically identical to the Inglis DLD cubic
(App.~A, Result~\xref{res:cubic-identity}{4}; recorded as a proven node in the
argument map, App.~B).

\textbf{$G$/$g$ collision: $g$ is the cubic kernel, not a length.}
Lower-case $g$ denotes the
skim kernel $g(x)\equiv3x^2-2x^3$ and rides as a subscript on the paired
prefactor $\kappa_g$. It is dimensionless. It is \emph{not} a gap width: the
opening is $G$, its through-wall length is $\ell_g$, and the post is $g_p$;
$G$ and $g$ differ only by case and denote different kinds of thing, so a
reader scanning a formula must read the case.
An earlier draft of this table used $g$ for the opening, which both duplicated
$G$ and collided with the kernel; that row is deprecated above. The gap
geometry needs no lower-case $g$: $\Lambda(i)\ge\max(G+g_p,\,2p)$ and
$\phi=G/\Lambda$ are stated entirely in $G$, $g_p$ and $\Lambda$.

\textbf{Two further symbol collisions, declared.} The lower-case $\phi$
above is the wall porosity; the upper-case $\Phi(i)=r(w_s(i))/r(\mathrm{GS})$
used in the folded-channel geometry is the \emph{fold factor} and is a
different quantity. The fold pitch $p=\mathrm{GS}+\max(F_{\min},2R_{\min})$
of that same geometry ($\mathrm{GS}$ the fold width, $F_{\min}$ the minimum
feature and $R_{\min}$ the minimum inner-corner radius, as the design
specification's symbol table defines them) shares a letter with the hydrostatic pressure $p$ of
Table~\xref{tab:units_fluid}{A.2}; pressure is always written $P$ or $p$ in a
pressure context and the fold pitch only inside that folded-channel geometry.
 

 
 

 
%==============================================================
% A.2 -- FLUID AND HYDRAULIC VARIABLES  (corrected)
% Every locator verified against per-page OCR (locatorcheck.py, gate 8b).
%==============================================================
\subsection{Fluid and Hydraulic Variables}
\label{app:units_fluid}

Table~\ref{tab:units_fluid} lists the flow, pressure and resistance
variables, together with the two external inputs $\sigma(i)$ and
$\sigma_0$. Two letters are overloaded and are disambiguated in notes below
the table: $R$ carries four unrelated meanings across this appendix, and $P$
is both the pressure and the normalised march flow $P(i)$ within this one
table.

% longtable, NOT a float: consistent with Table~\ref{tab:units_geom}, and safe if
% the table grows. Fractions sum to 0.92, budgeting 8 x \tabcolsep of inner gutter
% at 4pt; summing to 1.0 is what produces few-point overfull \hbox warnings.
% A bare `l' column beside an `X' is also avoided: `l' takes its natural width and
% can squeeze X below its minimum ("X Columns too narrow").
\begingroup\small\setlength{\tabcolsep}{4pt}
\begin{longtable}{@{}>{\raggedright\arraybackslash}p{0.115\textwidth}
>{\raggedright\arraybackslash}p{0.40\textwidth}
>{\raggedright\arraybackslash}p{0.115\textwidth}
>{\raggedright\arraybackslash}p{0.115\textwidth}
>{\raggedright\arraybackslash}p{0.175\textwidth}@{}}
\caption{Fluid-mechanical and hydraulic variables.}\label{tab:units_fluid}\\
\toprule Symbol & Quantity & SI unit & Practical unit & Source \\\midrule\endfirsthead
\multicolumn{5}{@{}l}{\footnotesize\itshape Table \thetable\ (continued)}\\
\toprule Symbol & Quantity & SI unit & Practical unit & Source \\\midrule\endhead
\bottomrule\endlastfoot
$Q_c$, $Q_c(i)$       & Center-channel volumetric flow rate after gap row $i$            & \si{\cubic\meter\per\second} & \si{\micro\liter\per\minute} & \cite[p.~4498]{b1} \\
$Q_s$, $Q_s(i)$       & Side-channel volumetric flow rate (one side, after gap $i$) & \si{\cubic\meter\per\second} & \si{\micro\liter\per\minute} & \cite[p.~4498]{b1} \\
$Q_{\mathrm{gap}}(i)$ & Volumetric flow through one gap on one side at row $i$    & \si{\cubic\meter\per\second} & \si{\nano\liter\per\minute} & \cite[p.~915]{b5} \\
$Q_{\mathrm{in}}$     & Total inlet flow rate. The lineage devices are operated at $30~\si{\micro\liter\per\minute}$; the \emph{symbol} is this work's & \si{\cubic\meter\per\second} & \si{\micro\liter\per\minute} & This work; rate \cite[p.~915]{b5} \\
$Q_m(i)$              & Middle-channel flow rate at gap $i$ (Dinh notation; $=Q_c(i-1)$)  & \si{\cubic\meter\per\second} & \si{\micro\liter\per\minute} & \cite[p.~915]{b5} \\
$P$, $p$              & Hydrostatic pressure. Gifford writes only the differential $\Delta P$; the bare symbol is this work's & \si{\pascal} & \si{\kilo\pascal} or psi & This work \\
$\Delta P$, $\Delta P(i)$ & Pressure drop across one channel segment              & \si{\pascal} & \si{\kilo\pascal} & \cite[p.~4498]{b1} \\
$p_{\text{center}}(L_{\mathrm{dev}})$ & Pressure in center channel at the outlet & \si{\pascal} & \si{\kilo\pascal} & Condition \cite[p.~4500]{b1}; symbol this work \\
$p_{\text{side}}(L_{\mathrm{dev}})$ & Pressure in side channel at the outlet & \si{\pascal} & \si{\kilo\pascal} & Condition \cite[p.~4500]{b1}; symbol this work \\
$R_c$, $R_c(i)$       & Hydraulic resistance of one center-channel segment        & \si{\pascal\second\per\cubic\meter} & --- & \cite[p.~4498]{b1} \\
$R_s$, $R_s(i)$       & Hydraulic resistance of one side-channel segment          & \si{\pascal\second\per\cubic\meter} & --- & \cite[p.~4498]{b1} \\
$R_{\mathrm{gap}}$    & Hydraulic resistance of one gap. Gifford describes the gaps as presenting ``a relatively small amount of resistance'' but defines no symbol; this one is this work's & \si{\pascal\second\per\cubic\meter} & --- & This work, after \cite[p.~4498]{b1} \\
$P(i)$, $S(i)$        & Normalised center and side flows along the march, $P(i)=Q_c(i)/Q_c(0)$ and $S(i)=Q_s(i)/Q_c(0)$. \emph{Not} the pressure $P$ above: see the collision note & --- & --- & This work; \cite{b12a} \\
$\sigma(i)$, $\sigma_0$ & Cumulative side-flow ratio along the march, $\sigma(i)=S(i)/P(i)$; $\sigma_0$ is the (fictitious) inlet side-flow. \textsc{external}: enters the mass identity $P(i)+2S(i)=1+2\sigma_0$ and is not derived here & --- & --- & This work; \cite{b12a} \\
$\mu$                 & Dynamic viscosity of fluid                                & \si{\pascal\second} & \si{\milli\pascal\second} & \cite[p.~4498]{b1} \\
$\rho$                & Fluid density. Water is taken as $1000~\si{\kilo\gram\per\cubic\meter}$ and platelet-rich plasma as $1025~\si{\kilo\gram\per\cubic\meter}$; neither value is measured here & \si{\kilo\gram\per\cubic\meter} & \si{\kilo\gram\per\cubic\meter} & This work (assumed) \\
$u(y)$, $v$           & Local fluid velocity                                      & \si{\meter\per\second} & \si{\milli\meter\per\second} & Standard \\
$u_{\max}$            & Peak velocity in parabolic profile                        & \si{\meter\per\second} & \si{\milli\meter\per\second} & Standard \\
\end{longtable}\endgroup

\textbf{R-symbol collision.} Four distinct objects use the
letter $R$ in this dissertation: the hydraulic resistances
$R_c, R_s, R_{\mathrm{gap}}$ (this table); the depth-correction
\emph{function} $\Ral$ (Section~\xref{app:units_cutoff}{``Cutoff-Law
Prefactor Decomposition''}); the Reynolds number $\mathrm{Re}$
(Table~\xref{tab:units_dimensionless}{A.3}); and the multi-variant ratio
$\mathcal{R}$ (Eqn~\xref{eqn:multi_variant}{16}). To keep them separate,
the depth correction is written $\Kdepth$ throughout, with $\Ral$ used
only as its interpretive alias when the aspect-ratio dependence is being
emphasized; $\mathcal{R}$ is always script.

\textbf{$P$-symbol collision.} The pressure $P$ of this table and the
normalised center flow $P(i)$ are different quantities that share a letter
within this same table. Pressure is always written with a $\Delta$ or a subscript
($\Delta P$, $p_{\text{center}}$, $p_{\text{side}}$) or as the bare $P$ with
no argument; the march quantity is \emph{always} written with its row index,
$P(i)$, and is paired with $S(i)$. The terminal filtrate fraction is
$F=1-P(N)$, never $1-P$.

\textbf{Outlet pressure equilibration (P1.4).} The two outlet pressures
$p_{\text{center}}(L_{\mathrm{dev}})$ and $p_{\text{side}}(L_{\mathrm{dev}})$
are taken equal. This is not an assumption introduced here: the lineage
source states the resistance-ratio calculation holds ``provided the
hydrostatic pressures in the output collection areas are kept equal''
(\cite[p.~4500]{b1}). It is the premise of the closed-form filtrate
identity (Eqn~\xref{eqn:closed_form}{12}), stated as P1.4 in the formal
argument of Section~\xref{app:units_quotes}{A.10}; App.~B's C4 carries it as
premise~(3) of P-C4.1.

\textbf{Unit equivalences.}
\begin{align*}
1~\si{\micro\liter\per\minute} &= 1.667 \times 10^{-11}~\si{\cubic\meter\per\second} \\
1~\si{\nano\liter\per\minute}  &= 1.667 \times 10^{-14}~\si{\cubic\meter\per\second} \\
1~\si{\kilo\pascal}            &= 1000~\si{\pascal} = 0.145~\text{psi} \\
1~\text{psi}                   &= 6894.76~\si{\pascal} \\
1~\si{\milli\pascal\second}    &= 10^{-3}~\si{\pascal\second}
\end{align*}
%==============================================================
% A.3 -- DIMENSIONLESS QUANTITIES  (corrected)
% Every locator verified against per-page OCR (locatorcheck.py, gate 8b).
% Every b12a key verified to exist in the derivation record.
%==============================================================
\subsection{Dimensionless Quantities}
\label{app:units_dimensionless}

Table~\ref{tab:units_dimensionless} collects the dimensionless groups. Two
entries are easily misread: $\fgap$ is defined \emph{per side and per gap},
which is the origin of the factor of two in the continuity relation, and $F$
has two distinct forms with different explanatory roles
(Section~\xref{app:units_equations}{``Equations: Step-by-Step Unit
Derivations''}). Three symbol collisions are declared in notes below the
table.

\begingroup\small\setlength{\tabcolsep}{4pt}
\begin{longtable}{@{}>{\raggedright\arraybackslash}p{0.125\textwidth}
>{\raggedright\arraybackslash}p{0.245\textwidth}
>{\raggedright\arraybackslash}p{0.315\textwidth}
>{\raggedright\arraybackslash}p{0.235\textwidth}@{}}
\caption{Dimensionless ratios and design parameters. All entries are strictly
dimensionless; SI and practical columns are omitted.}\label{tab:units_dimensionless}\\
\toprule Symbol & Name & Defining relation & Source \\\midrule\endfirsthead
\multicolumn{4}{@{}l}{\footnotesize\itshape Table \thetable\ (continued)}\\
\toprule Symbol & Name & Defining relation & Source \\\midrule\endhead
\bottomrule\endlastfoot
$\fgap$, $\fgap(i)$ & Per-gap filtration fraction (per-side, per-gap) & $Q_{\mathrm{gap}}(i) = \fgap(i) \cdot Q_c(i-1)$ & \cite[p.~915]{b5}, \cite[p.~4498]{b1} \\
$\fgap(\text{ref})$ & Reference $\fgap$ in the Xia 2022 corrected scaling & Xia 2022 corrected Eqn~1 & \cite[p.~1]{b3} \\
$\fgap^*$ & Empirical threshold $\fgap$ for cutoff (read visually) & ``the filtration fraction per gap below which the beads of a specific diameter are consistently'' retained & \cite[p.~4500]{b1} \\
$F_{\mathrm{outlet}}$ & Filtrate fraction from the outlet divider & $F_{\mathrm{outlet}}=2R_c/(R_s+2R_c)$; \emph{contains no} $\fgap$ & Eqn~\xref{eqn:closed_form}{12}; this work, under the outlet condition of \cite[p.~4500]{b1} \\
$F_{\mathrm{recursion}}$ & Filtrate fraction from the per-gap march & $F_{\mathrm{recursion}}=1-\prod_i(1-2\fgap(i))$ & Eqn~\xref{eqn:F_product}{13}; derived from Eqn~3 of \cite[p.~4498]{b1} \\
$F_{\mathrm{V1}}$, $F_{\mathrm{V3}}$ & Filtrate fraction at the Dinh V1 and V3 variants & Closed form, paper, or COMSOL & \cite[p.~919]{b5} \\
$\sigma_F$ & One-SD experimental uncertainty in $F$. \emph{Not} the side-flow ratio $\sigma(i)$ of Table~\xref{tab:units_fluid}{A.2}: see the collision note & Propagated from measurement & \cite[p.~919]{b5} \\
$\alpha=h/w_{\mathrm{ch}}$ & Channel aspect ratio & --- & This work; the resistance series it enters is \cite[p.~875]{b6} \\
$\ud$ & Reduced size ratio, $\ud=\dcrit/w_c$ & Sets the plateau ($\ud\ge0.035$) and foot ($\ud<0.012$) bins & \cite{b12a}\,(\texttt{kappa\_bins}) \\
$x$ & Reduced skim width, $x=\skimw/w_c$ & root of $\fgap = 3x^2 - 2x^3$ & \cite{b12a}\,(\texttt{skim\_invert}); \cite[p.~656]{b17} \\
$\beta$ & Cutoff $\fgap$-scaling exponent & $\beta=\partial\ln\dcrit/\partial\ln\fgap=\tfrac12$ (empirical plateau fit $\approx0.51$) & \cite{b12a}\,(\texttt{exponents}); App.~A, Result~\xref{res:exponents}{6} \\
$\gamma$ & Cutoff $w_c$-scaling exponent & $\gamma=\partial\ln\dcrit/\partial\ln w_c=1$ (empirical plateau fit $\approx1.02$) & \cite{b12a}\,(\texttt{exponents}); App.~A, Result~\xref{res:exponents}{6} \\
$\mathrm{Re}$ & Reynolds number & $\mathrm{Re}=\rho U L_{\mathrm{char}}/\mu$ & \cite{b13}, which defines ``$\mathrm{Re}$ is the Reynolds number'' \\
$\varepsilon$ & \emph{Source notation, not a dissertation variable.} In the deterministic-lateral-displacement literature $\varepsilon$ is the row-shift fraction, and the cross-family comparison of Section~\xref{app:units_crosslineage}{``Cross-Lineage Validation''} quotes prefactors ``across $\varepsilon=0.01$--$0.5$'' in that sense. It is never used as a CIF quantity here & quoted only & \cite{b17,b22} \\
$\Pi$ & Cumulative product (in $F_{\mathrm{recursion}}$) & $\Pi = \prod_i (1 - 2 \fgap(i))$ & Derived from Eqn~3 of \cite[p.~4498]{b1} \\
\end{longtable}\endgroup

\textbf{Convention.} $\fgap$ is the \emph{per-side, per-gap} extraction
fraction. The total bilateral extraction at gap row $i$ is $2\fgap(i)$, which
is the origin of the factor of $2$ in the continuity equation (Eqn~3 of
\cite[p.~4498]{b1}) and in the closed-form identity numerator $2R_c$.

\textbf{Aspect-ratio symbol.} This appendix writes the channel aspect ratio as
$\alpha=h/w_{\mathrm{ch}}$ and reserves $\beta$ for the cutoff-law scaling
exponent. This is the canonical dissertation convention and supersedes the
alternatives. Note that the reference mathematics takes the aspect ratio as a
parameter literally named \texttt{beta} (\texttt{c2(beta)},
\texttt{K\_depth(beta)}, evaluated at $0.80/1.40/2.69$). When reading code,
\texttt{beta} is this appendix's $\alpha$; when reading this appendix, $\beta$
is never an aspect ratio.

\textbf{$\sigma$-symbol collision.} $\sigma_F$ here is a \emph{standard
deviation}, propagated from measurement. It is not the side-flow ratio
$\sigma(i)$, $\sigma_0$ of Table~\xref{tab:units_fluid}{A.2}, which is a march
variable. The rule: a subscripted $\sigma$ with a measured quantity in the
subscript ($\sigma_F$, and $\sigma_{F_{V1}}$, $\sigma_{F_{V3}}$,
$\sigma_{\mathrm{ratio}}$ in the validation subsection) is always a dispersion;
$\sigma$ carrying a row index or the subscript $0$ is always the march
variable.

\textbf{Exponent vs.\ lamina-width collision.} $\beta$ and $\gamma$ are the
cutoff-law \emph{scaling exponents} (log--log slopes; symbolic $\tfrac12$ and
$1$; empirical plateau fits $\approx0.51$ and $\approx1.02$). These must not be
confused with the Inglis first-streamline \emph{lamina width}, which this
appendix writes as $b$ (Inglis writes $D_c=2\beta$; rendered $D_c=2b$ here,
Table~\xref{tab:units_geom}{A.1}). The prototype implementation overloads the
name \texttt{beta} for both (\texttt{scaling.beta} for the exponent,
\texttt{Dc=2*beta} for the width); this appendix disambiguates them as $\beta$
(exponent) and $b$ (width). The exponents are the falsifiable functional-form
test of the cutoff law: a single global power law over the full range is
rejected at reduced $\chi^2\approx30.6$, while the plateau exponents lie within
the confidence intervals of $\beta=\tfrac12$, $\gamma=1$, surviving
digitization noise.

%==============================================================
% A.4 -- CUTOFF-LAW PREFACTOR DECOMPOSITION  (complete, corrected)
%
% Requires, in the preamble:
%   \input{tex/kernel_pairing_values}   % generated by gen_kernel_pairing.py
% and \usepackage{longtable}.
%==============================================================
\subsection{Cutoff-Law Prefactor Decomposition (selectivity axis)}
\label{app:units_cutoff}

The selectivity axis is governed by the cutoff law
\begin{equation}
\dcrit = \kapp \, w_c \, \sqrt{\fgap/3},
\qquad
\skimw = w_c\sqrt{\fgap/3},
\label{eqn:cutoff_law}
\end{equation}
in which $\skimw$ is the leading term of the exact dividing-streamline
inversion of the skim cubic $\fgap = 3x^2 - 2x^3$ (Inglis $\varepsilon$-cubic,
\cite[p.~656]{b17}; $x=\skimw/w_c$), and $\kapp$ is the \emph{apparent}
prefactor that maps the idealized skim width onto the measured cutoff
diameter. \textbf{$\kapp$ is a plateau-region constant, not a universal
one.} It is constant only for $\ud\ge0.035$; below that the prefactor
becomes a function of the cutoff itself and Eqn~\ref{eqn:cutoff_law}
becomes implicit. Every use of the numerical value $\kappVal$ in this
appendix is a plateau-region use. The prefactor factorizes as
\begin{equation}
\kapp = \kster \cdot \Kin \cdot \Kdepth \cdot \sqrt{c}\cdot C
       = 2 \cdot 1 \cdot \Ral \cdot \sqrt{c}\cdot C .
\label{eqn:kappa_decomp}
\end{equation}

% longtable, NOT a float: eleven paragraph-length rows do not fit one page,
% and an oversized float is reported as "Float too large" and dropped. No bare
% l or c column sits beside an X column. Fractions sum to 0.92, budgeting the
% inner \tabcolsep gutters at 4pt.
\begingroup\footnotesize\setlength{\tabcolsep}{4pt}
\begin{longtable}{@{}>{\raggedright\arraybackslash}p{0.085\textwidth}
>{\raggedright\arraybackslash}p{0.185\textwidth}
>{\raggedright\arraybackslash}p{0.375\textwidth}
>{\raggedright\arraybackslash}p{0.095\textwidth}
>{\raggedright\arraybackslash}p{0.16\textwidth}@{}}
\caption{Cutoff-law prefactor factors. All strictly dimensionless. $\kster$,
$\Kin$ and the depth factor $\Kdepth$ are determined from first principles or
single-gap CFD; the product
$\sqrt{c}\,C$ is the single remaining calibrated (open)
quantity.}\label{tab:units_cutoff}\\
\toprule Symbol & Name & Value / defining relation & Status & Source \\\midrule\endfirsthead
\multicolumn{5}{@{}l}{\footnotesize\itshape Table \thetable\ (continued)}\\
\toprule Symbol & Name & Value / defining relation & Status & Source \\\midrule\endhead
\bottomrule\endlastfoot
$\kapp$ ($\equiv\kg$) & Apparent cutoff prefactor, cubic branch. \textsc{establishes} & $\dcrit = \kapp w_c\sqrt{\fgap/3}$; \emph{supplied} $\kapp = \kappVal$, 95\% CI $[\kappCIlo,\kappCIhi]$, $\mathrm{CV}=\kappCV\%$, plateau $\ud\gtrsim0.035$. \emph{Refit} from the digitised points gives $\kappFitFull$ ($=\kappFit$ to 4~d.p.); the two agree to four significant figures and differ by $\kappSuppliedGap$ relative. Use $\kappVal$ in prose and equations; use $\kappFit$ only where the fit itself is the subject & \textsc{fitted} & \cite[p.~4500, Fig.~3]{b1} digitization; \cite{b12a} \\
$\kF$ & Prefactor paired with the finite-aspect-ratio lamina kernel $F_{\mathrm{lam}}$ & $\kF=\kg\sqrt{\overline{c_2(\acal)}/3}=\kFval$ at the \emph{assumed} $80~\si{\micro\meter}$ calibration depth. \textbf{The kernel $F_{\mathrm{lam}}$ is exact; the scalar $\kF$ is not.} It is an array-level transferred prefactor obtained from the cubic-branch value using the mean $c_2$ over an assumed calibration set, so it inherits an unreported depth. \textbf{Not interchangeable with $\kapp$}; see the pairing rule below and the $\acal$ entry & \textsc{derived} from an \textsc{external} input & This work, computed from \cite{b12a} (\texttt{c2}, \texttt{plateau\_kappa}) \\
$\kster$ & Steric factor & $\boldsymbol{=\kstericVal}$: cutoff radius $=$ lamina width $\Rightarrow$ cutoff diameter $=2\times$lamina ($D_c=2b$; $b$ is defined in Table~\xref{tab:units_geom}{A.1}) & \textsc{derived}, corroborated & \cite[p.~656]{b17}; \cite[p.~917]{b5} main text; \cite{b12a} (\texttt{KAPPA\_STERIC}) \\
$\Kin$ & In-plane funnelling factor & $=\KinVal$ (single-gap CFD Runs A/B/C; no in-plane aperture funnelling) & \textsc{numerical} & \cite{b7}; \cite{b12a} (\texttt{K\_INPLANE}) \\
$\Kdth(\alpha)$ & Depth factor, \emph{analytic} & $\Kdth=\sqrt{3/c_2(\alpha)}=\lim_{x\to0}\skimw_{3D}/\skimw_{2D}$ from the rectangular-duct law. $\KdepthAnalyticShallow/\KdepthAnalyticMid/\KdepthAnalyticDeep$ at $\alpha=0.80/1.40/2.69$ & \textsc{derived} & \cite[p.~875]{b6}; \cite{b12a} (\texttt{c2}, \texttt{K\_depth}) \\
$\Kdcfd(\alpha)$ & Depth factor, \emph{CFD} & Single-gap CFD: $\KdepthCfdShallow/\KdepthCfdMid/\KdepthCfdDeep$ at the same three $\alpha$. Agreement $|\Kdcfd-\Kdth|/\Kdcfd\le0.6\%$ (actual $\KdepthDevShallow/\KdepthDevMid/\KdepthDevDeep\%$, i.e.\ at most $\KdepthDevMax\%$, the largest rounded up to a bound, at $\alpha=1.40$; the denominator is the CFD anchor, matching the gated \texttt{cfd\_kdepth\_obs}) & \textsc{numerical} & \cite{b7}; \cite{b12a} \\
$\acal$ & Calibration aspect-ratio array & $(0.8000,\,0.6400,\,0.5333)$ at $w_c=100/125/150~\si{\micro\meter}$ \emph{assuming} $h_{\mathrm{cal}}=80~\si{\micro\meter}$. The depth is \textbf{unreported}: a genuine free parameter, and an \textsc{external} input rather than a computed quantity. Two distinct sensitivities follow and must be named apart; see below & \textsc{external}, undischargeable & This work (assumed); the assumption is recorded in \cite{b12a} \\
$c$ & Per-device over-extraction factor. \textsc{competes} with the finite-particle account for the same residual & Enters as $\sqrt{c}$; the cumulative per-gap over-extraction tied to gap count $N$. \emph{Lower-case $c$ under the radical is not the finite-particle factor $C$ below}: see the collision note & \textsc{proposed} & \cite{b12a} \\
$C$ & Cumulative finite-particle factor & Body-averaged finite-sphere displacement vs.\ point streamline; \emph{open}, to be computed by a finite-body method of the MOPSA class (\cite[p.~2]{b19}). See the three-layer note below: the \emph{mechanism} is established, the \emph{method class} is available, the \emph{factor} is this work's parameterization & \textsc{proposed}, uncomputed & \cite[p.~359]{b21} (mechanism); \cite[p.~2]{b19} (method class) \\
$\sqrt{c}\,C$ & Open residual (two factors, three gaps). \textsc{establishes} the cutoff magnitude; its \emph{split} does not & $\Kdepth\sqrt{c}\,C=\kapp/\kster=\residualFolded$ (depth-folded, $\alpha$-independent); $\sqrt{c}\,C=\kapp/(\kster\Kdcfd)=\residualSepShallow/\residualSepMid/\residualSepDeep$ over $\alpha=0.80/1.40/2.69$, depth-separated. \emph{Never print the two conventions as one range} & \textsc{fitted} / \textsc{proposed} & \cite{b12a} \\
\end{longtable}\endgroup

\textbf{$c$/$C$ collision.} Lower-case $c$ is the per-device over-extraction
factor and upper-case $C$ the cumulative finite-particle factor. They are
different quantities that differ only by case, and they appear together in a
single product, $\sqrt{c}\,C$, which is the most misreadable arrangement
possible. The radical is the only typographic cue: $c$ is \emph{always} under
it, $C$ \emph{never}. When either appears alone, restate which it is.

\textbf{One open residual, two factors, three gaps.} The residual is
\emph{derived from a fitted prefactor}, not measured: $\kapp$ is calibrated
on the Gifford plateau~\cite{b1,b12a}, and published single-lineage data fix
the product $\sqrt{c}\,C$ but not its split into $c$ and $C$~\cite{b12a}.
Within the model's premises the transport lemma makes $\kster=\kstericVal$
exact, so the residual $\sqrt{c}\,C$ must arise where a premise fails; the three
model gaps the decomposition tracks are the flat-lamina idealisation, the wall
porosity $\phi=G/\Lambda$, and finite-particle displacement~\cite{b12a}.
Yet the factors number two: $C$ is the finite-particle displacement alone
(the anti-Zweifach--Fung attraction to the low-flow branch~\cite{b21}), and
because the decomposition $\kapp=\kster\Kin\Kdepth\,\sqrt{c}\,C$ closes by
construction ($\kster=\kstericVal$, $\Kin=\KinVal$), $\sqrt{c}\,C$ is the
entire residual~\cite{b12a}, so $\sqrt{c}$, the per-gap over-extraction, is
that residual divided by $C$ and carries the flat-lamina and porosity
contributions together without isolating either. This grouping follows from
the closure, not from any statement in the record, which names porosity as a
source but assigns it to no factor~\cite{b12a}; neither the single-gap solve
nor the yield identity resolves the split, which is what leaves the residual
open~\cite{b12a}.

\textbf{Calibration depth (caveat).} $\kapp = \kappVal$ is fit on the
Gifford Fig.~3 array, whose caption reports only the three center-channel
widths ($w_c = 100,125,150~\si{\micro\meter}$) and \emph{no depth}. The
value $\alpha = 0.8$ used for $\Kdepth$ borrows the $80~\si{\micro\meter}$
depth of the \emph{separate} Fig.~4 demonstration device at
$w_c = 100~\si{\micro\meter}$; it is therefore an assumption, and
$\Kdepth$ in fact varies across the three calibration widths
($\alpha = 0.8000/0.6400/0.5333$ at $w_c = 100/125/150~\si{\micro\meter}$,
giving $\Kdth = \KdepthCalNarrow/\KdepthCalMid/\KdepthCalWide$). The depth
correction is therefore $11$--$19\%$ across the calibration channels, not
negligible. Crucially, $\Kdepth$ is folded into the \emph{apparent}
prefactor as an interpretive factor and is \emph{not} applied to the
calibration data (which collapse more tightly on the parabolic $\skimw$),
so it does not shift the fitted $\kapp = \kappVal$; it enters only the
residual decomposition (depth-folded $\kapp/\kster = \residualFolded$
vs.\ depth-separated $\sqrt{c}\,C$). Gifford states three device depths, none of them for this array:
$80\pm3~\si{\micro\meter}$ for the Fig.~4 device (\cite[p.~4500, Fig.~4
caption]{b1}), $125~\si{\micro\meter}$ for the four Fig.~5 devices
(\cite[p.~4501, Fig.~5 caption]{b1}) and $150\pm5~\si{\micro\meter}$ for
the Fig.~6 platelet-enrichment device (\cite[p.~4502, Fig.~6 caption]{b1}),
so the borrowed $80~\si{\micro\meter}$ is the shallowest depth the paper
reports. \emph{The true array depth is not
known and is not inferred here}; the consequences of that ignorance are
quantified in the two $\acal$ sensitivities below.

\textbf{The two $\acal$ sensitivities, named apart.} Both are correct and
they answer different questions. Do not call either ``the
$\beta_{\mathrm{cal}}$ swing'' without saying which.
\begin{itemize}
\item $\Delta N_{\mathrm{kernel|depth}} = +26.9\%$: the row-count
inflation of an $F_{\mathrm{lam}}$ design over a cubic design at
$h_{\mathrm{cal}}=80~\si{\micro\meter}$ (against the cubic baseline
$N=2496$), falling to $+3.9\%$ at $150~\si{\micro\meter}$.
\textbf{Conditional on adopting $F_{\mathrm{lam}}$}: if the kernel is
never switched, $h_{\mathrm{cal}}$ does not enter $N$ at all, because
$\kg$ is fitted self-consistently on the cubic. This is the single
dominant risk the design code flags: the one number that
swings $N$ by roughly $27\%$. \emph{Provenance note: unlike every other
figure in this subsection, these three numbers come from a calibration-depth
sweep that is not reproduced in the derivation record and are not gated.}
\item $\Delta_{\mathrm{transfer}} = \transferDecl\%$: the within-array
transfer bias, $c_2(0.53)/c_2(0.80)-1$, for the array \emph{as declared}
$(0.80,0.64,0.53)$; $\transferExact\%$ using the exact ratio $0.5333$. The
exact value is canonical in this appendix (a two-decimal input should not
drive a four-decimal output); the declared-array figure is the one the
verification registry gates and is labelled as such wherever it appears.
\end{itemize}

\textbf{Steric criterion (adjudicated).} The cutoff radius equals the
extracted-lamina width (Dinh main text, $a<w$; \cite[p.~917]{b5}), and the
cutoff \emph{diameter} is therefore twice the lamina width, matching
Inglis's $D_c = 2\beta$ (\cite[p.~656]{b17}; $D_c=2b$ in this appendix's
notation) and the Inglis--Herman restriction
of particles ``geometrically restricted away from walls by a distance equal
to their radius'' (\cite[p.~1724]{b18}). Hence
$\boldsymbol{\kster = \kstericVal}$. The Dinh Fig.~1 caption phrasing
``about half the width of the extracted flow lamina'' (equivalent to
$\kster=1$) is \emph{deselected} in favor of the main text and Inglis.
\emph{This is an adjudication between two readings of the same paper, not a
preference}: Dinh's main text and Fig.~1 caption are formally inconsistent,
so one must be retracted, and the main text is selected because it agrees
with the two independent statements of the same criterion in \cite{b17} and
\cite{b18}.

\textbf{Finite-particle factor $C$: three layers, kept apart.} The
residual has physically motivated candidate explanations; its
finite-particle component has not yet been computed directly. The evidence
comes at three different strengths.
\begin{enumerate}
\item \emph{Established physical mechanism.} A finite sphere near a
bifurcation is drawn toward the \emph{low-flow-rate} branch, the
anti-Zweifach--Fung effect (Doyeux \emph{et al.}, \cite[p.~359]{b21}: the
attraction is ``directed towards the low flow rate branch''), which in CIF
is the side channel. Finite particles are therefore over-extracted
relative to a point-streamline model, so the streamline that just excludes
a particle is narrower than its radius and the measured cutoff diameter
exceeds the idealized steric value $2\skimw$. This is the physical content
of $C>1$.
\item \emph{Available computational methodology.} Finite-body trajectory
models of the MOPSA class (\cite[p.~2]{b19}) can represent particles as
finite rigid objects rather than point tracers; the paper ``treats the
particle as a two-dimensional rigid circular object instead of a single
point.''
\item \emph{The factor $C$ itself.} A dissertation-specific
parameterization of the remaining finite-particle contribution, deferred
to finite-particle work. It is \emph{not} independently established by
either \cite{b19} or \cite{b21}; those support the mechanism and the
method, not this multiplicative form.
\end{enumerate}
Its sign, mechanism and scale are corroborated on an unrelated device
family in Section~\xref{app:units_crosslineage}{``Cross-Lineage Validation''}.

\textbf{Foot region and retraction.} For $\ud<0.012$ the collapse degrades
to a \emph{sub-steric, recovery-suppressed} foot ($\kapp\to\footKappa$,
below the steric floor $\kster\Kdcfd\approx\stericFloorShallow$).
\emph{Both $\plateauKappa$ and $\footKappa$ are binned means, not
extremes}: $\plateauKappa$ is the mean of the $n=\plateauN$ plateau bin
($\ud\ge0.035$) and $\footKappa$ the mean of the $n=\footN$ foot bin
($\ud<0.012$). The fifteen individual points span $\kappaSpanLo$ to
$\kappaSpanHi$; that span is a different statistic and must not be quoted
as the ramp. Depth cannot account for the foot because
$\Kdepth\in[0.81,0.98]$ cannot move $\kapp$ from $\plateauKappa$ to
$\footKappa$.

\emph{Scope of the evidence, stated explicitly.} The foot bin is composed
of three points at $d=1.01~\si{\micro\meter}$ (Gifford Fig.~3 rigid beads,
$w_c=100/125/150~\si{\micro\meter}$), roughly $1~\si{\micro\meter}$
particles, well below the $2.8$--$4.8~\si{\micro\meter}$ platelet cutoffs
this framework targets, so the foot lies outside the regime the
dissertation's claims cover. The corroborating observation comes from a
\emph{different} device and a \emph{different} particle class: Dinh report
that ``recovery of the smallest 2 fL PLTs ($a \approx 1~\si{\micro\meter}$)
in the filtrate was consistently lower than the percent of extracted fluid
volume'' (\cite[p.~917]{b5}), and attribute it to ``crowding of the flow
laminae near the filtration gaps by retained'' platelets, which ``impeded
the free outflow of the smallest PLTs'' (\cite[p.~917]{b5}). That is a
concentration-dependent crowding effect in a platelet suspension, so its
bearing on the dilute rigid-bead foot is by analogy of sign only. Recast in
the language of this decomposition it corresponds to an effective sampled
lamina narrower than nominal
($\skimw_{\mathrm{eff}}<\skimw_{\mathrm{nominal}}$), so that
$\kapp=\kster(\skimw_{\mathrm{eff}}/\skimw_{\mathrm{nominal}})$ falls below
steric; the recasting is this work's inference, the crowding mechanism is
the source's, and the two are not rival accounts. Platelets are discoid and
outside the rigid-sphere envelope
(Section~\xref{app:units_particle}{``Particle and Separation Variables''}),
so the Dinh observation is offered as a mechanistic analogy of the same
sign, not as evidence within scope, and no cutoff number is drawn from it.
\textbf{The knee and foot mechanisms remain open}: the present data
establish that the prefactor is size-dependent, but not which physical
account of that dependence is correct.

\textbf{The $f_{\mathrm{eff}}$ retraction is narrow.} A
\emph{size}-dependent effective flow fraction $f_{\mathrm{eff}}(d)$ remains
a competing account of the knee, stated in its strongest form and not
resolved here: the published data do not discriminate between it and the
finite-particle account, because both predict the same sign and the same
qualitative size dependence. It is \emph{not} retracted. The cutoff law and
the yield law are independent of which account is correct, since neither
account enters them; only the \emph{interpretation} of
the residual turns on it. Retracted: $f_{\mathrm{eff}}(N)$. Open and
unresolved: $f_{\mathrm{eff}}(d)$.

\textbf{Scaling equivalence.} Eqn~\ref{eqn:cutoff_law} is equivalent to
$D_c \propto w_c\sqrt{\fgap}$ with proportionality constant
$\kapp/\sqrt{3} \approx \kappOverRootThree$; the $1/3$ is carried
explicitly inside the radical so that $\kapp = \kappVal$ is unambiguous.

% =====================================================================
% KERNEL-PREFACTOR PAIRING.
% Every COMPUTED number in the block below is a macro from
% tex/kernel_pairing_values, generated by gen_kernel_pairing.py from the
% derivation record and regenerated on every build. Do not type a computed
% value here; regenerate. Three kinds of literal remain, all deliberate:
%   * INPUTS -- the h_cal/w_c ratios 0.8000/0.6400/0.5333 (and their declared
%     two-decimal forms), and the Dinh geometry h = 140 um, f_gap = 1.04e-4.
%     An input is not generated; it is assumed, and is labelled as such.
%   * the flagged EXCEPTION, below.
% An earlier revision claimed "every number is a macro" while three computed
% leading-order values were still typed. Counting exceptions is fragile; the
% rule is now stated by KIND.
% The ONLY exception is flagged in the text where it appears: +26.9%, +3.9%
% and N = 2496 come from a calibration-depth sweep the record does not
% reproduce. Every other number here, including the leading-order values
% quoted to show what the approximation costs, is a macro.
% (Before this revision every one of the 31 numbers in this block was a typed
% literal sitting under a comment that said they were generated.)
% =====================================================================

\textbf{Kernel--prefactor pairing rule.}
$\kapp$ is meaningless without the kernel it was fitted with. Two kernels
are in use: the leading-order skim branch $g:\;\skimw=w_c\sqrt{\fgap/3}$
with prefactor $\kg=\kappVal$; and the exact lamina-flux kernel
$F_{\mathrm{lam}}$ with prefactor $\kF$. Declare a pair and use \emph{that
kernel at all three call sites}: the inversion, the schedule, and the scale
$\dcrit=\kappa_{\mathcal K}\skimw$.

\textbf{Conversion.} Using the \emph{array-mean} $c_2$ over the calibration
widths:
\begin{equation}
\kF=\kg\sqrt{\overline{c_2}/3}=\kg\big/\sqrt{3/\overline{c_2}},\qquad\overline{c_2}=\tfrac13\textstyle\sum_i c_2(\alpha_{\mathrm{cal},i}),
\label{eqn:kernel_pairing}
\end{equation}
with $\acal=h_{\mathrm{cal}}/w_c=0.8000/0.6400/0.5333$ at the \emph{assumed}
$h_{\mathrm{cal}}=80~\si{\micro\meter}$, giving
$\overline{c_2}=\ctwobar$ and $\kF=\kFval$.

\textbf{Two sources of numerical discrepancy in that formula.} Neither is
pedantry; both move the third significant figure.
\begin{enumerate}
\item \emph{Do not average the depth factors.}
$\sqrt{3/\overline{c_2}}=\depthEff$ is an \emph{effective} depth factor
derived from the mean $c_2$. The arithmetic mean of the three $\Kdepth$
values gives $\kF=\kFmeanK$ instead, $\kFmeanKerr\%$ off, because the mean
of the square roots is not the square root of the mean.
\item \emph{Do not evaluate at one aspect ratio.} Using $\Kdth(0.80)$
alone gives $\kFoneAlpha$, $\kFoneAlphaErr\%$ low. The conversion is an
array quantity.
\end{enumerate}

\textbf{What crossing costs, computed.} At the Dinh geometry
($h=140~\si{\micro\meter}$, $\fgap=1.04\times10^{-4}$). Throughout this
table $\skimw_g$ and $\skimw_F$ are the \emph{exact} inversions of the
cubic and of $F_{\mathrm{lam}}$ respectively, so that the two branches are
compared like for like, and the percentages are computed on those exact
widths at full precision, not from the rounded cells. The leading term
$w_c\sqrt{\fgap/3}$ of $\skimw_g$ is not used here: substituting it would
give $\dcritGGleadIn/\dcritGGleadOut$ for the first column, $\choiceLeadGap\%$
below the tabulated $\dcritGGIn/\dcritGGOut$, and would raise the
kernel-choice figures above the tabulated $+\choiceIn\%/+\choiceOut\%$.

\begin{center}\footnotesize
\begin{tabular}{@{}llrrrr@{}}
\toprule
station & $\alpha=h/w_c$ & $\kg\skimw_g$ & $\kF\skimw_{F}$ & $\kF\skimw_g$ & $\kg\skimw_{F}$ \\
\midrule
inlet ($w_c=75$) & $\alphaDIn$ & $\dcritGGIn$ & $\dcritFFIn$ & $\dcritFgIn$ & $\dcritgFIn$ \\
outlet ($w_c=52$) & $\alphaDOut$ & $\dcritGGOut$ & $\dcritFFOut$ & $\dcritFgOut$ & $\dcritgFOut$ \\
\bottomrule
\end{tabular}
\end{center}

The first two columns are the \emph{self-consistent} pairs; the last two
are crossings. Reading the inlet row: the kernel \emph{choice} alone
already moves the answer by $+\choiceIn\%$ ($\dcritGGIn$ vs
$\dcritFFIn~\si{\micro\meter}$), rising to $+\choiceOut\%$ at the outlet.
Crossing a prefactor onto the other kernel's lamina shifts the result by a
further $+\crossUp\%$ or $\crossDown\%$; one pair of figures covers both
rows because the lamina width cancels, leaving the pure prefactor ratio
$\kF/\kg-1$ and its inverse. The mismatch is not detected automatically and
is therefore a potentially silent modelling error.

\textbf{Which to use.} The specification is directive, not neutral: until
the calibration depth is measured, \emph{use the cubic branch with $\kg$},
which is self-consistent with the fit, and report the transfer-bias band
as an uncertainty rather than applying a correction whose calibration
cannot be verified. This dissertation's cutoff law
(Eqn~\ref{eqn:cutoff_law}) and its entire validation path are on the cubic
branch. The design code carries both pairs: its \texttt{design} function
defaults to the cubic pair, and the physics layer that the design-map test
marches (App.~B, P-C8.2) runs on $F_{\mathrm{lam}}$. The
pairing rule is what keeps the two apart. See the $\acal$ entry for the
row-count consequence of switching
($\Delta N_{\mathrm{kernel|depth}}=+26.9\%$).

\textbf{Rounding convention for $\acal$.} The array is declared in the
design code as $(0.80,\,0.64,\,0.53)$, but $0.53$ is a
two-decimal rounding of $h_{\mathrm{cal}}/w_c=0.5333$. This appendix uses
the \emph{exact} ratio: $\overline{c_2}=\ctwobar$ and $\kF=\kFval$. Under
the declared rounding they are $\ctwobarDecl$ and $\kFdecl$. The same split
appears in the transfer bias: $\transferExact\%$ exact against
$\transferDecl\%$ for the declared array, which is the figure the
verification registry gates. Both are correct for what they are; the exact
ratio is canonical here and the declared-array figure is labelled as such
wherever it appears.

%==============================================================
% A.5 -- PARTICLE AND SEPARATION VARIABLES  (complete, corrected)
%==============================================================
\subsection{Particle and Separation Variables}
\label{app:units_particle}

Table~\ref{tab:units_particle} gives the particle variables. The
distinction that matters is between a volume-equivalent sphere and a real
discoid platelet: the cutoff law returns the former.

% longtable for consistency with A.1-A.4. No bare l beside an X column, and the
% fractions sum to 0.92 to budget the inner \tabcolsep gutters.
\begingroup\small\setlength{\tabcolsep}{4pt}
\begin{longtable}{@{}>{\raggedright\arraybackslash}p{0.095\textwidth}
>{\raggedright\arraybackslash}p{0.46\textwidth}
>{\raggedright\arraybackslash}p{0.075\textwidth}
>{\raggedright\arraybackslash}p{0.115\textwidth}
>{\raggedright\arraybackslash}p{0.175\textwidth}@{}}
\caption{Particle-property variables, together with the one
separation-performance quantity they set: the cutoff diameter. The filtrate
fraction $F$ and its variants are in
Table~\xref{tab:units_dimensionless}{A.3}, and the validation dispersions in
the V\&V table.}\label{tab:units_particle}\\
\toprule Symbol & Quantity & SI unit & Practical unit & Source \\\midrule\endfirsthead
\multicolumn{5}{@{}l}{\footnotesize\itshape Table \thetable\ (continued)}\\
\toprule Symbol & Quantity & SI unit & Practical unit & Source \\\midrule\endhead
\bottomrule\endlastfoot
$\dcrit$ & \emph{Critical (cutoff) diameter} (canonical; output of Eqn~\xref{eqn:cutoff_law}{1}) & \si{\meter} & \si{\micro\meter} & \cite{b12a}; \cite[p.~656]{b17} \\
$D_c$, c.d. & Cutoff diameter (Xia/lineage symbol, $\equiv \dcrit$); Xia introduces the abbreviation as ``particles above a certain size (or `critical diameter', c.d.)'' & \si{\meter} & \si{\micro\meter} & \cite[p.~2 (Sci Rep PDF pagination)]{b2} \\
$a$ & Particle radius, as used in the \emph{adopted} capture criterion $a<w$ (main text). \emph{Not} the ``effective radius'' of the deselected Fig.~1 caption on p.~916; see the steric adjudication in Section~\xref{app:units_cutoff}{``Cutoff-Law Prefactor Decomposition''} & \si{\meter} & \si{\micro\meter} & \cite[p.~917]{b5} \\
$D=2a$ & Particle diameter (sphere) & \si{\meter} & \si{\micro\meter} & Standard \\
$V_{\mathrm{sph}}$ & \emph{Sphere} volume, $V_{\mathrm{sph}}=\tfrac43\pi a^3$. Used for the equivalent-sphere conversion only; real platelets are discoid and this is \emph{not} their volume formula & \si{\cubic\meter} & fL ($\equiv \si{\micro\meter\cubed}$) & Standard \\
$V_{12}$ & Dinh platelet-volume \emph{classification threshold}, $12$~fL: ``cell volumes larger than 12 fL are commonly considered `large''', the P-LCR cut used to define ``$>12$~fL recovery''. \emph{Not} a minimum retained volume & \si{\cubic\meter} & fL & \cite[p.~913]{b5} \\
\end{longtable}\endgroup

\textbf{$D$/$d$ collision: diameter-symbol reconciliation.} $\dcrit$ is canonical for the
cutoff law; $D_c$ is its lineage synonym (Xia). Neither is to be confused
with the deprecated depth symbol $d$ (use $h$), nor with the hydraulic
diameter $D_h$ (Table~\xref{tab:units_geom}{A.1}). For a sphere $D=2a$.

\textbf{Two further case collisions, declared.} $V$ is a volume; lower-case
$v$ is a local fluid velocity (Table~\xref{tab:units_fluid}{A.2}).
$V_{\mathrm{sph}}$ and $V_{12}$ are one quantity, a volume, read two ways.
A particle radius is written $a$ throughout and \emph{never} $r$. Lower-case
$r$ is not used as a variable anywhere in this appendix, and reserving it
keeps it clear of the resistances $R_c$, $R_s$, $R_{\mathrm{gap}}$ and their
aliases (Table~\xref{tab:units_fluid}{A.2}).

\textbf{Volume convention.} $1~\text{fL} = 1~\si{\micro\meter\cubed}$
(by definition: $1$~femtoliter $= 10^{-15}$~liter
$= 10^{-18}~\si{\meter\cubed} = (10^{-6})^3~\si{\meter\cubed}$).
For comparison with the rigid-sphere model, $12$~fL corresponds to an
equivalent-sphere diameter
$D = 2\,(3V_{\mathrm{sph}}/(4\pi))^{1/3} \approx 2.84~\si{\micro\meter}$.
\emph{This is a geometric conversion, not a measured platelet diameter.}
Dinh report the actual \emph{discoid} platelet as
$2a \approx 4.8~\si{\micro\meter}$ at $12$~fL, which in their own terms is the
major/equatorial axis of an oblate spheroid, i.e.\ the in-plane one
(\cite[p.~917]{b5}); the
$\sim\!1.7\times$ gap between the two reflects the sphere idealization and
therefore marks a shape-related \emph{limitation} of applying the
rigid-sphere model to platelets. The model is in scope; real platelets are
not validated against it. A cutoff diameter
predicted by Eqn~\xref{eqn:cutoff_law}{1} is therefore an
\emph{equivalent-sphere} diameter; for discoidal platelets it under-states
the physical in-plane diameter by this factor.


%==============================================================
% A.6 -- VALIDATION AND UNCERTAINTY QUANTITIES  (complete, corrected)
%
% Requires, in the preamble:
%   \input{tex/validation_values}   % generated by gen_validation_values.py
%==============================================================
\subsection{Validation and Uncertainty Quantities (V\&V~20, yield axis)}
\label{app:units_validation}

Table~\ref{tab:units_validation} lists the validation and uncertainty
quantities with the COMSOL and reproducibility diagnostics. The framework
is ASME V\&V~20~\cite{b25}: the comparison error $E=S-D$ is judged against
a combined validation uncertainty $\uval=\sqrt{\unum^2+\uinput^2+\uD^2}$,
of which the experimental term $\uD=\FsdVone$--$\FsdVthree$ is by far the
largest. Each Dinh variant was solved at two mesh levels a ${\sim}5$-fold
refinement apart (V1: $115{,}904$ and $24{,}783$ elements; V3: $104{,}206$
and $21{,}448$), fully converged in four Newton iterations,
\emph{these are the convergence-study meshes}; the production V3 solve used a
slightly different mesh of $104{,}651$ elements, so the two V3 figures are
different meshes rather than a discrepancy, and the
reported $F$ is identical at both~\cite[COMSOL reports; Block~19]{b12}. \emph{That
bounds $\unum$ at $5\times10^{-5}$ rather than measuring it}: five
significant figures is the precision to which the flow tables are printed,
so the true discretisation error is smaller and unresolved. Even as a bound it is at
least $\uDoverUnum\times$ below $\uD$ ($\uDoverUnumVOne\times$ at V1,
$\uDoverUnumVThree\times$ at V3).

The remaining term, $\uinput$, is bounded by fabrication tolerance in two
independent components; the boundary conditions do not contribute
separately, because $F$ is a ratio of filtrate flow to inlet flow, in
which the imposed inlet rate and the outlet pressure datum both cancel
identically. \emph{Out of plane}, the channel-depth tolerance shifts $F$
by at most $0.031\,\uD$ at V1 and $0.015\,\uD$ at V3 over
$\pm5~\si{\micro\meter}$~\cite[\S8.6]{b12}. \emph{In plane}, $F$ is
likewise a channel-resistance \emph{ratio} in which isotropic scaling of
the widths and lengths cancels to leading order, the only residual being
the aspect-ratio bracket of the resistance law
(Eqn~\xref{eqn:gifford5}{5}), which is second order in the depth-fixed
width perturbation and stays well below the out-of-plane term, so the
$\uvalExcessMax\%$ bound below already accommodates it. So $\uval$ exceeds $\uD$ by at most
$\uvalExcessMax\%$ ($\uvalExcessVOne\%$ at V1, $\uvalExcessVThree\%$ at
V3): to the two decimals given $E/\uval$ equals $E/\uD$, so the reported
distances are governed by the choice of $\uD$ scale, not by the small
$\unum$ or $\uinput$ terms. See the SD-vs-SEM note for that choice.

% longtable for consistency with A.1-A.5, and because the fixed p{} columns are
% what prevent an over-wide cell. Fractions sum to 0.92 to budget the gutters.
\begingroup\footnotesize\setlength{\tabcolsep}{4pt}
\begin{longtable}{@{}>{\raggedright\arraybackslash}p{0.175\textwidth}
>{\raggedright\arraybackslash}p{0.47\textwidth}
>{\raggedright\arraybackslash}p{0.085\textwidth}
>{\raggedright\arraybackslash}p{0.19\textwidth}@{}}
\caption{Quantities used in validation, COMSOL simulation, and
reproducibility infrastructure, including the ASME V\&V~20 uncertainty
terms.}\label{tab:units_validation}\\
\toprule Symbol & Quantity & Unit & Source \\\midrule\endfirsthead
\multicolumn{4}{@{}l}{\footnotesize\itshape Table \thetable\ (continued)}\\
\toprule Symbol & Quantity & Unit & Source \\\midrule\endhead
\bottomrule\endlastfoot
$|\Delta F_{\mathrm{1D}}|$ & Absolute difference: closed-form $F$ vs.\ recursion endpoint. \emph{A one-dimensional self-consistency check, not the mesh-to-mesh difference} & --- & \cite[\S8.4]{b12} \\
$\uval$ & Combined validation uncertainty, $\uval=\sqrt{\unum^2+\uinput^2+\uD^2}$; evaluates to $\approx\uD$ here & --- & \cite{b25}; \cite{b7} \\
$\unum$ & Numerical (discretisation) uncertainty. \emph{Bounded, not measured}: $\le5\times10^{-5}$, the reporting precision at which $F$ is unchanged across two mesh levels a ${\sim}5$-fold refinement apart, i.e.\ at least $\uDoverUnum\times$ below $\uD$ & --- & \cite{b7}; \cite[COMSOL reports]{b12} \\
$\uinput$ & Input uncertainty (geometry, boundary conditions). $\le0.031\,\uD$ from the channel-depth tolerance over $\pm5~\si{\micro\meter}$; isotropic in-plane fabrication error cancels to leading order in the resistance ratio & --- & \cite[\S8.6]{b12} \\
$\uD$ & Experimental (data) uncertainty, reported here as the published $\mathrm{SD}$ (a conservative choice; see the SD-vs-SEM note). V\&V~20's $\uD$ for a mean of $n$ has random component $\mathrm{SEM}=\mathrm{SD}/\sqrt{n}$, reported alongside & --- & \cite[p.~919]{b5} \\
$E$, $S$, $D$ & Validation comparison error $E=S-D$, simulation value $S$ minus data value $D$. \emph{All three letters are reused}: see the collision note below & --- & \cite{b25} \\
$T_{\mathrm{PE,1/2}}$ & PE half-life: gap rows for outlet pressure mismatch to halve & gaps (count) & \cite[\S8.7]{b12} \\
$\alpha_{\mathrm{3D}}$ & 3-D outlet resistance-ratio correction $(R_c/R_s)_{\mathrm{app}}/(R_c/R_s)_{\mathrm{Yang}}$ (yield-side). \emph{Not} the channel aspect ratio $\alpha$ of Table~\xref{tab:units_dimensionless}{A.3} & --- & \cite{b7}; \cite[\S8.8]{b12} \\
$\sigma_{F_{V1}}, \sigma_{F_{V3}}$ & One-SD experimental $F$ uncertainty (Dinh, $n{=}4$, panel B). A subscripted $\sigma$ naming a measured quantity is a dispersion, never the side-flow ratio $\sigma(i)$ of Table~\xref{tab:units_fluid}{A.2} & --- & \cite[p.~919]{b5} \\
$\sigma_{\text{ratio}}$ & Propagated uncertainty on multi-variant ratio & --- & \cite[p.~919]{b5}; \cite[\S15.4]{b12} \\
$N_{\text{elem}}, N_{\text{DOF}}$ & COMSOL mesh element / degrees-of-freedom count & count & \cite{b7}; \cite[Block~19]{b12} \\
$\epsilon_{\text{conv}}$ & Newton convergence residual tolerance. \emph{Not} the DLD row-shift fraction $\varepsilon$ of Table~\xref{tab:units_dimensionless}{A.3}; the two epsilons are one glyph & --- & \cite{b7} \\
$N_{\text{validators}}, N_{\text{stamps}}$ & Validator-suite assertion / regression-stamp count & count & \cite{b12a} \\
\end{longtable}\endgroup

\textbf{$E=S-D$ collision: three reused letters in one formula.} The
V\&V~20 comparison error is written $E=S-D$, and each letter already
denotes something else in this appendix. $S$ here is the \emph{simulation}
value of $F$; $S(i)$ in Table~\xref{tab:units_fluid}{A.2} is the normalised
side-channel flow along the march. $D$ here is the \emph{data} value of
$F$; $D=2a$ in Table~\xref{tab:units_particle}{A.5} is a particle diameter,
and $D_c$, $D_h$ are the cutoff and hydraulic diameters. The rule: inside
$E=S-D$, and nowhere else, $S$ and $D$ are the V\&V~20 simulation and data
values; they are dimensionless and always appear in that formula, or as the
reported ratio $E/\mathrm{SD}$, in which the roman $\mathrm{SD}$ is the
standard-deviation abbreviation and not the italic pair $S,D$. Elsewhere they keep their table meanings.

\textbf{$\alpha_{\mathrm{3D}}$ collision.} $\alpha=h/w_{\mathrm{ch}}$ is
the channel aspect ratio and $\acal$ its three-element calibration array
(Table~\xref{tab:units_dimensionless}{A.3}); both are geometric.
$\alpha_{\mathrm{3D}}$ is neither: it is a \emph{resistance-ratio}
correction on the yield axis, dimensionless but unrelated. The subscript is
the only cue, so restate it wherever the two could be read together.

\textbf{$\varepsilon$/$\epsilon$ collision.} \LaTeX{} offers two epsilons
and they are one glyph to a reader. $\varepsilon$ is source notation only,
the DLD row-shift fraction quoted from the deterministic-lateral-displacement
literature (Table~\xref{tab:units_dimensionless}{A.3}); it is never a CIF
quantity. $\epsilon_{\text{conv}}$ is this work's Newton convergence
tolerance. The rule: a bare epsilon is always the source's, and a
subscripted one is always this work's.

\textbf{$u$-symbol collision.} The letter $u$ carries three unrelated
meanings, two of them inherited and not negotiable. $u(y)$ and $u_{\max}$
are fluid velocities (Table~\xref{tab:units_fluid}{A.2}, standard fluid
mechanics). $\uval$, $\unum$, $\uinput$ and $\uD$ are ASME V\&V~20
uncertainties (this table, the standard's own notation). The third,
$\ud=\dcrit/w_c$, was moved off the letter entirely and is rendered
$\widetilde{d}$ (Table~\xref{tab:units_dimensionless}{A.3}) precisely
because $u_d$ and $u_D$ differ by subscript case alone.

\textbf{$N$-symbol collision.} Bare $N$ is the total gap-row count
(Table~\xref{tab:units_geom}{A.1}) and is load-bearing in the schedule and
cutoff laws. Every other $N$ in this appendix carries a subscript naming
what it counts: $N_{\text{elem}}$, $N_{\text{DOF}}$,
$N_{\text{validators}}$, $N_{\text{stamps}}$, all of which are
infrastructure diagnostics. A bare $N$ is never a count of anything but
gap rows.

\textbf{$\sigma$ collision.} A subscripted $\sigma$ that names a
\emph{measured} quantity is a dispersion, never a flow ratio:
$\sigma_{F_{V1}}$, $\sigma_{F_{V3}}$ and $\sigma_{\text{ratio}}$ are
experimental standard deviations (this table), while $\sigma(i)$ and its
inlet value $\sigma_0$ are the normalised side-channel flow along the march
(Table~\xref{tab:units_fluid}{A.2}). The two are told apart by subscript
\emph{content}, not case: a bare or row-indexed $\sigma$ is always the flow
ratio, a device- or ratio-subscripted $\sigma$ is always a standard
deviation.

\textbf{SD vs.\ SEM.} Dinh Fig.~4 reports \emph{mean $\pm$ standard
deviation} ($n{=}3$ panel~A; $n{=}4$ panel~B; \cite[p.~919]{b5}). The
$F$-validation is the two-device comparison in panel~B, so the datum $D$ is
a mean of $n{=}4$ (the V1 recovery $50.5\pm3.3\%$ is also reported in
panel~A at $n{=}3$; the panel~B quadruplicate is the one used here). ASME
V\&V~20 defines $\uD$ as the standard uncertainty appropriate for $D$ and
states that it ``includes all effects of averaging, includes all random and
systematic uncertainty components'' (\cite[\S1-6]{b25}); for a mean of $n$
the random component is therefore the standard deviation of the mean,
$\mathrm{SEM}=\mathrm{SD}/\sqrt{n}$, so the distance measured in units of
the standard's own random uncertainty is $E/\mathrm{SEM}$.

\textbf{The published $\mathrm{SD}$ is nevertheless retained here as
$\uD$}, as a deliberately \emph{conservative} choice on two grounds: it
over-covers that random component by $\sqrt{n}$, and it stands in for the
systematic experimental uncertainty that V\&V~20's $\uD$ formally includes
but that Dinh do not separately report. Both distances are given below so
the reader can see the choice rather than inherit it. Under the stricter
$\mathrm{SEM}$ scale the closed form and the 3-D solve remain within
$2\sigma$ (the largest, V3 COMSOL, is $\EsemComsolVThree\sigma$); on the
raw-$\mathrm{SD}$ scale they are within $1\sigma$.

% =====================================================================
% VALIDATION DISTANCES. Every number in the display below is a macro from
% tex/validation_values, generated by gen_validation_values.py from the same
% F values the document carries. Do not type a distance here; regenerate.
% Deliberately NOT generated, and not implied to be: the mesh element counts
% (COMSOL outputs), the 0.031/0.015 depth-tolerance sensitivities (a notebook
% sweep this record does not reproduce), and the +/-5 um tolerance and 5e-5
% reporting precision (inputs).
% =====================================================================
\textbf{Worked validation numbers (Dinh V1, V3).}
\begin{align*}
F^{\mathrm{closed}}_{V1} &= \FclosedVone, & F^{\mathrm{closed}}_{V3} &= \FclosedVthree
   &&\text{(closed form, no fit)}\\
F^{\mathrm{COMSOL}}_{V1} &= \FcomsolVone, & F^{\mathrm{COMSOL}}_{V3} &= \FcomsolVthree \\
F^{\mathrm{exp}}_{V1} &= \FexpVone \pm \FsdVone, & F^{\mathrm{exp}}_{V3} &= \FexpVthree \pm \FsdVthree
   &&\text{(SD, } n{=}4)\\
E/\mathrm{SD}\ \text{(1-D closed form)}:\quad V1 &= \EsdClosedVOne, & V3 &= \EsdClosedVThree \\
E/\mathrm{SD}\ \text{(3-D COMSOL)}:\quad V1 &= \EsdComsolVOne, & V3 &= \EsdComsolVThree \\
E/\mathrm{SEM}\ \text{(3-D COMSOL)}:\quad V1 &= \EsemComsolVOne, & V3 &= \EsemComsolVThree
   &&(\text{within } \uval).
\end{align*}

\textbf{What ``calibration-free'' claims.} The calibration-free 1-D closed
form is $\EsdClosedVOne\sigma$ (V1) and $\EsdClosedVThree\sigma$ (V3) above
the measurement, consistent within $2\sigma$ but not a match. The 3-D outlet solve is
closer: within $1\sigma$ on the $\mathrm{SD}$ scale and within $2\sigma$ on
the stricter $\mathrm{SEM}$ scale ($\EsemComsolVOne\sigma$ /
$\EsemComsolVThree\sigma$), and
$\alpha_{\mathrm{3D}}=\alphaThreeDVone/\alphaThreeDVthree$ quantifies the
outlet-resistance correction required to reconcile the two. The claim is
that the yield law requires \emph{no device-specific fitting}, not that the
one-dimensional form reproduces the measurement. State both distances
wherever one is stated. \emph{Whether this yield-side residual and the
cutoff-side residual of
Section~\xref{app:units_cutoff}{``Cutoff-Law Prefactor Decomposition''}
share a common three-dimensional origin is an open question and is deferred
to future work}; nothing in this appendix turns on the answer, and the two
are treated throughout as corrections on independent axes.

%==============================================================
% A.7 -- SELECTIVITY-AXIS VALIDATION  (complete, corrected)
%==============================================================
\subsection{Selectivity-Axis Validation (predicted vs.\ measured cutoff)}
\label{app:units_selval}

The yield axis is validated directly against a measured mean with a stated
dispersion (Section~\xref{app:units_validation}{``Validation and
Uncertainty Quantities''}). The selectivity axis is independently checked
using published devices that took no part in the $\kapp$ fit,
\textbf{although the available checks are materially weaker than the yield
validation}: the Dinh check is genuinely held out but resolves only a wide
interval bracket, and the Xia check is not independent of the calibration
at all. \textbf{Neither delivers a sharp held-out validation}, and both are
reported with the reasons why.

\textbf{Dinh 2024: weak interval-consistency check.} The framework predicts
a cutoff band $[\dinhBandLo,\,\dinhBandHi]~\si{\micro\meter}$ for the Dinh
device: the registered band, recomputed by the derivation record~\cite{b12a} and
carried throughout the formal argument. \textbf{Its width is an ambiguity,
not a gradient}, and the distinction matters because it is easy to read the
band as a cutoff that varies down the device. It does not, or at least the
lineage design says it should not.

\emph{Why the taper does not move the cutoff.} The centre channel narrows
for a reason unrelated to selectivity: Xia introduced the taper ``in order
to reduce the maximum amount of shear stress experienced by particles
within the device'' (\cite[p.~4 (Sci Rep PDF pagination)]{b2}), which is what the RS in RS-CIF
stands for. Narrowing $w_c$ alone would shrink the lamina and drag the
cutoff down with it, so the lineage ramps the extraction to compensate:
``In order to maintain a consistent c.d.\ along the length of an RS-CIF
device, the value of $\fgap$ can no longer be kept constant and must be
scaled relative to'' the narrowing (\cite[p.~9 (Sci Rep PDF pagination)]{b2}).
That scaling is Xia's eqn.~(1), stated here in its \emph{corrected} form,
$\fgap(i)=(w_c(\mathrm{ref})/w_c(i))^2\,\fgap(\mathrm{ref})$. The 2016
printing (\cite[p.~10 (Sci Rep PDF pagination)]{b2}) inverted the two
widths, and the 2022 Correction records that the article ``contains an
error in equation 1, where the terms'' $w_c(i)$ and $w_c(\mathrm{ref})$
``are inverted'' (\cite[p.~1]{b3}). \textbf{The direction is load-bearing
here}, so the version matters: with the widths inverted,
$\skimw\propto w_c(i)^2$ and the lamina shrinks \emph{faster} than the bare
taper, which is the opposite of compensation. In the corrected form it
cancels the taper exactly:
\[
\skimw(i)=w_c(i)\sqrt{\tfrac{\fgap(i)}{3}}
        =w_c(i)\,\frac{w_c(\mathrm{ref})}{w_c(i)}\sqrt{\tfrac{\fgap(\mathrm{ref})}{3}}
        =w_c(\mathrm{ref})\sqrt{\tfrac{\fgap(\mathrm{ref})}{3}},
\]
independent of $w_c$. A tapered device is designed to hold one cutoff, not
a range of them.

\emph{What the band's width therefore is.} For the bracket device Dinh
states a single initial $\fgap=1.04\times10^{-4}$, one of three initial
values its designs span (\cite[p.~917]{b5}), and describes such a value as
one of ``the
initial values for the width of the middle channel and the filtration
fraction'' (\cite[p.~915]{b5}); the paper states that the later values are
``iteratively adjusted with each subsequent gap $i$ \ldots\ to maintain the
desired size cut-off throughout the device'' (\cite[p.~916]{b5}), citing Xia
2016 for the method, but does not print them. The band is computed holding it constant, which makes
$\skimw\propto w_c$ and spreads the prediction across the taper. Of the
band's $\dinhBandWidth\times$ width, $75/52=1.44\times$ is that spread and
only $\dinhKneeRatio\times$ is the near-knee $\kapp$ ambiguity, where the fitted prefactor
has ramped off the plateau $\kg$ to a local $\dinhKappaLocal$. Under the
Xia-scaled reading, with $\fgap$ pinned at the initial station, the lamina,
and with it the cutoff at the plateau prefactor, is a single value and the
band collapses to $[\dinhCollapseLo,\,\dinhBandHi]$, the
$\kapp$ ambiguity alone. \emph{The wider band is retained because it is the
conservative envelope}: it contains the Xia-scaled reading, which is the one
Dinh describes, in its corrected form (Appendix~B, C5), and taking the
narrower one would assume that a schedule the paper describes but does not
print meets its stated intent exactly.

The measured rigid-bead bracket is $(1.0,\,4.7)~\si{\micro\meter}$ (Dinh
ESI Fig.~S1: $1~\si{\micro\meter}$ beads track the fluid,
$4.7~\si{\micro\meter}$ beads are retained). The predicted band lies
\emph{inside} the measured bracket on either reading. \textbf{This is a
weak test and must be reported as one:} the bracket is $4.7\times$ wide,
wider than the prediction it is supposed to discriminate; the device
operates near the knee, where held at its initial gap fraction the
self-consistent cutoff sits at the sub-steric foot, beneath the bracket, so
the band is anchor-contingent; and
the band is doubly sensitive to the assumed calibration depth and to the
unresolved $f_{\mathrm{eff}}$. It is a consistency check, not a sharp
validation.

\textbf{Xia 2016: cross-paper consistency check.} \emph{Recalculation} from
three Xia devices, using their stated cutoff and $\fgap$ values, gives
$\kapp=\xiaKappaA/\xiaKappaB/\xiaKappaC$ (mean $\xiaKappaMean$) at $140~\si{\micro\meter}$ depth
against the Gifford plateau $\kappVal$ fitted at an assumed
$80~\si{\micro\meter}$, which agrees to $\xiaAgreePct\%$ across papers and depths.
All three device triples are quoted data: OR-7 at
$\fgap=5.8\times10^{-4}$ for a $7~\si{\micro\meter}$ cutoff, and RS-8 and
RS-9 at $7.2\times10^{-4}$ and $8.6\times10^{-4}$ for $8$ and
$9~\si{\micro\meter}$, all at $w_c(\text{ref})=100~\si{\micro\meter}$, are
each stated verbatim in the source.

\textbf{This anchor is not independent, and the agreement is not
diagnostic.} Two separate reasons, and the second is the stronger one.
First, Xia reads its cutoff from the same Gifford figure the plateau was
fitted from, and Gifford is a Xia co-author, so the anchor and the fit
share a figure rather than merely a literature. Second, and decisively:
the $\xiaAgreePct\%$ by which the recomputed mean exceeds $\kappVal$ is
numerically the depth correction itself. Xia's devices are
$140~\si{\micro\meter}$ deep against the plateau's assumed
$80~\si{\micro\meter}$, and $1/\Kdth(1.40)-1=\depthCorrPct\%$, which is the
same correction as the $\xiaAgreePct\%$ agreement above (the two printed
figures differ only because they fall either side of a rounding boundary at
full precision), so
$\kappVal/\KdepthAnalyticMid$ reproduces the recomputed mean to three
significant figures. The observed agreement is therefore \emph{equally
consistent} with a depth-robust prefactor and with a depth correction of
exactly that size: it discriminates between them not at all. The
recalculation is real and rests on quoted data, but it is reported here as
a \emph{lead}, not as a cross-paper agreement, and the argument records
\texttt{xia\_anchor\_unverified} as open.

\textbf{Numerical uncertainty here is bounded differently.} The single-gap
CFD behind $\Kin$ and the $\Kdcfd$ anchors has no $h$-refinement study.
Its resolution rests instead on agreement with the exact analytic duct law
to $\le0.6\%$ at every point and aspect ratio (the depth factors themselves
agree to $\KdepthDevShallow/\KdepthDevMid/\KdepthDevDeep\%$), supported by mass
conservation to $0.015\%$ and an export-grid check that moves $\skimw$ by
$\le0.1\%$ under an $8\times$ coarsening of the \emph{sampling} pitch. A
sub-micron skim that lands on an analytic value to six parts in a
thousand cannot be under-resolved, so the discretisation error on $\skimw$
is bounded near $0.6\%$: some $5\times$ below the $\pm\kappCIhalfPct\%$ confidence
interval on $\kapp$ it feeds, and finer than the $4.7\times$-wide measured
bracket above can resolve.

\textbf{Fitted-prefactor uncertainty.} $\kapp$ carries a 95\% CI
$[\kappCIlo,\kappCIhi]$ and $\mathrm{CV}=\kappCV\%$ over the
$n=\plateauN$ plateau devices. The single global power law is rejected at
reduced $\chi^2\approx\chiGlobal$, robust across three independent
re-digitizations; that rejection is the falsifiable content of the
functional-form test. \emph{The plateau-only fit is itself
$\chi^2_{\mathrm{red}}\approx\chiPlateau$, which is not a good fit in absolute
terms}; the claim is relative, that $\chiGlobal$ against $\chiPlateau$ is a factor-$\chiRatio$
degradation which rejects the global law.
%==============================================================
% A.8 -- VALIDITY-ENVELOPE AND EXTERNAL INPUTS  (complete, corrected)
%
% STRUCTURE: converted from tabularx-in-a-float to longtable. As submitted this
% table reported "Float too large for page by 2739pt", one "X Columns too narrow",
% and 283 overfull \hbox warnings: a float must fit on one page and this one is
% several pages tall, and a bare `l' column beside an X column squeezes X below its
% minimum. The p{} fractions sum to 0.92 to budget the inner gutters.
%==============================================================
\subsection{Validity-Envelope and External Inputs}
\label{app:units_envelope}

The design framework operates inside a declared envelope. These quantities
are the boundary conditions on that envelope, and several are
\textsc{external}: supplied, and undischargeable by any validator in this
work. An appendix that catalogued the framework without them would omit
the conditions under which its own claims hold.

\begingroup\footnotesize\setlength{\tabcolsep}{4pt}
\begin{longtable}{@{}>{\raggedright\arraybackslash}p{0.155\textwidth}
>{\raggedright\arraybackslash}p{0.475\textwidth}
>{\raggedright\arraybackslash}p{0.115\textwidth}
>{\raggedright\arraybackslash}p{0.175\textwidth}@{}}
\caption{Validity-envelope quantities and external
inputs.}\label{tab:units_envelope}\\
\toprule Symbol & Role & Status & Source \\\midrule\endfirsthead
\multicolumn{4}{@{}l}{\footnotesize\itshape Table \thetable\ (continued)}\\
\toprule Symbol & Role & Status & Source \\\midrule\endhead
\bottomrule\endlastfoot
$u_{\mathrm{lo}}$ & Plateau guard: every row must satisfy $\ud=\dcrit/w_c\ge u_{\mathrm{lo}}$, keeping the design above the knee. \emph{A threshold on $\ud$, not a velocity and not a V\&V uncertainty}; see the $u$ note below & \textsc{external} & This work (supplied) \\
$u_{\mathrm{ff}}$ & Free-flow guard on the reduced size ratio $\ud$, same letter caveat & \textsc{external} & This work (supplied) \\
$\eta$ & Realized-yield factor $\eta=\etaVal$: a three-dimensional correction on the \emph{flow split}. \textbf{Not} an external free parameter: its shortfall $1-\eta=\etaShortfallPct\%$ is the same ${\sim}10\%$ one-dimensional-versus-three-dimensional residual on $F$ that the outlet solve closes, so it is a computed magnitude recorded here rather than a supplied constant. \emph{A single realized-yield magnitude, not a per-variant fit}: the actual residual is $\residPctVOne\%$ at V1 and $\residPctVThree\%$ at V3, so $\etaShortfallPct\%$ is representative and matches V1 closely rather than holding exactly at both. \textbf{Not} the lamina model gap either, which lives in $\kapp$; the two act on different axes and must not be traded against one another & \textsc{measured} (corroborated by the 3-D COMSOL outlet solve to within $1\sigma$, not derived by the one-dimensional treatment) & \cite{b12a}, \texttt{uncheckable\_eta} \\
$\phi_{\mathrm{cal}}$ & Calibration wall-porosity input, distinct from $\phi=G/\Lambda$, which is \emph{computed} from geometry; $\phi_{\mathrm{cal}}$ is supplied and no validator discharges it. See the $\phi$ note below & \textsc{external} & This work (supplied); \cite{b12a} \\
$\sigma_0$ & Fictitious inlet side-flow in the mass identity $P(i)+2S(i)=1+2\sigma_0$. \emph{Defined in Table~\xref{tab:units_fluid}{A.2}}; repeated here only to record that it is supplied & \textsc{external} & Table~\xref{tab:units_fluid}{A.2} \\
$f_{\mathrm{soft}}$, $\gamma_{\mathrm{GS}}$ & Schedule-softening and gap-scale process parameters. $f_{\mathrm{soft}}$ is a \emph{process} parameter and is not the filtrate fraction $\fgap$; see the $f$ note below & \textsc{external} & This work (supplied) \\
$\mathrm{AR}_{\max}$, $R_{\min}$ & Fabrication limits: maximum feature aspect ratio and minimum inner-corner radius & \textsc{external} & This work (supplied) \\
$c_{\mathrm{ent}}$, $L_{\mathrm{eq}}$ & Entrance coefficient and bend-equivalent length. \emph{The link length $\ell_g$ is defined in Table~\xref{tab:units_geom}{A.1}} and is not redefined here & \textsc{external} & This work (supplied) \\
\midrule
\multicolumn{4}{@{}p{0.93\textwidth}}{\emph{Three conditions a design must satisfy. Each is stated as a boundary rather than asserted to hold, because a real device can violate any of them; asserting they always hold would claim the device is always above the knee, which is false.}}\\
\midrule
Above-knee operation & If the device runs below the knee, $\kapp$ ceases to be a constant and becomes a function of the cutoff itself, making the design relation implicit & \textsc{bounds} & This work \\
Run-out match & A $3\times$ run-out mismatch collapses $F$ from $\runoutFmatched$ to $\runoutFmismatched$; the outlet lengths must be matched & \textsc{bounds} & \cite{b12a}, \texttt{runout\_mismatch} \\
Free row count & Pinning $N$ over-constrains the problem: two targets meet one free variable, and the orthogonality of cutoff and yield is lost & \textsc{bounds} & This work \\
\end{longtable}\endgroup

\textbf{$u$-symbol collision, extended.} Table~\xref{tab:units_validation}{A.6}
records three meanings for the letter $u$: the fluid velocities $u(y)$ and
$u_{\max}$, the V\&V~20 uncertainties $\uval$, $\unum$, $\uinput$, $\uD$,
and the reduced size ratio, which was moved \emph{off} the letter to
$\ud=\widetilde{d}$ precisely because $u_d$ and $u_D$ differ by subscript
case alone. The guards $u_{\mathrm{lo}}$ and $u_{\mathrm{ff}}$ are a
fourth use, and an awkward one: they are thresholds \emph{on} $\ud$, so
they are named after the very symbol this appendix stopped using. They keep
their names because they are the identifiers the design code and the
specification use, and renaming them here would break that correspondence,
but the reader should read $u_{\mathrm{lo}}$ and $u_{\mathrm{ff}}$ as
bounds on $\widetilde{d}$ and never as velocities or uncertainties.

\textbf{$f$-symbol collision.} $\fgap$ is the per-gap filtrate fraction and
is the central dimensionless quantity of the yield axis
(Table~\xref{tab:units_dimensionless}{A.3}). $f_{\mathrm{soft}}$ is a
schedule-softening \emph{process} parameter with no flow meaning, and
$F$ with its variants is the cumulative filtrate fraction. Three
$f$-shaped objects, distinguished by subscript and by case.

\textbf{$\phi$-symbol collision.} $\phi_{\mathrm{cal}}$, the wall porosity
\emph{supplied} to the calibration, is distinct from $\phi=G/\Lambda$, which
is \emph{computed} from the geometry (Table~\xref{tab:units_geom}{A.1}), and
and from the upper-case fold factor $\Phi(i)=r(w_s(i))/r(\mathrm{GS})$ of
the folded-channel geometry (Table~\xref{tab:units_geom}{A.1}): one letter,
separated by subscript and by case.

\textbf{$c$-symbol collision.} The letter $c$ carries three unrelated
meanings. In the cutoff decomposition $\sqrt{c}\,C$
(Table~\xref{tab:units_cutoff}{A.4}) lower-case $c$ is the per-device
over-extraction factor and upper-case $C$ the cumulative finite-particle
factor, distinguished by case alone. $c_2$ is the near-wall coefficient of
the velocity expansion (Table~\xref{tab:units_dimensionless}{A.3}). And
$c_{\mathrm{ent}}$, in this table, is a dimensionless entrance-length
coefficient with no relation to either. All three are dimensionless, so
unit disagreement will not separate them; the subscript is the only cue.

\textbf{$L$-symbol note.} $L_{\mathrm{eq}}$, the bend-equivalent length, is
a member of the length family alongside $L_{\mathrm{dev}}$, $L_c(i)$ and
$L_s(i)$ (Table~\xref{tab:units_geom}{A.1}): all are lengths in metres and
all carry a subscript naming what they measure.

\textbf{On the} \textsc{external} \textbf{label.} In this table it has the
second of the two senses Section~\xref{app:units_geom}{A.1} gives it: it marks
a quantity this work supplies rather than derives, and which no validator in the record
discharges. That is a statement about \emph{this} record, not a claim that
the quantity is underivable: $\phi_{\mathrm{cal}}$ and the fabrication
limits are physically measurable, and a later campaign that measured them
would move them out of this table. The label is a boundary on what the
present argument establishes.

%==============================================================
% A.9 -- EQUATIONS: STEP-BY-STEP UNIT DERIVATIONS  (complete, corrected)
%
% CHANGES FROM THE SUBMITTED DRAFT, each justified in the report:
%   * depth-factor denominator corrected to |Kdcfd - Kdth|/Kdcfd. The draft wrote
%     /Kdth while printing the /Kdcfd values, the same formula-versus-value
%     mismatch A.4 carried and had corrected.
%   * the CFD deviations and fourteen other computed figures are now macros.
%   * every cross-subsection \ref became \xref, so this file builds standalone.
%   * the self-reference (Section \ref{app:units_equations} pointing at itself)
%     now names the subsubsection that actually quotes the band.
%   * the cutoff law and its decomposition are displayed UNNUMBERED here and
%     cited to A.4, which already numbers both. One equation, one number.
% This list is the change set FROM SUBMISSION, not a running log; later rounds
% corrected the lineage f_gap maximum, attributed the row pitch to Gifford's
% measured Fig.4 value, and anchored the hydraulic-electric analogy.
%==============================================================
\subsection{Equations: Step-by-Step Unit Derivations}
\label{app:units_equations}

This section provides a dimensional \emph{consistency check} for each
load-bearing equation. The notation $[X]$ denotes ``the SI unit of $X$.''
\emph{Dimensional consistency is necessary, not sufficient}: it shows an
equation is well formed, not that its physical assumptions hold. The
assumptions are stated separately with each equation.

\subsubsection{Gifford 2014 Eqn~1: flow recursion}

\begin{equation}
Q_s(i) = Q_s(i-1) + \fgap(i) \cdot Q_c(i-1)
\label{eqn:gifford1}
\end{equation}

\emph{Dimensional analysis:}
\begin{align*}
[\text{LHS}] &= \si{\meter\cubed\per\second} \\
[\text{RHS}] &= \si{\meter\cubed\per\second} + \underbrace{[\fgap]}_{1}\!\cdot\!\underbrace{[Q_c]}_{\si{\meter\cubed\per\second}}
              = \si{\meter\cubed\per\second}
\end{align*}
$[\text{LHS}] = [\text{RHS}].$~\checkmark
\emph{Index note:} Gifford writes Eqn~1 with no row index on
$\fgap$ (his constant-$\fgap$ regime); the $(i)$ here is the
tapered-$\fgap$ generalization. $Q_c(i)$ and $Q_s(i)$ are the flows
following gap row~$i$, as Gifford defines them, so $\fgap(i)$ is the
fraction taken at row~$i$ from the central flow reaching it, $Q_c(i-1)$,
which is Dinh's $Q_m(i)$ (\cite{b5}). Every equation in this appendix
uses this labelling.

\subsubsection{Gifford 2014 Eqn~3: continuity}

\begin{equation}
\frac{\Delta P(i)}{\Delta P(i-1)} = 1 - 2\fgap(i)
\label{eqn:gifford3}
\end{equation}

\emph{Dimensional analysis:} $[\text{LHS}]=\si{\pascal}/\si{\pascal}=1$;
$[\text{RHS}]=1-2\cdot1=1$.~\checkmark The factor of $2$ is bilateral
extraction (each of two side channels takes $\fgap$ of $Q_c$).
\emph{Index note:} Gifford's verbatim form carries no index on $\fgap$;
the $(i)$ is the tapered generalization.

\subsubsection{Gifford 2014 Eqn~5: Yang series for hydraulic resistance}

\begin{equation}
R = \frac{12 \mu L}{w_{\mathrm{ch}} h^3}
   \left[
     1 - \frac{192\,h}{\pi^5 w_{\mathrm{ch}}}
         \sum_{n=1,3,5,\ldots}^{\infty} \frac{1}{n^5}
         \tanh\!\left(\frac{n \pi w_{\mathrm{ch}}}{2h}\right)
   \right]^{-1}
\label{eqn:gifford5}
\end{equation}

\emph{Dimensional analysis} (convention $h \le w_{\mathrm{ch}}$):
prefactor $[12\mu L]/[w_{\mathrm{ch}} h^3]
= \si{\pascal\second\meter}/\si{\meter^4}
= \si{\pascal\second\per\meter\cubed}$; the bracket is dimensionless, so
$[R] = \si{\pascal\second\per\meter\cubed}$.~\checkmark

\textbf{Attribution.} Gifford's Eqn~5 carries citation [20] $=$
Yang/\"{U}ndar/Zahn (\cite{b6}), whose Eqn~3 is the rectangular-duct
series; Gifford's ref [19] (Sethu/Sin/Toner) is the \emph{simplified}
$R(w)$ treatment Eqn~5 improves on, not its source.

\textbf{Limit checks.} Parallel plates ($h\ll w_{\mathrm{ch}}$): series
$\to 0$, $R\to12\mu L/(w_{\mathrm{ch}} h^3)$.~\checkmark Square channel
($h=w_{\mathrm{ch}}$): $R\to 28.4\,\mu L/h^4$ (Bruus 2008, \cite{b20}).
Independently corroborated by the registry's square-duct anchor
$12/B(1) = \squareDuctR$, computed from the duct series with no reference
to \cite{b20}.~\checkmark
\emph{Note:} the channel-width argument is $w_{\mathrm{ch}}\in\{w_c,w_s\}$,
not the lamina width $w$.

\subsubsection{Dinh 2024 p.~915: definition of $\fgap$}

\begin{equation}
Q_{\mathrm{gap}}(i) = \fgap(i) \cdot Q_c(i-1)
\label{eqn:dinh_fgap}
\end{equation}

\emph{Dimensional analysis:} $\fgap(i)=Q_{\mathrm{gap}}(i)/Q_c(i-1)$,
units $\si{\meter\cubed\per\second}/\si{\meter\cubed\per\second}=1$.~\checkmark
\emph{Index note:} Dinh writes the flow reaching gap~$i$ as $Q_m(i)$
(\cite{b5}); in Gifford's labelling, used throughout this appendix, it
is $Q_c(i-1)$.

\textbf{Numerical sanity (Dinh V1):} $Q_c\approx30~\si{\micro\liter\per\minute}$,
$\fgap=1.04\times10^{-4}\Rightarrow Q_{\mathrm{gap}}=\QgapDinh~\si{\nano\liter\per\minute}$
per gap per side.

\subsubsection{Xia 2022 corrected scaling for matched cutoffs}

\begin{equation}
\fgap(i) = \left( \frac{w_c(\text{ref})}{w_c(i)} \right)^2 \fgap(\text{ref})
\label{eqn:xia_corrected}
\end{equation}

\emph{Dimensional analysis:} ratio squared is dimensionless; both sides
dimensionless.~\checkmark At fixed cutoff: $\fgap\,w_c^2=$ const, i.e.\
$\dcrit\propto w_c\sqrt{\fgap}$, consistent with the cutoff law of
Section~\xref{app:units_cutoff}{``Cutoff-Law Prefactor Decomposition''}.
\emph{This is the \textbf{corrected} form}: the 2016 printing inverted
$w_c(i)$ and $w_c(\mathrm{ref})$, and the 2022 Correction records that the
article ``contains an error in equation 1, where the terms'' $w_c(i)$ and
$w_c(\mathrm{ref})$ ``are inverted'' (\cite[p.~1]{b3}). The direction is
load-bearing: inverted, the lamina would shrink \emph{faster} than the
taper instead of holding constant.

\subsubsection{Skim cubic and dividing-streamline width (selectivity)}
\label{sec:skim_band}

\begin{equation}
\fgap = 3x^2 - 2x^3,\qquad x=\skimw/w_c;\qquad
\skimw = w_c\!\left[\sqrt{\fgap/3} + \tfrac{1}{9}\fgap + \cdots\right]
\label{eqn:skim_cubic}
\end{equation}

\emph{Dimensional analysis:} $\fgap$ and $x$ are dimensionless; the cubic
balances $1=1$.~\checkmark $[\skimw]=[w_c]\cdot1=\si{\meter}$.~\checkmark
This is the Inglis $\varepsilon$-cubic (\cite{b17}): the
parabolic-profile flux fraction equals $3x^2-2x^3$. Inglis states the
premise, ``by assuming a conventional parabolic flow profile through the
gap'' (\cite[p.~656]{b17}), and writes the result in monic cube-root form,
$y^3-\tfrac{3}{2}y^2+\tfrac{\varepsilon}{2}=0$ with $y=b/g$; multiplying by
$-2$ and identifying $\varepsilon\!\leftrightarrow\!\fgap$,
$y\!\leftrightarrow\!x$ recovers $\fgap=3x^2-2x^3$ exactly. \emph{The
polynomial is a derived identity, not a quoted formula}, which is why the
source-quote index anchors the premise rather than the cubic.

\textbf{Accuracy of the leading inversion, with the band named.} The
leading inversion $\skimw=w_c\sqrt{\fgap/3}$ is within $0.5\%$ of the
exact root across the Gifford \emph{calibration/design} band
($\fgap\le\fgapDesignCap$, device~\#33; worst case $\skimErrDesign\%$),
and within $0.9\%$ across the wider \emph{numerical sweep} band
($\fgap=\fgapBandLo$ to $\fgapBandHi$; worst case $\skimErrBand\%$ at the
top). Quote the band with the figure: the two differ by nearly a factor
of two.

\subsubsection{Cutoff law and prefactor decomposition (selectivity)}

The cutoff law and its prefactor decomposition are numbered in
Section~\xref{app:units_cutoff}{``Cutoff-Law Prefactor Decomposition''}
(Eqns~\xref{eqn:cutoff_law}{6} and~\xref{eqn:kappa_decomp}{7}) and are
restated here \emph{unnumbered}, for dimensional analysis only, so that
one equation carries one number:
\[
\dcrit = \kapp\, w_c\, \sqrt{\fgap/3},\qquad
\kapp = \kster\,\Kin\,\Kdepth\,\sqrt{c}\,C .
\]

\emph{Dimensional analysis:} $\kapp$ is dimensionless (product of
dimensionless factors), $[\sqrt{\fgap/3}]=1$, so
$[\dcrit]=[w_c]=\si{\meter}$.~\checkmark Each factor is a pure number
(Table~\xref{tab:units_cutoff}{A.4}); $\kster=\kstericVal$,
$\Kin=\KinVal$, and the depth-folded residual
$\Kdepth\sqrt{c}\,C=\kapp/\kster=\residualFolded$ (depth-separated
$\sqrt{c}\,C=\kapp/(\kster\Kdcfd)=\residualSepShallow/\residualSepMid/\residualSepDeep$
at $\alpha=0.80/1.40/2.69$).

\subsubsection{Depth-correction factor: analytic and CFD}

\begin{equation}
\Kdth(\alpha) = \sqrt{\frac{3}{c_2(\alpha)}}
 = \lim_{x\to0}\frac{\skimw_{3D}(\alpha)}{\skimw_{2D}},
\qquad
\frac{|\Kdcfd-\Kdth|}{\Kdcfd}\le 0.6\%
\label{eqn:depth_factor}
\end{equation}

\emph{Dimensional analysis:} ratio of two widths, dimensionless.~\checkmark
$\Kdth$ is evaluated from the rectangular-duct velocity field of aspect
ratio $\alpha=h/w_{\mathrm{ch}}$; $\Kdcfd$ is the single-gap CFD
measurement. \textbf{These are two distinct quantities and must not be
averaged or interchanged.} Analytic:
$\KdepthAnalyticShallow/\KdepthAnalyticMid/\KdepthAnalyticDeep$. CFD:
$\KdepthCfdShallow/\KdepthCfdMid/\KdepthCfdDeep$. Actual deviations
$\KdepthDevShallow\%$, $\KdepthDevMid\%$, $\KdepthDevDeep\%$; the $0.6\%$
figure is the \emph{bound}, not the observed value. \emph{The denominator
is $\Kdcfd$}, the measured factor, because the residual is defined against
the measurement; dividing by $\Kdth$ instead shifts the middle deviation
from $\KdepthDevMid\%$ to $\KdepthDevMidTh\%$ and would no longer match the gated
value. The analytic set is canonical for derivations; the CFD set is
canonical for the depth-separated residual, which is defined against it.

\subsubsection{Filtrate fraction: two identities, two roles}

The filtrate fraction has two representations. They evaluate to the same
realized $F$ for a self-consistent device satisfying the model
assumptions, but they expose different dependencies and should not be
written as one undifferentiated chain of equalities.

\paragraph{Outlet divider (yield law).}
\begin{equation}
F_{\mathrm{outlet}} = \frac{2 R_c}{R_s + 2 R_c}
\label{eqn:closed_form}
\end{equation}

\emph{Dimensional analysis:} numerator and denominator both
$\si{\pascal\second\per\meter\cubed}$, so $F$ is dimensionless.~\checkmark

\textbf{Derivation under PE (P1.4) and $Q_s(0)=0$ (P1.5).} With outlet
pressure equilibration $p_{\text{center}}(L)=p_{\text{side}}(L)$, the
current-divider rule on the equivalent network (two side channels $R_s$ in
parallel, sharing the outlet node with the center channel $R_c$) gives
$Q_s(\text{one side}) = R_c Q_{\mathrm{in}}/(R_s+2R_c)$, hence
$F = 2Q_s/Q_{\mathrm{in}} = 2R_c/(R_s+2R_c)$. The hydraulic--electric
analogy underpinning this is \cite{b13,b14}.

\textbf{Decoupling.} $\fgap$ does not appear in
Eqn~\ref{eqn:closed_form}: $F_{\mathrm{outlet}}$ is fixed by outlet
geometry ($R_c,R_s$) alone, while the cutoff is fixed locally by
$\skimw=w_c\sqrt{\fgap/3}$. \emph{This absence is the decoupling.} The
two axes are therefore separately addressable, and the orthogonality holds
\emph{because $N$ is free}; pinning $N$ re-couples them.

\paragraph{Cumulative march (internal extraction history).}
\begin{equation}
F_{\mathrm{recursion}} = 1 - \prod_{i=1}^{N} \bigl[ 1 - 2 \fgap(i) \bigr]
\label{eqn:F_product}
\end{equation}

\emph{Dimensional analysis:} each factor dimensionless; product
dimensionless.~\checkmark Small-$\fgap$ limit:
$F\approx1-\exp(-2\sum_i\fgap(i))$; constant $\fgap$:
$F\approx1-\exp(-2N\fgap)$. This form describes how cumulative extraction
unfolds along the ladder; it does \emph{not} exhibit the decoupling,
because $\fgap$ appears explicitly.

\subsubsection{Hagen--Poiseuille / Ohm's-law analogue}

\begin{equation}
\Delta P = R \cdot Q
\label{eqn:HP_ohm}
\end{equation}

\emph{Dimensional analysis:}
$[R\cdot Q]=\si{\pascal\second\per\meter\cubed}\cdot\si{\meter\cubed\per\second}=\si{\pascal}=[\Delta P]$.~\checkmark
The load-bearing identity of the hydraulic--electrical analogy
($P\!\leftrightarrow\!V$, $Q\!\leftrightarrow\!I$, $R\!\leftrightarrow\!R$):
``pressure drop, flow rate, and flow resistance in microchannels are
analogous to electrical voltage, current, and resistance, respectively''
(\cite[p.~3312]{b14}), and the circuit-analogy design paper states the same
mapping as ``correlations of pressure to voltage, volumetric flow rate to
current, and hydraulic to electric resistance'' (\cite[p.~515]{b13}).

\subsubsection{Reynolds number and flow regime}

\begin{equation}
\text{Re} = \frac{\rho U L_{\mathrm{char}}}{\mu}
\label{eqn:reynolds}
\end{equation}

\emph{Dimensional analysis:} dimensionless (using
$\si{\pascal\second}=\si{\kilo\gram\per\meter\per\second}$).~\checkmark

\textbf{CIF regime check (Dinh, water-like fluid).}
Dinh's channel depth is $h=140~\si{\micro\meter}$ and the middle channel
narrows to $w_c\approx52~\si{\micro\meter}$ at the outlet, giving a
hydraulic diameter
$D_h = 2w_c h/(w_c+h) \approx \DhOutlet~\si{\micro\meter}$. At
$Q=30~\si{\micro\liter\per\minute}$ through the outlet cross-section (the
whole flow through the narrowest section, a worst case that bounds the
design code's local-flow value), $U=Q/A\approx\Uoutlet~\si{\meter\per\second}$,
so the \emph{channel-scale}
$\text{Re}=\rho U D_h/\mu \approx \ReWorstChannel$ (rising to
$\approx\ReHighFlow$ at $90~\si{\micro\liter\per\minute}$). PRP is more
viscous than water, so its Reynolds number is lower in proportion; none of
the three papers states PRP's viscosity. The
Reynolds number \emph{at the gap} is obtained not by rescaling the length
but by rescaling the \emph{flux}: a gap carries only the diverted fraction
$\fgap$ of the channel flow, so
$\mathrm{Re}_{\mathrm{gap}}\approx\fgap\,\mathrm{Re}_{\mathrm{channel}}
\approx10^{-3}\times\ReWorstChannel\approx\ReGapScale \ll 1$, evaluated at
$\fgap\approx10^{-3}$, which is the \emph{top} of the published lineage
range rather than a margin above it. Gifford's 33-device \emph{calibration}
array caps lower, at $\fgapDesignCap$ (device~\#33); his separate
bead-separation devices were built with ``a range of values of'' $\fgap$
``$(2.2$--$10.0 \times 10^{-4})$'' (\cite[p.~4500]{b1}), reaching
$1.0\times10^{-3}$, and no published lineage device exceeds it. The
argument record's scope gate evaluates the same product. The design code's
Stokes gate (M14) instead computes the gap's own Reynolds number, over the
gap's $G\times h$ section at the first row's gap flow, which exceeds this
product by $(w_c+h)/(G+h)$ (about $1.2$ for Dinh's $19~\si{\micro\meter}$ gaps
and outlet width), and a lamina-shear $\mathrm{Re}_\delta$, and in a cutoff
design refuses either above $10^{-2}$. \textbf{Do not confuse this with the wider band
$\fgapBandLo$ to $\fgapBandHi$ quoted under ``Skim cubic and
dividing-streamline width'' above}, which is the span of a deliberately
widened numerical sweep used to demonstrate decoupling over a decade, not
a range any built device occupies. \emph{The margin is worth stating
plainly}: at the top of that sweep, $\fgap=\fgapBandHi$, the same scaling
would give $\mathrm{Re}_{\mathrm{gap}}\approx0.010$ (the design code's own
measure about a fifth higher, just past its gate), which is where the
creeping-flow premise stops having room. The premise is comfortable for
built devices and marginal at the top of the test sweep. \emph{(A
length-rescaling estimate gives different numbers and should not be
conflated with this one: at the same $U$, using the
$19~\si{\micro\meter}$ gap width gives $\mathrm{Re}\approx1.3$ and using
the outlet lamina width $\skimw\approx0.31~\si{\micro\meter}$ gives
$\approx0.02$. The flux argument is the one the scope gate uses.)} The
steric-skimming junction is therefore in the creeping (Stokes) regime
while the channel flow is low-Reynolds laminar.
\textbf{The framework is therefore not globally Stokes}: it is scoped to
rigid, approximately spherical particles in a Newtonian, nearly
constant-viscosity fluid, with creeping-flow conditions \emph{at the gap
and lamina scales that govern particle selection}. Deformable particles,
cell crowding and whole-blood rheology are outside the validated envelope.
\emph{(An earlier revision used $L_{\mathrm{char}}=30~\si{\micro\meter}$
and $U=0.22~\si{\meter\per\second}$ and overstated the regime as
``Stokes''; the steric derivation is unaffected because it rests on the
gap-scale $\text{Re}\ll1$.)}

\subsubsection{Entrance length and the quasi-developed assumption}

\begin{equation}
\frac{L_e}{D_h}\simeq\Bigl[0.619^{1.6}+(0.0567\,\mathrm{Re})^{1.6}\Bigr]^{1/1.6}
\ \xrightarrow{\ \mathrm{Re}\to0\ }\ 0.619
\label{eqn:entrance}
\end{equation}

\emph{Dimensional analysis:} ratio of two lengths, dimensionless.~\checkmark

\textbf{The kernels assume fully developed flow, and strictly it is not.}
The correct laminar entrance length has a \emph{Stokes floor} and does not
vanish as $\mathrm{Re}\to0$; the naive $L_e=0.05\,\mathrm{Re}\,D_h$ goes
to zero in this regime and is wrong here. The correlation is Durst \emph{et al.}~\cite{b23} for developing
\emph{pipe} flow, applied here through the hydraulic diameter $D_h$; the
$0.619$ constant is that correlation's low-Reynolds limit. The
rectangular-duct coefficient differs slightly, and \cite{b23} is the one
citation in this section whose source is not held in the verification
corpus, so the constant is taken as published. The conclusion below does
not turn on it.

\textbf{Evaluated on this document's own geometry}, rather than quoted:
at the Dinh inlet ($w_c=75~\si{\micro\meter}$, $h=140~\si{\micro\meter}$,
$D_h=\DhInlet~\si{\micro\meter}$) the floor gives
$L_e\approx\LeInlet~\si{\micro\meter}$; at the outlet
($w_c=52~\si{\micro\meter}$, $D_h=\DhOutlet~\si{\micro\meter}$) it gives
$L_e\approx\LeOutlet~\si{\micro\meter}$. Against the lineage row pitch $\Lambda=G+g_p=46.1~\si{\micro\meter}$, which
the record measures from Gifford's Fig.~4 ($G=20.5$, $g_p=25.6$), this is
$\LeOverLamGInlet$ at the inlet and $\LeOverLamGOutlet$ at the outlet. \emph{Dinh publish no row
pitch}: their Fig.~2 gives the gap width, the taper and the depth but no
post width or spacing, so an estimated Dinh pitch near
$49~\si{\micro\meter}$ (giving $\LeOverLamInlet$ and $\LeOverLamOutlet$)
is an inference, not a source figure, and the traceable Gifford value is
quoted here instead. On either pitch
the ratio exceeds unity at the inlet and sits at unity at the outlet.
\textbf{The entrance length is
therefore comparable to the row pitch and exceeds it over much of the
device: the flow does not classically redevelop between rungs.}

This is expected to be benign, for a stated reason: each rung perturbs the
base flow by a relative amount of order $\fgap$ (about $10^{-4}$, and at most
$10^{-3}$ across the published lineage range), so the error this makes in a
segment's resistance is expected to be small. The design specification puts
that error at $O(\fgap^2)\sim10^{-8}$; that order needs the first-order term
to cancel, which this record does not show, so it is not used here as a
bound. The conclusion does not turn on
the exact value of $L_e$: it is the same whether $L_e/\Lambda$ is $\LeOverLamOutlet$ or
$\LeOverLamGInlet$. Report $L_e/\Lambda$; do not gate
on it. Gate instead on the perturbation, which the design specification identifies with the design
code's gap-dominance gate (M13).

\subsubsection{Platelet volume and the Dinh threshold}

\begin{equation}
V_{\mathrm{sph}} = \frac{4}{3} \pi a^3
\label{eqn:platelet_volume}
\end{equation}

\emph{Dimensional analysis:} $[V_{\mathrm{sph}}]=[a]^3=\si{\meter\cubed}$.~\checkmark
\textbf{Numerical ($12$~fL):} $a\approx\aTwelveFL~\si{\micro\meter}$,
$D=2a\approx\DTwelveFL~\si{\micro\meter}$ (volume-equivalent sphere; a
geometric conversion, not a measured diameter). Dinh report
$2a\approx4.8~\si{\micro\meter}$ for the actual discoid $12$~fL platelet
(\cite[p.~917]{b5}); see the volume-convention note in
Section~\xref{app:units_particle}{``Particle and Separation Variables''}.

\subsubsection{Multi-variant ratio (Dinh V1, V3)}

\begin{equation}
\mathcal{R} = \frac{\ln(1 - F_{V3})}{\ln(1 - F_{V1})}
\label{eqn:multi_variant}
\end{equation}

\emph{Dimensional analysis:} $\ln$ of dimensionless argument;
ratio dimensionless.~\checkmark \textbf{Origin.} With
$F=1-\exp(-2N\fgap)$, $\ln(1-F)=-2N\fgap$ and
$\mathcal{R}=(N_{V3}\,f_{\mathrm{gap},V3})/(N_{V1}\,f_{\mathrm{gap},V1})$,
sensitive only to the cumulative $\fgap$-product.
% =====================================================================
% OPTION A (recommended): fold the unit-consistency summary into A.9 as a
% closing subsubsection. APPEND this to A9_FINAL.tex and DELETE A.12.
%
% Why here rather than as its own subsection: fourteen of the seventeen rows
% point at equations numbered in THIS file, so as a separate subsection they
% are cross-file \ref pointers. That is exactly the fragility that broke the
% build when the cutoff law was unnumbered and app:units_summary still
% referenced the retired label. Folded in, those fourteen become intra-file
% references that resolve in a standalone A.9 build, and only the three
% Appendix-A.4 rows need \xref.
%
% The CONVENTION column is the reason the table is not merely A.9's
% checkmarks reprinted: a dimensional check is only meaningful under a stated
% convention, and three of these carry one that a reader would otherwise have
% to reconstruct from the body.
% =====================================================================
\subsubsection{Unit-consistency summary}
\label{app:units_summary}

Each equation above was checked inline. The table gathers them so the
whole set can be scanned at once, and records the convention each check
runs under, because a dimensional identity that holds only under a
convention is not established until the convention is named. Three rows
come from Section~\xref{app:units_cutoff}{``Cutoff-Law Prefactor
Decomposition''} and are checked here rather than there.

\begingroup\small\setlength{\tabcolsep}{4pt}
\begin{longtable}{@{}>{\raggedright\arraybackslash}p{0.30\textwidth}
>{\raggedright\arraybackslash}p{0.17\textwidth}
>{\raggedright\arraybackslash}p{0.17\textwidth}
>{\raggedright\arraybackslash}p{0.28\textwidth}@{}}
\caption{Unit-consistency check for every load-bearing equation, both axes
and the validation layer, with the convention each check assumes. Every
entry matches; the column that carries information is the
last.}\label{tab:units_check}\\
\toprule Equation & LHS & RHS & Convention assumed \\\midrule\endfirsthead
\multicolumn{4}{@{}l}{\footnotesize\itshape Table \thetable\ (continued)}\\
\toprule Equation & LHS & RHS & Convention assumed \\\midrule\endhead
\bottomrule\endlastfoot
Eqn~\ref{eqn:gifford1} (Gifford flow recursion) & \si{\meter\cubed\per\second} & \si{\meter\cubed\per\second} & $\fgap$ carries the row index; Gifford's own form has none \\
Eqn~\ref{eqn:gifford3} (Gifford continuity) & 1 & 1 & bilateral extraction, hence the factor $2$ \\
Eqn~\ref{eqn:gifford5} (Yang series) & \si{\pascal\second\per\meter\cubed} & \si{\pascal\second\per\meter\cubed} & $h\le w_{\mathrm{ch}}$; $w_{\mathrm{ch}}\in\{w_c,w_s\}$, not the lamina $w$ \\
Eqn~\ref{eqn:dinh_fgap} (Dinh $\fgap$ definition) & \si{\meter\cubed\per\second} & \si{\meter\cubed\per\second} & per gap, per side \\
Eqn~\ref{eqn:xia_corrected} (Xia corrected scaling) & 1 & 1 & the \emph{corrected} form; the 2016 printing inverts the widths \\
Eqn~\ref{eqn:skim_cubic} (skim cubic / $\skimw$) & 1 \;/\; \si{\meter} & 1 \;/\; \si{\meter} & two statements in one display: the cubic is dimensionless, $\skimw$ is a length \\
Eqn~\xref{eqn:cutoff_law}{6} (cutoff law, $\dcrit$) & \si{\meter} & \si{\meter} & leading skim $w_c\sqrt{\fgap/3}$, the form $\kg$ was fitted with \\
Eqn~\xref{eqn:kappa_decomp}{7} ($\kapp$ decomposition) & 1 & 1 & every factor a pure number \\
Eqn~\xref{eqn:kernel_pairing}{8} (kernel--prefactor conversion) & 1 & 1 & \emph{array-mean} $\overline{c_2}$ over the calibration widths, not a per-device $c_2$ \\
Eqn~\ref{eqn:depth_factor} (depth factor $\Kdth$) & 1 & 1 & deviation denominator is $\Kdcfd$, the measured factor \\
Eqn~\ref{eqn:closed_form} ($F_{\mathrm{outlet}}$) & 1 & 1 & outlet pressure equilibration and $Q_s(0)=0$ \\
Eqn~\ref{eqn:F_product} ($F_{\mathrm{recursion}}$) & 1 & 1 & cumulative over $N$ rows; not the outlet identity \\
Eqn~\ref{eqn:HP_ohm} (Hagen--Poiseuille analogue) & \si{\pascal} & \si{\pascal} & $P\!\leftrightarrow\!V$, $Q\!\leftrightarrow\!I$, $R\!\leftrightarrow\!R$ \\
Eqn~\ref{eqn:reynolds} (Reynolds number) & 1 & 1 & channel scale; the gap value rescales the \emph{flux}, not the length \\
Eqn~\ref{eqn:entrance} (entrance length $L_e/D_h$) & 1 & 1 & the $\mathrm{Re}\to0$ Stokes floor, $0.619$ \\
Eqn~\ref{eqn:platelet_volume} (sphere volume) & \si{\meter\cubed} & \si{\meter\cubed} & volume-equivalent sphere, not a measured discoid diameter \\
Eqn~\ref{eqn:multi_variant} (multi-variant ratio) & 1 & 1 & constant-$\fgap$ exponential limit \\
\end{longtable}\endgroup

\noindent All seventeen are unit-consistent. \emph{Dimensional consistency
is necessary and not sufficient}, and the conventions in the last column
are where the substantive assumptions live: the check confirms each
equation is well formed under them, not that they hold.

%==============================================================
% A.10 -- SOURCE-QUOTE VALIDATION INDEX  (complete, corrected)
%==============================================================
\subsection{Source-Quote Validation Index (verbatim)}
\label{app:units_quotes}

Direct quotations anchoring each unit convention and adjudicated
criterion, organized by source. Quotes are minimal; the
page/figure/equation locator does the identifying work. A consolidated
claim$\rightarrow$source$\rightarrow$anchor table follows
(Table~\ref{tab:verbatim}).

\paragraph{What needs a quotation, and what does not.} Not every statement
in this appendix should carry one, and forcing quotations onto
mathematical or computed claims would weaken rather than strengthen them.
The standard applied here is: \emph{every externally attributable premise
has a source anchor; every mathematical claim has a derivation; every
computed claim has a reproducible computation; every measured claim has a
traceable datum; every interpretive claim has an explicit evidentiary
status.} Identities such as $D=2a$, $D_h=2wh/(w+h)$ and
Eqn~\xref{eqn:closed_form}{12}, the series inversion, the CFD values, the
computed Reynolds numbers and the uncertainty ratios belong to the second,
third and fourth categories and are marked as such rather than quoted.

\paragraph{How this index is maintained.} Table~\ref{tab:verbatim} is
\emph{generated} from the quotation register, not typed. Every row is
verified on every build: its text against the primary source and, where its
page can be resolved, its page against the paginated source; rows marked
$^{\dagger}$ are text-verified only. Locators are corrected in
the register and propagate here on the next build.

\subsubsection*{Gifford 2014~\cite{b1}}
\textbf{p.~4498} ($\fgap$ regime): Gifford describes ``devices that remove
less than $0.05\%$ of the fluid from the central channel at each gap,''
i.e.\ $\fgap < 5\times10^{-4}$. \emph{This is the regime the source
reports, not a design target it states.}\\
\textbf{p.~4498} (Eqn~1): $Q_s(i)=Q_s(i-1)+\fgap Q_c(i-1)$ (no row index on
$\fgap$).\\
\textbf{p.~4498} (Eqn~3): continuity, from ``the total volumetric flow
within the device before and after each gap must be constant.''\\
\textbf{p.~4498} (Eqn~5): the rectangular-channel resistance carries
citation [20] ($=$ Yang/\"{U}ndar/Zahn, \cite{b6}); ``$L$ is the length,
$w$ is the width and $d$ is the depth of the channel.'' Clarified on
\textbf{p.~4499}: $L$ is ``the size (along the direction of flow) of an
obstacle in the device.''\\
\textbf{p.~4500} (calibration): ``33 test devices to explore a wide range
of $\fgap$ values ranging linearly from'' $6.40\times10^{-5}$ (device \#1)
to $5.76\times10^{-4}$ (device \#33); $\fgap^*$ read visually;
$\fgap=3.34\times10^{-4}$ selected between beads of diameter
$4.7~\si{\micro\meter}$ and $8.3~\si{\micro\meter}$ (Fig.~4 caption).\\
\textbf{p.~4500, Fig.~3 caption} (depth \emph{unreported}): states only
$w_c = 100,125,150~\si{\micro\meter}$, no depth (the
$80~\si{\micro\meter}$ value belongs to the Fig.~4 device).\\
\textbf{p.~4503}: ``the size cutoff is determined primarily by the width of
the streamlines diverted into the side (filtrate) channels.''

\subsubsection*{Yang 2006~\cite{b6} (resistance series)}
\textbf{p.~875}: its Eqn~3 is ``the analytical solution of flow resistance
within a rectangular duct'': the series Gifford's Eqn~5 cites and that
Eqn~\xref{eqn:gifford5}{5} reproduces.

\subsubsection*{Xia 2016~\cite{b2} / 2022 Correction~\cite{b3}}
\textbf{Correction p.~1} (verbatim): the article ``contains an error in
equation 1, where the terms $w_c(i)$ and $w_c(\text{ref})$ are inverted'';
it should read $\fgap(i)=(w_c(\text{ref})/w_c(i))^2\,\fgap(\text{ref})$.

\subsubsection*{Lezzar 2022~\cite{b4} (third lineage device)}
\textbf{p.~2}: the CIF design in that study targets a $6~\si{\micro\meter}$
critical diameter, and ``any CIF design must incorporate thousands of
gaps.'' The unpublished gap count $N$ is the quantity to which the
per-gap over-extraction factor $\sqrt{c}$ is tied
(Section~\xref{app:units_cutoff}{``Cutoff-Law Prefactor Decomposition''}). \emph{Verification note:} two of the
three registered Lezzar quotes are the softest verifications in the
register, one passing under ligature-tolerant matching and one against a
page image rather than a text layer; they are not equivalent in strength
to the exact-match anchors elsewhere in this index.

\subsubsection*{Dinh 2024~\cite{b5}}
\textbf{p.~913} ($12$~fL threshold): ``PLTs with cell volumes larger than
12 fL are commonly considered `large','' the P-LCR classification cut.
\emph{A classification threshold, not a cutoff.}\\
\textbf{p.~915} ($\fgap$ definition): a CIF device operates ``by siphoning
a relatively small (and precisely calculated) fraction, $\fgap(i)$, of the
flow in the middle channel, $Q_m(i)$,'' so
$Q_{\mathrm{gap}}(i)=\fgap(i)Q_m(i)$.\\
\textbf{p.~916}: the taper, ``progressively increasing the width of the
side channels and decreasing the width of the middle channel.''\\
\textbf{p.~916, Fig.~1 caption (DESELECTED, $\kster=1$):} ``about half the
width of the extracted flow lamina'', superseded by the main text and
Inglis; the two Dinh readings are formally inconsistent, so one must be
retracted.\\
\textbf{p.~917, main text (ADOPTED for $\kster$):} particles with a radius
smaller than the lamina width ($a<w$) are carried into the side
channels; larger particles ($a>w$) remain in the middle channel.\\
\textbf{p.~917 (discoid platelet):} ``recovery of 12 fL PLTs
($2a\approx4.8~\si{\micro\meter}$)'', the physical discoid diameter
(cf.\ $2.84~\si{\micro\meter}$ equivalent-sphere; see
Section~\xref{app:units_particle}{``Particle and Separation Variables''}).\\
\textbf{p.~917, Fig.~caption (geometry):} gaps $19~\si{\micro\meter}$-wide;
middle channel narrows $75\to52~\si{\micro\meter}$; depth of all channels
$140~\si{\micro\meter}$.\\
\textbf{p.~919, Fig.~4 caption (validation dispersion):} ``Values shown are
mean $\pm$ standard deviation (panel A: $n=3$; panel B: $n=4$).'' The same
caption reports the filtrate volume recoveries, $50.5\pm3.3\%$ and
$64.9\pm3.6\%$.

\subsubsection*{Yamada \& Seki 2005~\cite{b15} (geometric exclusion)}
\textbf{p.~1233--1234}: a particle centre cannot approach a sidewall closer
than one particle radius, and the method uses that fact. This is the
geometric-exclusion premise beneath $\kster$; it is \emph{not} a source for
the finite-particle factor $C$, which is cited to \cite{b19} and
\cite{b21}.

\subsubsection*{Davis 2008~\cite{b22} (cross-family DLD comparison)}
\textbf{pp.~14, 18, 22--24}: the DLD critical diameter is about twice the width
of the first streamline; the correlation with terms $A=1.4$ and $B=0.48$
is fitted to critical diameters measured on ``about 20 different devices''
and was used to design the thresholds the thesis reports, and its parabolic
approximation is checked against simulation. Used only in the cross-family
comparison (Section~\xref{app:units_crosslineage}{``Cross-Lineage Validation''}), never to fix a CIF
constant. \emph{Verification note:} the archived text carries the
thesis's printed page numbers, and these locators are verified against them.

\subsubsection*{Inglis 2006~\cite{b17} (steric basis)}
\textbf{p.~656}: the boundary-streamline fraction satisfies the
parabolic-profile cubic $\varepsilon = 3y^2 - 2y^3$ (identical to the CIF
skim cubic), and the critical diameter is twice the first-streamline
width, which Inglis writes $D_c = 2\beta$ (its Eqn~1; Greek $\beta$ is his
symbol, rendered $b$ here per the notation convention), giving
$\kster=\kstericVal$.

\subsubsection*{Inglis \& Herman 2013~\cite{b18} (lamina / branch flow)}
\textbf{p.~1724}: a particle is ``restricted away from walls by a distance
equal to their radius,'' so a boundary stream of width $b$ sets a cutoff
radius $b$ (equivalently $\kster=\kstericVal$); the abstract reports
``high filtration ratios (14--29\%) on beads and fungal spores'' with ``a
sharp size cutoff.''\\
\textbf{p.~1727} (Methods): ``The measured depth of the SU-8 mold was
48'' \si{\micro\meter}, against $50~\si{\micro\meter}$ designed.\\
\textbf{p.~1728, Table~1}: design versus measured cutoff diameters
($7.2/5.0=1.44$; $8.9/6.0=1.48$; $9.7/8.0=1.21$).\\
\textbf{p.~1729} (Discussion): ``the experimentally observed cutoff
diameters are significantly higher than the cutoff sizes determined from
our design process,'' and the shallower channel ``increases the main
channel resistance more than it does the branch channel resistance''
(see Section~\xref{app:units_crosslineage}{``Cross-Lineage Validation''}).

\subsubsection*{MOPSA 2017~\cite{b19} (finite-particle method class)}
\textbf{p.~2}: the simulation ``treats the particle as a two-dimensional
rigid circular object instead of a single point'': the finite-body
\emph{method class} available for computing $C$, not an establishment of
the factor itself.

\subsubsection*{Doyeux \emph{et al.} 2011~\cite{b21} (factor-$C$ mechanism)}
\textbf{p.~359} (abstract): near an asymmetric bifurcation the
finite-sphere attraction is ``directed towards the low flow rate branch''
(anti-Zweifach--Fung): the \emph{mechanism} that raises the measured
cutoff above the steric value ($C>1$), not a derivation of the
multiplicative form. \emph{Locator note:} the register cites \emph{JFM}
pagination while the archived artifact is the arXiv preprint, which
paginates differently; the text verifies, the page does not.

\subsubsection*{Bruus 2008~\cite{b20} (resistance limit)}
The rectangular-duct series reduces to $R = 28.4\,\mu L/h^4$ at the square
cross-section $h=w$ (Eqn~\xref{eqn:gifford5}{5} limit check). This claim is
\emph{self-derived} rather than quoted: the registry's independent
square-duct anchor is $12/B(1)=28.4542$, checked by the derivation record
rather than by the quotation gate.

\subsubsection*{Formal Argument}
\textbf{P1.4} (PE outlet condition):
$p_{\text{center}}(L_{\mathrm{dev}})=p_{\text{side}}(L_{\mathrm{dev}})$ at the
outlet. \emph{The physical condition is not introduced here.} The lineage
source states that its resistance-ratio calculation holds ``provided the
hydrostatic pressures in the output collection areas are kept equal''
(\cite[p.~4500]{b1}). What is this work's is the \emph{formalisation} of that
condition as an equality of outlet-node pressures, and its use as the premise
of the closed-form filtrate identity
(Eqn~\xref{eqn:closed_form}{12}).\\
\textbf{P1.5} (inlet seal): $Q_s(0)=0$; side channels accumulate filtrate
only through the gap network. \emph{This one is a numbered assumption of this
work and carries no external anchor.}\\
\emph{Both are stated and used as premises in
Section~\xref{app:units_equations}{``Equations: Step-by-Step Unit Derivations''}.}

\paragraph{Consolidated index.} Table~\ref{tab:verbatim} maps every
load-bearing claim to its source location and a short verbatim anchor, so
a third party can confirm each in seconds. It is generated from the
quotation register; do not edit it by hand.

% The block below is GENERATED from the quotation register by gen_a10.py and is
% inlined here so this subsection is one self-contained file. Do not edit inside the
% markers: gate 8c re-splices the generator output into a copy and diffs, so a hand
% edit here fails the build. To change a locator, fix the REGISTER and rebuild with
%     python3 gen_a10.py quote_register.json cif_argument.yml --inline A1_A10_FINAL.tex
% It is a longtable and carries its own caption and label -- it must NOT be wrapped
% in \begin{table}, because a longtable is not a float.
% >>> BEGIN GENERATED A.10 TABLE (gen_a10.py) -- DO NOT EDIT THIS BLOCK
% GENERATED by gen_a10.py from quote_register.json -- DO NOT EDIT BY HAND.
% Every row here is covered by gate 8 (verbatim text) and gate 8b (page locator).
% To change a locator, fix the REGISTER and rebuild.
% Requires \usepackage{longtable} in the preamble.
\begingroup\scriptsize\setlength{\tabcolsep}{4pt}\renewcommand{\arraystretch}{1.15}
\begin{longtable}{@{}>{\raggedright\arraybackslash}p{0.30\textwidth}>{\raggedright\arraybackslash}p{0.17\textwidth}>{\raggedright\arraybackslash}p{0.45\textwidth}@{}}
\caption{Claim $\rightarrow$ source (location) $\rightarrow$ verbatim anchor. Generated from the quotation register. Every row's \emph{text} is verified against the primary source on every build; \emph{pages} are checked where they can be resolved: against per-page images, or, for the Davis thesis, against the page numbers its text carries. Rows marked $^{\dagger}$ are text-verified but not page-verified.}\label{tab:verbatim}\\
\toprule
Claim in glossary & Source (location) & Verbatim anchor (short) \\
\midrule\endfirsthead
\multicolumn{3}{@{}l}{\footnotesize\itshape Table \thetable\ (continued)}\\
\toprule
Claim in glossary & Source (location) & Verbatim anchor (short) \\
\midrule\endhead
\midrule\multicolumn{3}{r@{}}{\footnotesize\itshape continued on the next page}\\
\endfoot
\bottomrule\endlastfoot
V\&V-20 defines u\_D to include all effects of averaging & \cite[p.~1-6]{b25} & ``includes all effects of averaging, includes all random and systematic uncertainty components'' \\  % A-VV1
Critical diameter is twice the first-streamline width (independent corroboration) & \cite[p.~14]{b22} & ``about twice the width of the first streamline'' \\  % P-C1.6
DLD critical-diameter correlation, terms A and B & \cite[p.~23]{b22} & ``with terms A=1.4 and B=0.48'' \\  % P-C6.1
That correlation was used to design the thresholds the thesis reports & \cite[p.~24]{b22} & ``equation 2-22 was used to design all the thresholds discussed herein'' \\  % P-C6.2
Critical diameters measured across about twenty devices & \cite[p.~22]{b22} & ``about 20 different devices over a range of G from 1.3 micron to 38 micron'' \\  % P-C6.3a
Parabolic approximation checked against simulation & \cite[p.~18]{b22} & ``The Femlab simulation shows this simple parabolic approximation is a reasonable estimation for the velocity\,\ldots'' \\  % P-C6.3b
Definition of the per-gap extraction fraction & \cite[p.~915]{b5} & ``siphoning a relatively small (and precisely calculated)'' \\  % A-D1
V1 filtrate-volume recovery, the yield-axis experimental datum (panel B, n=4) & \cite[p.~919]{b5} & ``Volume recovery in filtrate was 50.5 $\pm$ 3.3\% for f* gap = 1.04'' \\  % A-D10
Dinh selects f*\_gap as an initial value alongside the initial channel width & \cite[p.~915]{b5} & ``the initial values for the width of the middle channel and the filtration fraction'' \\  % A-D11
Dinh reports three initial f*\_gap values across three designs & \cite[p.~917 (Fig. 2 caption)]{b5} & ``for f* gap values of 1.04 $\times$ 10$-$4 , 1.12 $\times$ 10$-$4, and 1.28 $\times$ 10$-$4, respectively'' \\  % A-D12
V3 filtrate-volume recovery, the yield-axis experimental datum (panel B, n=4) & \cite[p.~919]{b5} & ``64.9 $\pm$ 3.6\% for f* gap = 1.28'' \\  % A-D13
Dinh describes the taper as widening side channels and narrowing the middle & \cite[p.~916]{b5} & ``progressively increasing the width of the side channels and decreasing'' \\  % A-D14
A platelet-size classification threshold (not the cutoff) & \cite[p.~913]{b5} & ``cell volumes larger than 12 fL are commonly considered'' \\  % A-D2
Physical discoid diameter of a 12 fL platelet & \cite[p.~917]{b5} & ``recovery of 12 fL PLTs (2a $\approx$ 4.8 \si{\micro\meter}) in the filtrate'' \\  % A-D3
Reported dispersion is standard deviation, with panel sample sizes & \cite[p.~919]{b5} & ``Values shown are mean $\pm$ standard deviation (panel A: n = 3; panel B: n = 4)'' \\  % A-D4
Device operating flow rate & \cite[p.~915]{b5} & ``a CIF device operating at a flow rate of 30'' \\  % A-D5
Smallest platelets are recovered in the filtrate below the fluid rate & \cite[p.~917]{b5} & ``recovery of the smallest 2 fL PLTs (a $\approx$ 1 \si{\micro\meter}) in the filtrate was consistently lower than the percent of extracted fluid volume'' \\  % A-D6
Dinh's own explanation: crowding by retained platelets & \cite[p.~917]{b5} & ``crowding of the flow laminae near the filtration gaps by retained'' \\  % A-D7
Crowding impedes the outflow of the smallest platelets & \cite[p.~917]{b5} & ``impeded the free outflow of the smallest PLTs'' \\  % A-D8
Platelet diameter is the major/equatorial axis of an oblate spheroid & \cite[p.~917]{b5} & ``range in diameter (major/equatorial axis, 2a) from about 2.0'' \\  % A-D9
\textbf{DESELECTED} ($\kster=1$ reading), superseded by the main text and Inglis; the two Dinh readings are formally inconsistent, so one must be retracted: Caption reading places the cutoff at half the lamina width & \cite[p.~916]{b5} & ``Particles with an effective radius smaller than about half the width of the extracted flow\,\ldots'' \\  % P-C1.2-caption
\textbf{ADOPTED} for $\kster=2$: Main-text capture criterion, at a = lamina width & \cite[p.~917]{b5} & ``For particles with a radius smaller than the width of the extracted flow lamina (a\,\ldots'' \\  % P-C1.2-lower
Device geometry: gap width, middle-channel taper, channel depth & \cite[p.~917 (Fig. 2 caption)]{b5} & ``All filtration gaps in the linear section were 19 \si{\micro\meter}-wide. The middle channel narrowed from\,\ldots'' \\  % P-C7.1-geom
Platelets are discoid, not spherical & \cite[p.~913]{b5} & ``1 \si{\micro\meter}-thick oblate spheroids (discoids) ranging from 2 to 7 \si{\micro\meter} in diameter'' \\  % P-LIM.1
Discoid platelets tumble near the channel walls & \cite[p.~921]{b5} & ``the characteristic tumbling motion oblate spheroids undergo when flowing near the channel walls'' \\  % P-LIM.2
Bead bracket, lower end: a 1 um bead is recovered in the filtrate at the fluid rate & \cite[Fig. S1(A)]{b5}$^{\dagger}$ & ``fluid volume recovery in filtrate was 52\% and bead recovery in filtrate was 54\%'' \\  % P-C7.2
Bead bracket, upper end: a 4.7 um bead is retained against the fluid & \cite[Fig. S1(B)]{b5}$^{\dagger}$ & ``fluid volume recovery in filtrate: 50\%, bead recovery in filtrate: 1.8\%'' \\  % P-C7.3
Finite particle size shifts the effective separation line at a bifurcation & \cite[p.~368]{b21}$^{\dagger}$ & ``might play a role, due to the finite size of the flowing objects'' \\  % P-C3d.5
Finite-sphere attraction is toward the low-flow branch (mechanism for C > 1) & \cite[abstract]{b21}$^{\dagger}$ & ``such an attraction does not occur but is, on the contrary, directed towards the low\,\ldots'' \\  % P-C6.4
The apparent high-flow enrichment is an inlet-distribution effect & \cite[abstract]{b21}$^{\dagger}$ & ``the enrichment in particles in the high flow rate branch is shown to be mainly\,\ldots'' \\  % P-C6.5
Flow recursion, Eqn 1 & \cite[p.~4498]{b1} & ``Qs(i) = Qs(i $-$ 1) + fgapQc(i $-$ 1)'' \\  % A-G1
Gifford built four devices spanning f\_gap from 2.2e-4 to 10.0e-4 & \cite[p.~4500]{b1} & ``four individual devices with a range of values of fgap (2.2--10.0 $\times$ 10$-$4'' \\  % A-G10
Gifford names Eqn 5 variables: length, width, depth of the channel & \cite[p.~4498]{b1} & ``L is the length, w is the width and d is the depth of the\,\ldots'' \\  % A-G11
Gifford clarifies L as an obstacle size along the flow direction & \cite[p.~4499]{b1} & ``the size (along the direction of flow) of an obstacle in the device'' \\  % A-G12
The Fig. 4 device is 80 um deep with a 100 um central channel & \cite[p.~4500 (body; cf. the Fig. 4 caption, which gives 80 \si{\micro\meter} $\pm$ 3 \si{\micro\meter})]{b1} & ``The depth of the device shown in Fig. 4 was 80 \si{\micro\meter}; the width of the central channel was 100 \si{\micro\meter}'' \\  % A-G13
Gifford describes the gaps as wide, about 20 um & \cite[p.~4498]{b1} & ``wide (~20 \si{\micro\meter}) gaps'' \\  % A-G14
Gifford names the total number of gap rows & \cite[p.~4499]{b1} & ``the necessary total number of rows of gaps'' \\  % A-G15
Gifford defines f*\_gap as the per-gap fraction below which beads are retained & \cite[p.~4500]{b1} & ``the filtration fraction per gap below which the beads of a spe'' \\  % A-G16
Continuity, Eqn 3, from incompressibility & \cite[p.~4498]{b1} & ``the total volumetric flow within the device before and after each gap must be constant'' \\  % A-G2
The reported per-gap extraction regime (not a stated design target) & \cite[p.~4498]{b1} & ``devices that remove less than 0.05\% of the fluid from the central channel at each gap, i.e. fgap < 5 $\times$ 10$-$4'' \\  % A-G3
The 33-device calibration array and its f\_gap span & \cite[p.~4500]{b1} & ``33 test devices to explore a wide range of fgap values ranging linearly from'' \\  % A-G4
The f\_gap selection bracket between two bead sizes & \cite[p.~4500]{b1} & ``retain beads with diameter 4.7 \si{\micro\meter} and diameter 8.3 \si{\micro\meter}'' \\  % A-G5
Calibration figure gives three channel widths and no depth & \cite[p.~4500]{b1} & ``Effect of central channel width (wc = 100 \si{\micro\meter}, 125 \si{\micro\meter}, 150 \si{\micro\meter})'' \\  % A-G6
Centre- and side-channel flow rates defined & \cite[p.~4498]{b1} & ``a simple recursive eqn (1) for the volumetric flow rate in the central'' \\  % A-G7
Centre and side segment resistances defined & \cite[p.~4498]{b1} & ``the resistance of a side channel segment, Rs(i), and the central channel segment resistance, Rc'' \\  % A-G8
Outlet hydrostatic pressures are kept equal (the P1.4 condition) & \cite[p.~4500]{b1} & ``provided the hydrostatic pressures in the output collection areas are kept equal'' \\  % A-G9
Cutoff is set by the diverted-lamina width, not the gap width & \cite[p.~4503]{b1} & ``the size cutoff is determined primarily by the width of the streamlines diverted into the\,\ldots'' \\  % P-C1.1
Platelets are approximately discoid & \cite[p.~4502]{b1} & ``approximately discoid'' \\  % P-LIM.3
Platelet-rich plasma behaviour in the lineage device & \cite[p.~4502]{b1} & ``we observed the behavior of PRP flowing through the same three arrays'' \\  % P-LIM.4
Pressure, flow and resistance map to voltage, current and resistance & \cite[p.~3312]{b14} & ``pressure drop, flow rate, and flow resistance in microchannels are analogous to electrical voltage, current,\,\ldots'' \\  % A-HE1
Inglis derives the streamline width from a parabolic profile through the gap & \cite[p.~656]{b17} & ``By assuming a conventional parabolic flow profile through the gap, with zero velocity at the\,\ldots'' \\  % A-IN1
Critical diameter is twice the first-streamline width, D\_c = 2b & \cite[p.~656]{b17} & ``The critical particle diameter, Dc, for determining which path a particle will follow is then\,\ldots'' \\  % P-C2.1
First-streamline width defined by integrating the flow profile from the post edge & \cite[p.~656]{b17} & ``we integrate the flow profile'' \\  % P-C2.3
First-streamline construction, second statement & \cite[p.~655]{b17} & ``is the width of the first streamline'' \\  % P-C2.4
Branch-flow filtration ratios and a sharp size cutoff & \cite[p.~1724]{b18} & ``high filtration ratios (14--29\%) on beads and fungal spores'' \\  % A-IH1
Measured cutoff diameters exceed the design values & \cite[p.~1729]{b18} & ``The experimentally observed cutoff diameters are significantly higher than the cutoff sizes determined from our design process.\,\ldots'' \\  % A-IH2
Fabricated mould depth fell short of the design depth & \cite[p.~1727]{b18} & ``The measured depth of the SU-8 mold was 48'' \\  % A-IH3
A shallow mold raises main-channel resistance more than branch resistance & \cite[p.~1729]{b18} & ``increases the main channel resistance more than it does the branch channel resistance'' \\  % A-IH4
More branch flow raises the critical size at each branch & \cite[p.~1729]{b18} & ``more fluid flow in the branches, increasing the critical size at each branch'' \\  % A-IH5
Excluding a particle can need a streamline narrower than its radius & \cite[p.~1729]{b18} & ``excluding a particular size particle from a branch may require a streamline width that is as much as 50\% smaller than the particle radius'' \\  % A-IH6
Inglis and Herman name a second contributing reason & \cite[p.~1729]{b18} & ``A second contributing reason is that the devices, as fabricated'' \\  % A-IH7
The finite-particle attraction runs contrary to the Zweifach-Fung effect & \cite[p.~1729]{b18} & ``contrary to the situation in the Zweifach-Fung effect'' \\  % A-IH8
The branches are narrow, so a shallow mould raises main resistance more & \cite[p.~1729]{b18} & ``because the branches are only 25'' \\  % A-IH9
Same mechanism and exclusion premise, independent (non-lineage) group & \cite[p.~1724]{b18} & ``Assuming that particles do not affect streamlines and are geometrically restricted away from walls by\,\ldots'' \\  % P-C1.5
The same source's own caveat on that mechanism & \cite[p.~1724]{b18} & ``This is a frequent assumption in the microfluidic particle separation literature, that is not always\,\ldots'' \\  % P-C1.5-hedge
The same near-wall lamina integrated over the same rectangular-duct field & \cite[p.~1724]{b18} & ``the cutoff size is determined by the width of the boundary layer in the main\,\ldots'' \\  % P-C2.2
Third lineage device: the cutoff is set by the flow fraction, not the gap width & \cite[p.~2 (Sci Rep PDF pagination)]{b4} & ``It is the magnitude of this flow fraction---not the width of the gaps---that determines the\,\ldots'' \\  % P-C1.4
Third lineage device: 6 um critical-diameter target & \cite[p.~2 (Sci Rep PDF pagination)]{b4} & ``The critical diameter of the CIF design used in this study was 6'' \\  % P-C1.4b
Gap count is unpublished; sqrt(c) is tied to it & \cite[p.~2 (Sci Rep PDF pagination)]{b4} & ``any CIF design must incorporate thousands of gaps to achieve the desired'' \\  % P-C1.4c
Finite-body method class: the particle treated as a rigid disc, not a point & \cite[p.~2]{b19} & ``treats the particle as a two-dimensional rigid circular object instead of a single point'' \\  % A-M1
Oh maps pressure to voltage, flow rate to current, and hydraulic to electric resistance & \cite[p.~515]{b13} & ``correlations of pressure to voltage, volumetric flow rate to current, and hydraulic to electric resistance'' \\  % A-HE2
Xia introduces the cutoff-diameter abbreviation c.d. & \cite[p.~2 (Sci Rep PDF pagination)]{b2} & ``particles above a certain size (or `critical diameter', c.d.)'' \\  % A-X1
RS-8 and RS-9 device parameters: cutoff, reference width and reference f\_gap & \cite[p.~6 (Sci Rep PDF pagination)]{b2} & ``nominal c.d. of 8\si{\micro\meter} (``RS-8'', wc(ref)= 100 \si{\micro\meter} and fgap(ref)= 7.2$\times$ 10$-$4 ) and 9 \si{\micro\meter} (``RS-9'', wc(ref)= 100 \si{\micro\meter} and fgap(ref)= 8.6$\times$ 10$-$4'' \\  % A-X2
OR-7 reference width and reference f\_gap & \cite[p.~10 (Sci Rep PDF pagination)]{b2} & ``OR-7 was patterned using wc= 100\si{\micro\meter} and fgap= 5.8$\times$ 10$-$4'' \\  % A-X3
The f\_gap ramp exists to hold the cutoff constant along a tapering channel & \cite[p.~9 (Sci Rep PDF pagination)]{b2} & ``In order to maintain a consistent c.d. along the length of an RS-CIF device, the value of fgap can no longer be kept constant and must be scaled relative to'' \\  % A-X4
The centre-channel taper exists to reduce shear, not to move the cutoff & \cite[p.~4 (Sci Rep PDF pagination)]{b2} & ``in order to reduce the maximum amount of shear stress experienced by particles within the\,\ldots'' \\  % A-X5
The 7 um cutoff is inherited by citation, not measured & \cite[p.~10 (Sci Rep PDF pagination)]{b2} & ``a combination of parameters we found corresponding to a c.d. of 7'' \\  % P-C3b.2
The citation points at the same figure the plateau is fitted from & \cite[p.~9 (Sci Rep PDF pagination)]{b2} & ``see Fig. 3 of Gifford et al.'' \\  % P-C3b.3
Equation 1 of the 2016 paper is corrected: the two widths were inverted & \cite[p.~1]{b3} & ``This Article contains an error in equation 1, where the terms wc(i) and wc(ref) are inverted'' \\  % P-C5.2
A particle centre is excluded from the wall by one radius & \cite[p.~1234]{b15} & ``the center position of a particle cannot be present within a certain distance from the\,\ldots'' \\  % P-C1.3
Geometric-exclusion premise, context sentence & \cite[p.~1233]{b15} & ``particles are concentrated and aligned onto the sidewalls'' \\  % P-C1.3b
The rectangular-duct resistance series the lineage resistance law cites & \cite[p.~875]{b6} & ``analytical solution of flow resistance within a rectangular duct'' \\  % A-Y1
\end{longtable}\endgroup

\vspace{2pt}\footnotesize $^{\dagger}$Text verified against the primary source; \emph{page} not machine-verified (a figure anchor rather than a page, or preprint-versus-journal pagination).
% <<< END GENERATED A.10 TABLE

\textbf{Self-derived (not from a single source).} The skim-cubic inversion
$\skimw=w_c[\sqrt{\fgap/3}+\tfrac19\fgap+\cdots]$ is obtained by
series-reverting $\fgap=3x^2-2x^3$ (verified term-by-term). The worked
V\&V-20 distances follow from the tabulated $F$-values with $n{=}4$. The
Reynolds figures ($D_h\approx75.8~\si{\micro\meter}$,
$U\approx0.069~\si{\meter\per\second}$,
$\mathrm{Re}\approx\ReWorstChannel$) follow from the Dinh geometry by
$D_h=2w_ch/(w_c+h)$ and $U=Q/A$. The square-duct limit $12/B(1)=28.4542$
is a registry anchor. The kernel-pairing table is generated from the
registry. All were re-derived independently.


%==============================================================
% A.11 -- CROSS-LINEAGE VALIDATION  (complete, corrected)
%
% The two tables stay FLOATS, deliberately. A.8's envelope table had to become a
% longtable because it ran several pages and a float that cannot fit is dropped
% with "Float too large for page". These two are three and four data rows; they
% fit, and they compile with zero overfull boxes and no X-column warning.
%
% CHANGES:
%   * A-IH2 quoted as the source writes it. The draft rendered "determined from
%     [the] design" where Inglis & Herman write "determined from our design
%     process": a bracketed insertion that also dropped a word.
%   * the two resistance sentences, quoted verbatim with a page but previously
%     outside the register, are now anchored (A-IH4, A-IH5).
%   * "Inglis" disambiguated. In the DLD parenthetical it meant the Inglis CUBIC,
%     a computation, while three paragraphs above "Inglis & Herman" is an
%     experimental family. One surname, two papers, two roles.
%   * every gated figure is a macro; \ref into A.4 became \xref.
%   * item 1's finite-particle MAGNITUDE and the Zweifach-Fung contrast anchored to
%     Inglis & Herman's own p.1729 words (A-IH6, A-IH8), replacing the coined
%     "anti-Zweifach-Fung" shorthand in the body with their phrasing; the TWO-reason
%     count anchored to their "A second contributing reason" (A-IH7). The two-effect
%     mapping is the source's structure, not this work's imposition.
%   * "cutoff" scoped to "cutoff DIAMETER", with A.4's sub-steric foot named as a
%     below-cutoff regime so it cannot read as a counterexample.
%   * the DLD bracket described as two computations running in OPPOSITE senses that
%     agree to a fraction of a percent at the top of the range, not one curve
%     sweeping the full bracket. Union bounds are now min/max OF THE LEGS, so the
%     bracket and the legs cannot print different values for one quantity.
%==============================================================
\subsection{Cross-Lineage Validation (Davis DLD; Inglis \& Herman branch flow)}
\label{app:units_crosslineage}

The cutoff prefactor's open residual predicts that a measured cutoff
diameter sits \emph{above} the idealized steric value
$\dcrit^{\text{ideal}}=2\skimw=2w_c\sqrt{\fgap/3}$. This is corroborated
across three device families: CIF
(Shevkoplyas lineage), deterministic lateral displacement (Davis), and
branch-flow filtration (Inglis \& Herman). \emph{Within the CIF lineage
itself there is no independence}: Gifford, Xia, Dinh and Lezzar share a
laboratory, and it is the two non-CIF families that carry the
out-of-lineage check.

\begin{table}[ht]
\centering
\caption{Inglis \& Herman 2013, Table~1: design vs.\ measured cutoff
diameter and filtration ratio (FR). The measured cutoff exceeds the design
in every case; the mean ratio is $\approx\ihRatioMean$. $\pm$ values are CV
over three runs.}
\label{tab:crosslineage}
\small
\begin{tabularx}{\textwidth}{@{}YYYYY@{}}
\toprule
Design $D_c$ (\si{\micro\meter}) & Design FR & Measured FR & Measured $D_c$ (\si{\micro\meter}) & Measured/Design \\
\midrule
5.0 & 13.3\% & $14.5\pm0.2\%$ & 7.2 & \ihRatioA \\
6.0 & 18.3\% & $20.4\pm0.5\%$ & 8.9 & \ihRatioB \\
8.0 & 28.4\% & $28.5\pm0.4\%$ & 9.7 & \ihRatioC \\
\midrule
\multicolumn{4}{@{}l}{\emph{Mean measured/design ratio}} & $\approx\ihRatioMean$ \\
\bottomrule
\end{tabularx}
\end{table}

\begin{table}[ht]
\centering
\caption{Three-family excess over the steric floor $\kster=2$, all
expressed as ratios to that floor so they compare directly. The magnitudes
differ, and that difference is a result.}
\label{tab:threefamily}
\small
\begin{tabularx}{\textwidth}{@{}lYY@{}}
\toprule
Family (source) & Excess over the steric floor & Basis \\
\midrule
CIF (Shevkoplyas) & $\kapp/\kster=\residualFolded$ & Gifford Fig.~3 plateau fit \\
DLD (Davis) & $\kapp/2 = \dldExcessLo$--$\dldExcessHi$ & recomputed from Davis eqn~2-22~\cite{b22} \\
Branch flow (Inglis \& Herman) & measured/design $=\ihRatioLo$--$\ihRatioHi$ & \cite{b18}, Table~1 \\
\bottomrule
\end{tabularx}
\end{table}

\textbf{What agrees, and what does not.} The consilience is in the
\emph{sign} and the \emph{floor}, not the magnitude. Every family's
measured cutoff \emph{diameter} exceeds the steric floor $\kster=2$ where the
three are compared: on CIF's plateau, and across the whole reported range of
the other two. (Below CIF's knee the apparent
prefactor falls under the floor, to the sub-steric foot of
Section~\xref{app:units_cutoff}{``Cutoff-Law Prefactor Decomposition''}; that
regime is the knee's own, and is not a counterexample.) \textbf{The magnitudes do not agree, and that is a result rather than
a discrepancy to be harmonised away}: CIF's excess is roughly \emph{twice}
the DLD and branch-flow families'. Correctly normalised, DLD shows a
smooth, knee-free apparent prefactor across $\varepsilon=0.01$--$0.5$,
bracketed by $\kapp=\dldKappaLo$--$\dldKappaHi$. \emph{Those two bounds are
two computations of the same DLD prefactor, not two laboratories, and they
run in opposite senses}: Davis's eqn~2-22 falls with $\varepsilon$, from
$\davisKappaHi$ to $\davisKappaLo$, while the exact Inglis \emph{cubic}
(\cite{b17}, the skim construction of
Section~\xref{app:units_equations}{``Equations: Step-by-Step Unit Derivations''})
rises, from $\inglisKappaLo$ to $\inglisKappaHi$. At the top of the range the two agree to
$\dldConvergePct\%$, so the width of the bracket is a between-method gap at
small $\varepsilon$, not size-dependence within either. Each is monotone
and neither shows a knee. Read ``Inglis'' there as the cubic, not as Inglis \& Herman, whose
branch-flow data are the separate third family above. Within the examined
DLD comparison datasets the normalized prefactor therefore does not
exhibit the CIF knee: the size-dependence is CIF-specific \emph{within the
validation corpus}, not demonstrably a universal feature of
streamline-skimming devices. \emph{The factor-of-two difference is an
observed quantitative distinction between CIF and the comparison families,
not an identified physical cause.} \emph{Note the two bases}: Table~\ref{tab:threefamily}'s DLD excess
$\dldExcessLo$--$\dldExcessHi$ is the Davis design correlation alone, which
is the leg the ``roughly twice'' comparison rides on, whereas the wider
bracket above adds the theoretical cubic and is used only to establish the
absence of a knee. Halving the bracket would give a different and
incomparable range. \textbf{Compare the depth-\emph{folded}
$\residualFolded$ with the other two families, not the depth-separated
triplet}: the published measured/design ratios are depth-inclusive.

\textbf{Interpretation of the branch-flow leg.} The authors state that the
``experimentally observed cutoff diameters are significantly higher than
the cutoff sizes determined from our design process'' (\cite[p.~1729]{b18}),
and they themselves name two contributing reasons for it (``A second
contributing reason is that the devices, as fabricated\ldots'',
\cite[p.~1729]{b18}), which correspond to the two non-unity factors of
Eqn~\xref{eqn:kappa_decomp}{7}. \emph{The
correspondence is to mechanism classes, not to this work's specific
factorization}, and the published data fix the \emph{product}
$\sqrt{c}\,C$ without fixing the split between its two factors. That
degeneracy is a property of the available measurements, not a gap in the
analysis: no published dataset varies $c$ and $C$ independently.
\begin{enumerate}
\item \textbf{Finite-particle displacement ($C$).} A finite sphere near a
bifurcation is drawn toward the \emph{low-flow-rate} branch, ``contrary to
the situation in the Zweifach-Fung effect'' (Doyeux \emph{et al.},
\cite{b21}; the contrast is Inglis \& Herman's own phrase,
\cite[p.~1729]{b18}). They invoke the mechanism themselves, agreeing with
Doyeux that ``excluding a particular size particle from a branch may
require a streamline width that is as much as 50\% smaller than the
particle radius'' (\cite[p.~1729]{b18}): the streamline that just excludes
the particle is narrower than its radius, raising the measured cutoff.
\item \textbf{Over-extraction from increased branch flow ($\sqrt{c}$).}
The realized SU-8 mold depth was $48~\si{\micro\meter}$ against a
$50~\si{\micro\meter}$ design (\cite[p.~1727]{b18}). Inglis \& Herman
attribute the resulting cutoff shift not to a skim-width change but to
resistance: the shallower channel ``increases the main channel resistance
more than it does the branch channel resistance'' ``because the branches
are only 25'' \si{\micro\meter} wide, which will mean ``more fluid flow in
the branches, increasing the critical size at each branch''
(\cite[p.~1729]{b18}). \emph{They give no magnitude for that shift}: the
48-versus-50 \si{\micro\meter} depth and the 25 \si{\micro\meter} branch
width are geometry, not a quantification of the cutoff change, which is
why the published data fix the product and not the split. This is per-gap \emph{over-extraction}, the
$\sqrt{c}$ factor, not the finite-aspect-ratio skim-width factor
$\Kdepth$. \emph{Note that this makes the $\ihRatioLo$--$\ihRatioHi$ range
a mixed measurement}: part of their excess is a fabrication deviation, so
it is not a pure finite-particle figure.
\end{enumerate}
The independent branch-flow literature therefore corroborates the sign and
the plausible physical classes of the CIF residual; it does not identify
the CIF-specific numerical factors $c$ and $C$ separately.



%==============================================================
% A.12 -- MASTER CHEAT-SHEET  (complete, corrected)
%
% CORRECTIONS, each traced to the subsection it contradicted:
%   * K_dcfd denominator was /Kdth while printing the /Kdcfd values, the same
%     formula-versus-value mismatch A.4 and A.9 had corrected. Now /Kdcfd,
%     matching the gated cfd_kdepth_obs.
%   * the DEPRECATED lower-case g was reintroduced as "gap width". A.1 retires
%     it: "lower-case g is needed for the cubic kernel; use G." Row now names
%     G, g_p and Lambda only.
%   * entrance-length figures were the pre-correction ones (47-61, Lambda 49,
%     0.96-1.2). A.9 now leads with Gifford's MEASURED pitch 46.1.
%   * u_D read "primary = SD", the framing A.6 retired once ASME V&V-20 S1-6
%     was read: the standard's u_D "includes all effects of averaging", so SD
%     is a CONSERVATIVE stand-in, not the standard's primary quantity.
%   * the Xia anchor was stated flatly at 5.17; A.7 downgraded it to a lead
%     because the agreement is numerically the depth correction itself.
%   * every computed figure is a macro; \ref into other subsections is \xref.
%   * longtable, not a float: thirty-odd rows do not fit one page.
%==============================================================
\subsection{Master Cheat-Sheet}
\label{app:units_cheatsheet}

A single consolidated lookup table for both performance axes, the cutoff
prefactor decomposition, the unification campaign, and the validation
layer. Values in \textbf{bold} are fixed constants; $w_c$-dependent and
device-dependent quantities are given as relations. \emph{This table
summarises; where it and a subsection differ, the subsection governs.}

% SYMLINT: SUMMARY -- this table REPEATS symbols defined in A.1-A.11; it does not
% define them. Without this directive the collision lint reads its thirty-odd rows as
% redefinitions and reports one duplicate per row.
\begingroup\scriptsize\setlength{\tabcolsep}{4pt}\renewcommand{\arraystretch}{1.2}
\begin{longtable}{@{}>{\raggedright\arraybackslash}p{0.11\textwidth}
>{\raggedright\arraybackslash}p{0.40\textwidth}
>{\raggedright\arraybackslash}p{0.09\textwidth}
>{\raggedright\arraybackslash}p{0.31\textwidth}@{}}
\caption{CIF master cheat-sheet. ``Unit'' is the practical CIF-lineage
convention; all derivations use SI.}\label{tab:cheatsheet}\\
\toprule Symbol & Value / defining relation & Unit & One-line meaning \\\midrule\endfirsthead
\multicolumn{4}{@{}l}{\scriptsize\itshape Table \thetable\ (continued)}\\
\toprule Symbol & Value / defining relation & Unit & One-line meaning \\\midrule\endhead
\bottomrule\endlastfoot
\multicolumn{4}{@{}l}{\emph{Geometry}}\\
$w_c$ & design input & \si{\micro\meter} & Center-channel width \\
$w,\ \skimw$ & $\skimw=w_c\sqrt{\fgap/3}$ (theory); $w$ empirical & \si{\micro\meter} & Diverted-lamina width ($b\equiv\skimw\approx w$) \\
$h$ & design input; Eqn~\xref{eqn:gifford5}{5} is written with $h\le w_{\mathrm{ch}}$, the two exchanged when $h>w_{\mathrm{ch}}$ & \si{\micro\meter} & Channel depth \\
$G$, $g_p$, $\Lambda$ & design inputs; $\Lambda=G+g_p$. \emph{Lower-case $g$ is the cubic kernel, not a length} (Table~\xref{tab:units_geom}{A.1}) & \si{\micro\meter} & Link width, post width, row pitch \\
$\phi$ & $\phi=G/\Lambda$ (computed); cf.\ $\phi_{\mathrm{cal}}$ (external) & --- & Wall porosity; one of three residual model gaps \\
$N$ & design input & --- & Number of gap rows \\
\addlinespace
\multicolumn{4}{@{}l}{\emph{Yield axis}}\\
$\fgap$ & $Q_{\mathrm{gap}}=\fgap Q_c$; Gifford's stated practical regime is $\fgap<5\times10^{-4}$ & --- & Per-side, per-gap extraction fraction \\
$R_c,R_s$ & Eqn~\xref{eqn:gifford5}{5} (Yang series, \cite{b6}) & \si{\pascal\second\per\meter\cubed} & Segment hydraulic resistances \\
$F_{\mathrm{outlet}}$ & $=2R_c/(R_s+2R_c)$, with no $\fgap$; this absence \emph{is} the decoupling & --- & Outlet divider under PE (P1.4) and $Q_s(0)=0$ (P1.5) \\
$F_{\mathrm{recursion}}$ & $=1-\prod_i(1-2\fgap(i))$ & --- & Cumulative march; equals $F_{\mathrm{outlet}}$ for a self-consistent device \\
$\Pi$ & $\prod_i(1-2\fgap(i))$ & --- & Retained fraction \\
\addlinespace
\multicolumn{4}{@{}l}{\emph{Selectivity axis (cutoff)}}\\
$\dcrit$ & $\dcrit=\kapp w_c\sqrt{\fgap/3}$ & \si{\micro\meter} & Critical (cutoff) diameter \\
$\ud$ & $\ud=\dcrit/w_c$; plateau $\ge\plateauUwc$, foot $<\footUwc$ & --- & Reduced size ratio (sets the bins) \\
$x$ & root of $\fgap=3x^2-2x^3$; $x=\skimw/w_c$ & --- & Reduced skim width (Inglis cubic) \\
$\beta,\gamma$ & $\partial\ln\dcrit/\partial\ln\fgap=\tfrac12\,(\approx\expBeta)$; $\partial\ln\dcrit/\partial\ln w_c=1\,(\approx\expGamma)$ & --- & Cutoff scaling exponents ($\neq$ lamina width $b$) \\
$\kapp\equiv\kg$ & supplied \textbf{\kappVal} (95\% CI $[\kappCIlo,\kappCIhi]$, CV $\kappCV\%$); refit $\kappFit$ & --- & Apparent prefactor, \textbf{cubic branch} \\
$\kF$ & \textbf{\kFval} $=\kg\sqrt{\overline{c_2}/3}$ at assumed $h_{\mathrm{cal}}=80~\si{\micro\meter}$. \textbf{Not interchangeable with $\kg$}: crossing shifts $\dcrit$ by $\crossUp\%$ / $\crossDown\%$; the kernel choice alone shifts it by $\choiceIn\%$ to $\choiceOut\%$ & --- & Prefactor, \textbf{exact $F_{\mathrm{lam}}$ branch} \\
$\kster$ & \textbf{\kstericVal} ($D_c=2b$, Inglis's $2\beta$; Dinh $a<w$) & --- & Steric factor (radius $=$ lamina) \\
$\Kin$ & \textbf{\KinVal} (single-gap CFD) & --- & In-plane funnelling (none) \\
$\Kdth$ & analytic $\KdepthAnalyticShallow/\KdepthAnalyticMid/\KdepthAnalyticDeep$ at $\alpha=0.80/1.40/2.69$\,$^{\dagger}$ & --- & Depth factor, theorem \\
$\Kdcfd$ & CFD $\KdepthCfdShallow/\KdepthCfdMid/\KdepthCfdDeep$; $|\Kdcfd-\Kdth|/\Kdcfd\le0.6\%$, at most $\KdepthDevMax\%$. \emph{The denominator is the CFD anchor}, matching the gated value & --- & Depth factor, CFD \\
$\sqrt{c}\,C$ & $\Kdepth\sqrt{c}\,C=\kapp/\kster=\mathbf{\residualFolded}$ (depth-folded); $\sqrt{c}\,C=\residualSepShallow/\residualSepMid/\residualSepDeep$ (depth-separated). Two factors; \emph{three} model gaps & --- & Open calibrated residual \\
$\acal$ & $(0.80,0.64,0.5333)$ at assumed $h_{\mathrm{cal}}=80~\si{\micro\meter}$; $\Delta_{\mathrm{transfer}}=\transferDecl\%$. The $+26.9\%$ kernel/depth sensitivity is \emph{ungated}, from a specification sweep this record does not reproduce & --- & Calibration array; \textbf{unreported depth} \\
\addlinespace
\multicolumn{4}{@{}l}{\emph{Particle}}\\
$a,\ D$ & $D=2a$; $V_{\mathrm{sph}}=\tfrac43\pi a^3$ & \si{\micro\meter} & Radius / diameter (rigid sphere) \\
$V_{12}$ & $12$~fL classification threshold $\to D\approx\DTwelveFL~\si{\micro\meter}$ equivalent sphere; $2a\approx4.8~\si{\micro\meter}$ (Dinh discoid) & fL & \textbf{Not} a minimum retained volume \\
\addlinespace
\multicolumn{4}{@{}l}{\emph{Regime and envelope}}\\
$\mathrm{Re}$ & channel $\approx\ReWorstChannel$; gap/lamina $\approx\ReGapScale\ll1$ & --- & Not globally Stokes; creeping at the selecting scale \\
$L_e/\Lambda$ & $L_e\approx\LeOutlet$--$\LeInlet~\si{\micro\meter}$ (outlet to inlet) against the Gifford Fig.~4 pitch $\Lambda=G+g_p=46.1~\si{\micro\meter}$, i.e.\ $L_e/\Lambda\approx\LeOverLamGOutlet$--$\LeOverLamGInlet$ & --- & Exceeds the row pitch; expected benign either way (per-rung perturbation of order $\fgap$) \\
$u_{\mathrm{lo}}$, $\sigma_0$, $\phi_{\mathrm{cal}}$ & supplied & --- & External inputs, undischargeable \\
$\eta$ & $\eta=\etaVal$, a computed magnitude, not a free parameter (Table~\xref{tab:units_envelope}{A.8}) & --- & Realized-yield factor, \textsc{measured} \\
\addlinespace
\multicolumn{4}{@{}l}{\emph{Validation (V\&V-20)}}\\
$E$ & $E=S-D$ (sim $-$ data) & --- & Validation comparison error; $S$ and $D$ are reused letters \\
$\alpha_{\mathrm{3D}}$ & $(R_c/R_s)_{\mathrm{app}}/(R_c/R_s)_{\mathrm{Yang}}$; $\alphaThreeDVone$/$\alphaThreeDVthree$ & --- & Outlet correction reconciling 1-D with 3-D \\
$\uD$ & the published $\mathrm{SD}$, a \emph{conservative} choice. V\&V~20's $\uD$ for a mean of $n$ has random component $\mathrm{SEM}=\mathrm{SD}/\sqrt{n}$, reported alongside & --- & Data uncertainty \\
$\uval$ & $\sqrt{\unum^2+\uinput^2+\uD^2}$ & --- & Combined validation uncertainty \\
$F$ check & 1-D closed form $\EsdClosedVOne\sigma$ / $\EsdClosedVThree\sigma$; 3-D COMSOL $\EsdComsolVOne\sigma$ / $\EsdComsolVThree\sigma$ on the SD scale & --- & The 3-D solve lands within $1\sigma$ on the SD scale, $2\sigma$ on SEM \\
Cutoff check & Dinh band $[\dinhBandLo,\dinhBandHi]$ inside measured $(1.0,4.7)~\si{\micro\meter}$; Xia $\kapp\approx\xiaKappaMean$ vs $\plateauKappa$, reported as a \emph{lead} & \si{\micro\meter} & Both consistency checks: Dinh held out but weak, Xia not independent \\
\end{longtable}\endgroup

\noindent{\footnotesize $^{\dagger}$ The $\kapp$ fit uses the Gifford
Fig.~3 array, whose depth is \emph{unreported}; $\alpha=0.8$ assumes the
$80~\si{\micro\meter}$ depth of the \emph{separate} Fig.~4 device, of which
Gifford writes ``The depth of the device shown in Fig.~4 was 80 \si{\micro\meter};
the width of the central channel was 100 \si{\micro\meter}''
(\cite[p.~4500]{b1}), and $\Kdth$ varies across the three
calibration widths ($\KdepthCalNarrow/\KdepthCalMid/\KdepthCalWide$ at
$w_c=100/125/150~\si{\micro\meter}$, i.e.\ $\alpha=0.80/0.64/0.5333$). See
Section~\xref{app:units_cutoff}{``Cutoff-Law Prefactor Decomposition''}.}

%==============================================================
% BIBLIOGRAPHY
%==============================================================
\begin{thebibliography}{99}

\bibitem{b1}
S.~C. Gifford, A.~M. Spillane, S.~M. Vignes, and S.~S. Shevkoplyas,
``Controlled incremental filtration: A simplified approach to design
and fabrication of high-throughput microfluidic devices for selective
enrichment of particles,'' \emph{Lab Chip}, vol.~14, no.~23,
pp.~4496--4505, 2014, doi: 10.1039/C4LC00785A.

\bibitem{b2}
H.~Xia, B.~C. Strachan, S.~C. Gifford, and S.~S. Shevkoplyas,
``A high-throughput microfluidic approach for 1000-fold leukocyte
reduction of platelet-rich plasma,'' \emph{Sci. Rep.}, vol.~6,
p.~35943, 2016, doi: 10.1038/srep35943.

\bibitem{b3}
H.~Xia, B.~C. Strachan, S.~C. Gifford, and S.~S. Shevkoplyas,
``Author Correction: A high-throughput microfluidic approach for
1000-fold leukocyte reduction of platelet-rich plasma,''
\emph{Sci. Rep.}, vol.~12, p.~14689, 2022, doi:
10.1038/s41598-022-18449-5.

\bibitem{b4}
D.~L. Lezzar \emph{et~al.}, ``A high-throughput microfluidic device
based on controlled incremental filtration to enable
centrifugation-free, low extracorporeal volume leukapheresis,''
\emph{Sci. Rep.}, vol.~12, p.~13798, 2022, doi:
10.1038/s41598-022-16748-5.

\bibitem{b5}
M.~T.~P. Dinh, A.~Mukhamedshin, K.~Abhishek, F.~W. Lam, S.~C. Gifford,
and S.~S. Shevkoplyas, ``Separation of platelets by size in a
microfluidic device based on controlled incremental filtration,''
\emph{Lab Chip}, vol.~24, no.~4, pp.~913--923, 2024, doi:
10.1039/D3LC00842H.

\bibitem{b6}
S.~Yang, A.~\"{U}ndar, and J.~D. Zahn, ``A microfluidic device for
continuous, real time blood plasma separation,'' \emph{Lab Chip},
vol.~6, no.~7, pp.~871--880, 2006, doi: 10.1039/B516401J.

\bibitem{b7}
COMSOL Multiphysics\textregistered\ v.~6.3, COMSOL AB, Stockholm,
Sweden, \url{www.comsol.com}.

\bibitem{b12}
A.~H. Routt, ``CIF reference implementation \texttt{v3.20}'' [software and
data: computational notebook and archived COMSOL reports], \emph{Zenodo},
2026, doi: 10.5281/zenodo.20146562. Notebook sections and blocks are cited
by number; the four five-gap reports (V1 and V3 at production and
${\sim}5\times$ coarser meshes) are cited as \emph{COMSOL reports}.

\bibitem{b12a}
A.~H. Routt, ``CIF verified derivation record'' [software], release tag
and content hash recorded in the dissertation reproducibility archive,
2026. \emph{Numerical authority: every constant is recomputed on each run
and pinned by content hash, and the whole record is regenerated and
re-verified on every commit.}

\bibitem{b13}
K.~W. Oh, K.~Lee, B.~Ahn, and E.~P. Furlani, ``Design of pressure-driven
microfluidic networks using electric circuit analogy,'' \emph{Lab Chip},
vol.~12, no.~3, pp.~515--545, 2012, doi: 10.1039/C2LC20799K.

\bibitem{b14}
Z.~Li, C.~Liu, and J.~Sun, ``Hydraulic--electric analogy for design and
operation of microfluidic systems,'' \emph{Lab Chip}, vol.~23,
pp.~3311--3327, 2023, doi: 10.1039/D3LC00265A.

\bibitem{b15}
M.~Yamada and M.~Seki, ``Hydrodynamic filtration for on-chip
particle concentration and classification utilizing microfluidics,''
\emph{Lab Chip}, vol.~5, no.~11, pp.~1233--1239, 2005, doi:
10.1039/B509386D.

% b8-b11 (design-automation) and b16 (PDMS shrinkage) were declared here but never
% cited in this appendix; they belong to the introduction and the sensitivity
% chapter. Removed so the reference list matches the text.
% ---------------- Inglis references (steric basis, kappa_steric = 2) ----------------
\bibitem{b17}
D.~W. Inglis, J.~A. Davis, R.~H. Austin, and J.~C. Sturm, ``Critical
particle size for fractionation by deterministic lateral displacement,''
\emph{Lab Chip}, vol.~6, no.~5, pp.~655--658, 2006, doi:
10.1039/B515371A.

\bibitem{b18}
D.~W. Inglis and N.~Herman, ``A scalable approach for high throughput
branch flow filtration,'' \emph{Lab Chip}, vol.~13, no.~9,
pp.~1724--1731, 2013, doi: 10.1039/C3LC50192B.

% ---------------- finite-particle mechanism and method ----------------
\bibitem{b19}
J.~Wang, V.~G.~J. Rodgers, P.~Brisk, and W.~H. Grover, ``MOPSA: A
microfluidics-optimized particle simulation algorithm,''
\emph{Biomicrofluidics}, vol.~11, no.~3, p.~034121, 2017, doi:
10.1063/1.4989860.

\bibitem{b20}
H.~Bruus, \emph{Theoretical Microfluidics}. Oxford, U.K.: Oxford
Univ.\ Press, 2008.

\bibitem{b21}
V.~Doyeux, T.~Podgorski, S.~Peponas, M.~Ismail, and G.~Coupier,
``Spheres in the vicinity of a bifurcation: elucidating the
Zweifach--Fung effect,'' \emph{J. Fluid Mech.}, vol.~674,
pp.~359--388, 2011, doi: 10.1017/S0022112010006567.

\bibitem{b22}
J.~A. Davis, \emph{Microfluidic Separation of Blood Components through
Deterministic Lateral Displacement}, Ph.D. dissertation, Princeton
University, 2008. (Eqn~2-22; recomputed here for the cross-family
comparison, Table~\ref{tab:threefamily}.)

\bibitem{b23}
F.~Durst, S.~Ray, B.~\"{U}nsal, and O.~A. Bayoumi, ``The development
lengths of laminar pipe and channel flows,'' \emph{J. Fluids Eng.},
vol.~127, no.~6, pp.~1154--1160, 2005, doi: 10.1115/1.2063088.
(Source of the entrance-length correlation, Eqn~\ref{eqn:entrance}, and of
its low-Reynolds limit $0.619$.)

\bibitem{b25}
\emph{Standard for Verification and Validation in Computational Fluid
Dynamics and Heat Transfer}, ASME V\&V 20-2009 (R2021), American Society
of Mechanical Engineers, New York, 2009.

\end{thebibliography}

\end{document}